% Leçon 7 — Expression orale : délinquance juvénile — Chinois 1ère A — Cours signé OnBuch+
% Programme officiel MINESEC. Compiler avec : tectonic ZHA07-oral-delinquance-juvenile.tex
\newcommand{\DOCMATIERE}{Chinois}
\newcommand{\DOCNIVEAU}{1ère A}
\newcommand{\DOCMODULE}{Module 1 : Vie familiale et intégration sociale}
\newcommand{\DOCLECON}{ZHA07}
\newcommand{\DOCTITRE}{Leçon 7 — Expression orale : délinquance juvénile}
% OnBuch+ — préambule « cours signés » ENRICHI (compilé avec tectonic / XeLaTeX)
% Système visuel « Pop » d'OnBuch+ : fond crème (jamais blanc), bordures encre
% épaisses, ombres « dures » décalées, titres Archivo Black en capitales.
% Version MATHÉMATIQUES : graphes (pgfplots/tikz), boîtes pédagogiques
% (définition, propriété, méthode, attention, le savais-tu, exercice résolu,
% exercice, TP, bilan), page de garde et signature OnBuch+ ancrée en bas.
%
% Macros attendues dans main.tex AVANT \input{preamble} :
%   \DOCMATIERE  \DOCNIVEAU  \DOCMODULE  \DOCLECON (numéro)  \DOCTITRE
\documentclass[11pt]{article}

\usepackage{fontspec}
\usepackage[french]{babel} % français = langue principale ; le chinois via \zh{..} / textebox (xeCJK)
\usepackage{xeCJK}
\usepackage{pifont} % \ding{51} \ding{55} utilisés par les modèles
\usepackage[a4paper,margin=2cm,top=3.4cm,bottom=2.8cm,headheight=1.4cm,footskip=1.3cm]{geometry}
\usepackage{amsmath,amssymb,amsfonts}
\usepackage{mathtools} % \xmapsto, \coloneqq... souvent utilisés par les modèles
\usepackage{cancel} % simplifications barrées (bilans de forces)
% \abs{z} (module, valeur absolue) et \norm{u} (norme), utilisés par les modèles
\providecommand{\abs}{}\let\abs\relax\DeclarePairedDelimiter{\abs}{\lvert}{\rvert}
\providecommand{\norm}{}\let\norm\relax\DeclarePairedDelimiter{\norm}{\lVert}{\rVert}
\usepackage[table]{xcolor} % [table] : fournit aussi \rowcolors (alternance de lignes)
\usepackage{pagecolor}
\usepackage{tikz}
\usetikzlibrary{arrows.meta,positioning,calc,shapes.geometric,shapes.misc,shapes.arrows,decorations.pathmorphing,decorations.pathreplacing,decorations.markings,patterns,shadows,mindmap,trees,fit,backgrounds,matrix,intersections,angles,quotes}
\usepackage{pgfplots}
\usepackage[version=4]{mhchem} % équations nucléaires \ce{^{235}_{92}U}
\usepackage[european,siunitx]{circuitikz} % schémas électriques
\pgfplotsset{compat=1.18}
\usepgfplotslibrary{fillbetween,groupplots} % aires entre courbes (calcul intégral)
\usepackage{enumitem}
\usepackage{fancyhdr}
% [most] charge toutes les bibliothèques tcolorbox (listings, minted, raster,
% documentation...) et gonfle inutilement la mémoire TeX sur un document long
% (~300 boîtes) : on ne charge que ce qui est réellement utilisé.
\usepackage[skins,breakable]{tcolorbox}
\usepackage{graphicx}
\usepackage{colortbl}
\usepackage{booktabs}
\usepackage{tabularx}
\usepackage{diagbox} % case d'en-tête diagonale des tableaux à double entrée
% Colonnes extensibles centrée / gauche / droite (C, L, R), souvent utilisées par les modèles
\newcolumntype{C}{>{\centering\arraybackslash}X}
\newcolumntype{L}{>{\raggedright\arraybackslash}X}
\newcolumntype{R}{>{\raggedleft\arraybackslash}X}
\usepackage{array}
\usepackage{multicol}
\usepackage{multirow}
\usepackage{siunitx}
% parse-numbers=false : accepte aussi \SI{2,5 \times 10^{-3}}{...}
\sisetup{locale=FR,output-decimal-marker={,},group-minimum-digits=4,parse-numbers=false}
\DeclareSIUnit\ohm{\ensuremath{\symup{\Omega}}} % U+2126 absent de la police
\DeclareSIUnit\division{div} % réglages d'oscilloscope (ms/div, V/div)
\DeclareSIUnit\tr{tr} % tours (vitesse de rotation, ex. \SI{1500}{\tr\per\minute})
\DeclareSIUnit\molar{M} % concentration molaire (ex. \SI{3}{\molar}, \SI{1}{\micro\molar})
\DeclareSIUnit\an{an} % année (le modèle confond parfois avec \year, un compteur TeX) — ex. \SI{5}{\kilo\meter\per\an}
\DeclareSIUnit\FCFA{FCFA} % franc CFA (ex. \SI{500}{\FCFA\per\kilo\gram})
\DeclareSIUnit\degreeBrix{\textdegree Brix} % teneur en sucre d'un sirop/jus (réfractométrie)
\DeclareSIUnit\month{mois} % le modèle utilise parfois \month, un compteur TeX, comme unité de durée
\usepackage{qrcode}
% \degree : commande fréquemment utilisée par les modèles pour un angle ou une
% température en dehors de \SI{}{} (non fournie par les paquets chargés).
\providecommand{\degree}{^{\circ}}
% \qed (fin de démonstration), utilisé par les modèles sans amsthm
\providecommand{\qed}{\hfill$\square$}
% \barre : double barre des valeurs interdites dans un tableau de variations (tkz-tab)
\providecommand{\barre}{\ensuremath{\|}}
% environnement proof (amsthm non chargé) utilisé par les modèles
\ifdefined\proof\else\newenvironment{proof}[1][Démonstration]{\par\noindent\textit{#1.}\ }{\qed\par}\fi
\usepackage{lastpage}
\usepackage{hyperref}
\hypersetup{hidelinks}

% ============================================================
% Palette « Pop » OnBuch+ — mêmes hex que la classe P (plus_ui.dart)
% ============================================================
\definecolor{poporange}{HTML}{F59321}
\definecolor{popdark}{HTML}{E07A0C}
\definecolor{popcream}{HTML}{FFF6E8}
\definecolor{popink}{HTML}{1C1714}
\definecolor{poppurple}{HTML}{7A5AE0}
\definecolor{popgreen}{HTML}{2E9E6A}
\definecolor{poppink}{HTML}{E0457B}
\definecolor{popblue}{HTML}{2F7DE1}
\definecolor{popgold}{HTML}{C98A12}
\definecolor{popmuted}{HTML}{8A7F76}
\definecolor{popline}{HTML}{EADFD2}
\definecolor{popgray}{HTML}{8A7F76}
\colorlet{popgrey}{popgray}
% Alias tolérants (le modèle nomme parfois les couleurs différemment)
\colorlet{popwhite}{white}
\colorlet{popblack}{popink}
\colorlet{popred}{poppink}
\colorlet{popyellow}{popgold}
% Teintes claires (fonds de tableaux/encadrés) — le modèle nomme parfois
% les couleurs avec un suffixe L (ex. popblueL) sans les avoir définies.
\colorlet{poporangeL}{poporange!20}
\colorlet{popdarkL}{popdark!20}
\colorlet{poppurpleL}{poppurple!20}
\colorlet{popgreenL}{popgreen!20}
\colorlet{poppinkL}{poppink!20}
\colorlet{popblueL}{popblue!20}
\colorlet{popgoldL}{popgold!20}
\colorlet{popredL}{popred!20}
\colorlet{popgrayL}{popgray!20}
% Teintes claires pour fonds de tableaux / zones
\colorlet{popblueL}{popblue!10!white}
\colorlet{poporangeL}{poporange!14!white}
\colorlet{popgreenL}{popgreen!12!white}
\colorlet{poppurpleL}{poppurple!12!white}
\colorlet{poppinkL}{poppink!12!white}
\colorlet{popgoldL}{popgold!14!white}

\pagecolor{popcream}

% ============================================================
% Typographie
% ============================================================
\defaultfontfeatures{Ligatures=TeX}
\setmainfont{PlusJakartaSans-Medium.ttf}[Path=../../../fonts/,BoldFont=PlusJakartaSans-Bold.ttf]
% Symboles, API et emojis absents de la police principale : repli sur DejaVu Sans (caractères actifs)
\newfontface\symfont{DejaVuSans.ttf}[Path=../../../fonts/]
\makeatletter
\newcommand\symactive[1]{\catcode#1=13 \begingroup\lccode`\~=#1 \lowercase{\endgroup\def~}{{\symfont\char#1}}}
\newcommand\symblank[1]{\catcode#1=13 \begingroup\lccode`\~=#1 \lowercase{\endgroup\def~}{}}
\newcommand\symhyph[1]{\catcode#1=13 \begingroup\lccode`\~=#1 \lowercase{\endgroup\def~}{\mbox{-}}}
\newcount\symc
\newcommand\symrange[2]{\symc=#1 \@whilenum\symc<#2 \do{\expandafter\symactive\expandafter{\the\symc}\advance\symc by 1 }}
\newcommand\symblankrange[2]{\symc=#1 \@whilenum\symc<#2 \do{\expandafter\symblank\expandafter{\the\symc}\advance\symc by 1 }}
\makeatother
\symrange{"0250}{"02FF}\symactive{"2713}\symactive{"2714}\symactive{"2717}\symactive{"2718}\symactive{"2610}\symactive{"2611}\symactive{"2612}
\symrange{"2600}{"27BF}
\catcode"2705=13 \begingroup\lccode`\~="2705 \lowercase{\endgroup\def~}{{\symfont\char"2713}}\symblank{"26BD}\symblank{"26A1}
\symrange{"2460}{"257F}\symactive{"223C}\symactive{"2248}\symactive{"2260}
\symhyph{"2011}\symhyph{"2E1A}\symblankrange{"1F300}{"1FAFF}\symblank{"FE0F}
% Caractères chinois : WenQuanYi Zen Hei (sous-ensemble GB2312 + pinyin), gras/italique simulés
\setCJKmainfont{WenQuanYiZenHei-subset.ttf}[Path=../../../fonts/,AutoFakeBold=3,AutoFakeSlant=0.2]
\xeCJKsetup{CJKspace=false}
% Pinyin : voyelles à caron (ǎ ǐ ǒ ǔ ǖ ǘ ǚ ǜ) absentes de la police principale -> caractères actifs
\newfontface\pyfont{WenQuanYiZenHei-subset.ttf}[Path=../../../fonts/]
\newcommand\pyactive[1]{\catcode#1=13 \begingroup\lccode`\~=#1 \lowercase{\endgroup\def~}{{\pyfont\char#1}}}
\pyactive{"01CD}\pyactive{"01CE}\pyactive{"01CF}\pyactive{"01D0}\pyactive{"01D1}\pyactive{"01D2}\pyactive{"01D3}\pyactive{"01D4}\pyactive{"01D5}\pyactive{"01D6}\pyactive{"01D7}\pyactive{"01D8}\pyactive{"01D9}\pyactive{"01DA}\pyactive{"01DB}\pyactive{"01DC}
% Maths en sans-serif (Fira Math) avec chiffres et lettres droites de Plus
% Jakarta, pour que les formules se fondent dans le texte.
\usepackage[math-style=ISO,bold-style=ISO]{unicode-math}
\setmathfont{FiraMath-Regular.otf}[Path=../../../fonts/]
\setmathfont{PlusJakartaSans-Medium.ttf}[Path=../../../fonts/,range={up/num,up/{Latin,latin},\mathup}]
\usepackage{icomma} % virgule décimale sans espace en mode math
% Caractères mathématiques Unicode absents des polices de texte (Plus Jakarta,
% Archivo Black) : rendus par leur équivalent mathématique, en texte comme en
% formule — sinon ils s'affichent en carré vide (ex. « Arithmétique dans ℤ »).
\usepackage{newunicodechar}
\newunicodechar{ℝ}{\ensuremath{\mathbb{R}}}
\newunicodechar{ℤ}{\ensuremath{\mathbb{Z}}}
\newunicodechar{ℂ}{\ensuremath{\mathbb{C}}}
\newunicodechar{ℕ}{\ensuremath{\mathbb{N}}}
\newunicodechar{ℚ}{\ensuremath{\mathbb{Q}}}
\newunicodechar{𝔻}{\ensuremath{\mathbb{D}}}
\newunicodechar{α}{\ensuremath{\alpha}}
\newunicodechar{β}{\ensuremath{\beta}}
\newunicodechar{γ}{\ensuremath{\gamma}}
\newunicodechar{δ}{\ensuremath{\delta}}
\newunicodechar{ε}{\ensuremath{\varepsilon}}
\newunicodechar{θ}{\ensuremath{\theta}}
\newunicodechar{λ}{\ensuremath{\lambda}}
\newunicodechar{μ}{\ensuremath{\mu}}
\newunicodechar{σ}{\ensuremath{\sigma}}
\newunicodechar{τ}{\ensuremath{\tau}}
\newunicodechar{φ}{\ensuremath{\varphi}}
\newunicodechar{ψ}{\ensuremath{\psi}}
\newunicodechar{ω}{\ensuremath{\omega}}
\newunicodechar{Σ}{\ensuremath{\Sigma}}
% majuscules produites par \MakeUppercase dans les titres (α-aminés -> Α)
\newunicodechar{Α}{\ensuremath{\alpha}}
\newunicodechar{Β}{\ensuremath{\beta}}
\newunicodechar{Γ}{\ensuremath{\Gamma}}
\newunicodechar{Δ}{\ensuremath{\Delta}}
\newunicodechar{Ε}{\ensuremath{\varepsilon}}
\newunicodechar{Θ}{\ensuremath{\Theta}}
\newunicodechar{Λ}{\ensuremath{\Lambda}}
\newunicodechar{Μ}{\ensuremath{\mu}}
\newunicodechar{Τ}{\ensuremath{\tau}}
\newunicodechar{Φ}{\ensuremath{\Phi}}
\newunicodechar{Ψ}{\ensuremath{\Psi}}
\newunicodechar{Ω}{\ensuremath{\Omega}}
\newunicodechar{↦}{\ensuremath{\mapsto}}
\newunicodechar{∈}{\ensuremath{\in}}
\newunicodechar{∉}{\ensuremath{\notin}}
\newunicodechar{∧}{\ensuremath{\wedge}}
\newunicodechar{⇔}{\ensuremath{\Leftrightarrow}}
\newunicodechar{⟺}{\ensuremath{\Longleftrightarrow}}
\newunicodechar{⇒}{\ensuremath{\Rightarrow}}
\newunicodechar{⟹}{\ensuremath{\Longrightarrow}}
\newunicodechar{∘}{\ensuremath{\circ}}
\newunicodechar{∪}{\ensuremath{\cup}}
\newunicodechar{∩}{\ensuremath{\cap}}
\newunicodechar{≡}{\ensuremath{\equiv}}
\newunicodechar{⊂}{\ensuremath{\subset}}
\newunicodechar{⊥}{\ensuremath{\perp}}
\newunicodechar{∃}{\ensuremath{\exists}}
\newunicodechar{∀}{\ensuremath{\forall}}
\newunicodechar{∛}{\ensuremath{\sqrt[3]{\,}}}
\newunicodechar{∜}{\ensuremath{\sqrt[4]{\,}}}
\newunicodechar{⃗}{\ensuremath{{}^{\to}}}
\newunicodechar{ⁿ}{\ifmmode^{n}\else\textsuperscript{\ensuremath{n}}\fi}
\newunicodechar{ˣ}{\ifmmode^{x}\else\textsuperscript{\ensuremath{x}}\fi}
\newunicodechar{ᵇ}{\ifmmode^{b}\else\textsuperscript{\ensuremath{b}}\fi}
\newunicodechar{ʳ}{\ifmmode^{r}\else\textsuperscript{\ensuremath{r}}\fi}
\newunicodechar{ᵘ}{\ifmmode^{u}\else\textsuperscript{\ensuremath{u}}\fi}
\newunicodechar{ʸ}{\ifmmode^{y}\else\textsuperscript{\ensuremath{y}}\fi}
\newunicodechar{ᵃ}{\ifmmode^{a}\else\textsuperscript{\ensuremath{a}}\fi}
\newunicodechar{ᵉ}{\ifmmode^{e}\else\textsuperscript{\ensuremath{e}}\fi}
\newunicodechar{ᵏ}{\ifmmode^{k}\else\textsuperscript{\ensuremath{k}}\fi}
\newunicodechar{ᵖ}{\ifmmode^{p}\else\textsuperscript{\ensuremath{p}}\fi}
\newunicodechar{ᴺ}{\ifmmode^{N}\else\textsuperscript{\ensuremath{N}}\fi}
\newunicodechar{ᶜ}{\ifmmode^{c}\else\textsuperscript{\ensuremath{c}}\fi}
\newunicodechar{ʰ}{\ifmmode^{h}\else\textsuperscript{\ensuremath{h}}\fi}
\newunicodechar{ᵠ}{\ifmmode^{q}\else\textsuperscript{\ensuremath{q}}\fi}
\newunicodechar{ₙ}{\ifmmode_{n}\else\textsubscript{\ensuremath{n}}\fi}
\newunicodechar{ₐ}{\ifmmode_{a}\else\textsubscript{\ensuremath{a}}\fi}
\newunicodechar{ᵢ}{\ifmmode_{i}\else\textsubscript{\ensuremath{i}}\fi}
\newunicodechar{ᵣ}{\ifmmode_{r}\else\textsubscript{\ensuremath{r}}\fi}
\newunicodechar{ₖ}{\ifmmode_{k}\else\textsubscript{\ensuremath{k}}\fi}
\newunicodechar{ᵦ}{\ifmmode_{\beta}\else\textsubscript{\ensuremath{\beta}}\fi}
\newunicodechar{ₓ}{\ifmmode_{x}\else\textsubscript{\ensuremath{x}}\fi}
\newfontfamily\popdisplay{ArchivoBlack-Regular.ttf}[Path=../../../fonts/]
\newfontfamily\poplogo{SpaceGrotesk-Bold.ttf}[Path=../../../fonts/]
\newfontfamily\popsemi{PlusJakartaSans-SemiBold.ttf}[Path=../../../fonts/]
\newfontfamily\popxbold{PlusJakartaSans-ExtraBold.ttf}[Path=../../../fonts/]
\newfontfamily\popxboldls{PlusJakartaSans-ExtraBold.ttf}[Path=../../../fonts/,LetterSpace=3.0]

\setlist[enumerate]{leftmargin=1.7em,itemsep=5pt,topsep=4pt}
\setlist[itemize]{leftmargin=1.7em,itemsep=4pt,topsep=4pt,label={\color{poporange}\textbullet}}
\setlength{\parindent}{0pt}
\setlength{\parskip}{5pt}
\renewcommand{\arraystretch}{1.25}

% pgfplots : style Pop par défaut
\pgfplotsset{
  every axis/.append style={axis line style={popink,line width=1pt},tick style={popink},
    grid style={popline,line width=0.5pt},label style={font=\small},tick label style={font=\footnotesize}},
  popcurve/.style={poporange,line width=1.8pt,no markers,smooth},
}
% Même style disponible dans un \draw TikZ simple (hors pgfplots)
\tikzset{popcurve/.style={poporange,line width=1.8pt}}
\tikzset{popgrid/.style={popline,very thin}}
\colorlet{popcurve}{poporange} % le nom du style est parfois utilisé comme couleur

% ============================================================
% Pill / squiggle
% ============================================================
\newcommand{\poppill}[3]{%
  \tikz[baseline=-0.55ex]\node[draw=popink,line width=1.2pt,rounded corners=2.6pt,%
    fill=#2,inner xsep=6.5pt,inner ysep=3pt]{%
    \popxboldls\fontsize{7}{7}\selectfont\color{#3}\MakeUppercase{#1}};%
}
\newcommand{\popsquiggle}[1]{%
  \tikz[baseline=-0.4ex]{
    \draw[#1,line width=2.6pt,line cap=round]
      plot [smooth] coordinates {(0,0.09) (0.34,0.01) (0.68,0.21) (1.02,0.01) (1.36,0.10)};
  }%
}
% Mot-clé surligné (marqueur orange sous le texte)
\newcommand{\cle}[1]{{\popxbold\color{popdark}#1}}

% ============================================================
% En-tête / pied
% ============================================================
\pagestyle{fancy}
\fancyhf{}
\renewcommand{\headrulewidth}{0pt}
\renewcommand{\footrulewidth}{0pt}
\lhead{\raisebox{-0.15ex}{\poplogo\fontsize{13}{13}\selectfont\color{popink}OnBuch\color{poporange}+}}
\rhead{\poppill{\DOCMATIERE\ \textbullet\ \DOCNIVEAU\ \textbullet\ Leçon \DOCLECON}{white}{popink}}
\lfoot{\popxbold\fontsize{7.6}{7.6}\selectfont\color{popmuted}\MakeUppercase{Cours signé OnBuch+}\ \textbullet\ {\popsemi\DOCTITRE}}
\rfoot{\tikz[baseline=-0.6ex]\node[rounded corners=8.5pt,fill=popink,minimum height=17pt,minimum width=17pt,inner xsep=5pt,inner sep=0]{\popxbold\fontsize{8}{8}\selectfont\color{popcream}\thepage\,/\,\pageref*{LastPage}};}

% ============================================================
% Titres
% ============================================================
\newcommand{\chaptitre}[2]{%
  \begin{tcolorbox}[enhanced,colback=poporange,colframe=popink,boxrule=2.6pt,arc=11pt,
      drop shadow=popink,left=15pt,right=15pt,top=12pt,bottom=11pt,before skip=3pt,after skip=18pt]
    \poppill{#2}{popink}{popcream}\par\vspace{9pt}
    {\raggedright\hyphenpenalty=10000\exhyphenpenalty=10000\popdisplay\fontsize{22}{24}\selectfont\color{popcream}\MakeUppercase{#1}\par}
  \end{tcolorbox}
}
% Section numérotée : pastille orange avec le numéro + titre Archivo
\newcounter{popsec}
\newcommand{\coursec}[1]{%
  \stepcounter{popsec}\setcounter{popsub}{0}%
  \par\vspace{16pt}\noindent
  \begin{minipage}{\linewidth}
  \tikz[baseline=-0.7ex]\node[draw=popink,line width=1.6pt,fill=poporange,rounded corners=4pt,minimum size=22pt,inner sep=0]{\popdisplay\fontsize{12}{12}\selectfont\color{popcream}\thepopsec};%
  \hspace{8pt}{\popdisplay\fontsize{15}{17}\selectfont\color{popink}\MakeUppercase{#1}}\par
  \vspace{4pt}\noindent{\color{poporange}\rule{36pt}{3.6pt}}
  \end{minipage}\par\nopagebreak\vspace{8pt}
}
\newcounter{popsub}
\newcommand{\courssub}[1]{%
  \stepcounter{popsub}%
  \par\vspace{9pt}\noindent{\popxbold\fontsize{11.5}{13}\selectfont\color{popink}{\color{poporange}\thepopsec.\thepopsub}\hspace{6pt}#1}\par\nopagebreak\vspace{4pt}
}

% ============================================================
% Boîtes pédagogiques — même recette : fond blanc, bordure épaisse colorée,
% ombre dure encre, badge plein. Titre optionnel [#1].
% ============================================================
\tcbset{popbase/.style={breakable,enhanced,colback=white,boxrule=2.3pt,arc=9pt,
  shadow={3pt}{-3pt}{0pt}{popink},left=14pt,right=14pt,top=11pt,bottom=11pt,before skip=11pt,after skip=15pt,
  fontupper=\popsemi\fontsize{10.6}{14.5}\selectfont\color{popink},
  fontlower=\popsemi\fontsize{10.6}{14.5}\selectfont\color{popink}}}
% Coche dessinée (le glyphe ✓ n'existe pas dans les polices du document)
\newcommand{\popcheck}[1]{\tikz[baseline=-0.5ex]\draw[#1,line width=1.6pt,line cap=round,line join=round](-0.13,0.0)--(-0.04,-0.09)--(0.13,0.1);}
\providecommand{\coche}{\popcheck{popgreen}}
\providecommand{\resultat}[1]{\par\smallskip\noindent\fcolorbox{popgreen}{popgreenL}{\parbox{\dimexpr\linewidth-2\fboxsep-2\fboxrule\relax}{\textbf{Résultat.} #1}}\par\smallskip}
\newcommand{\popboxhead}[3]{\poppill{#1}{#2}{white}\ifx\relax#3\relax\else\hspace{6pt}{\popxbold\fontsize{10.6}{12}\selectfont\color{popink}#3}\fi\par\vspace{7pt}}

\newtcolorbox{aretenir}[1][]{popbase,colframe=popgreen,before upper={\popboxhead{À retenir}{popgreen}{#1}}}
% Texte en chinois (document de lecture, dialogue, extrait) : cadre
% d'étude. Un titre optionnel (source, auteur) s'affiche dans l'en-tête.
\newtcolorbox{textebox}[1][]{popbase,colframe=popblue,before upper={\popboxhead{文本 · Texte}{popblue}{#1}},after upper={}}
% Marques ✓ / ✗ (forme correcte / fautive) souvent utilisées par les modèles
\providecommand{\cross}{\textcolor{poppink}{\ensuremath{\times}}}
\providecommand{\xmark}{\cross}\providecommand{\crossmark}{\cross}\providecommand{\cmark}{\ensuremath{\checkmark}}
% Ligne de réponse / justification à compléter (inventée par les modèles)
\providecommand{\justification}[1]{\par\noindent\textit{Justification~:} \dotfill\par}
% \textipa (package tipa non chargé) : la police ne couvre pas l'API -> simple passage
\providecommand{\textipa}[1]{#1}
% Soulignement (ulem non chargé) : \uline, \ul
\providecommand{\uline}[1]{\underline{#1}}\providecommand{\ul}[1]{\underline{#1}}
% \glossaire{\item ...} inventée par les modèles : liste de glossaire
\providecommand{\glossaire}[1]{\begin{itemize}#1\end{itemize}}
% \alert (beamer) inventée par les modèles : texte mis en évidence
\providecommand{\alert}[1]{\textbf{\textcolor{poppink}{#1}}}
% variantes ASCII d'ordinaux écrites en macro
\providecommand{\iere}{\textsuperscript{re}}\providecommand{\ieme}{\textsuperscript{e}}\providecommand{\ere}{\textsuperscript{re}}\providecommand{\eme}{\textsuperscript{e}}\providecommand{\er}{\textsuperscript{er}}\providecommand{\ie}{\textsuperscript{e}}
% Ordinaux français écrits en macro par les modèles : 1\ière, 3\ième
\newcommand{\ière}{\textsuperscript{re}}\newcommand{\ième}{\textsuperscript{e}}
% \exemple (inventée par les modèles) : étiquette « Exemple : »
\providecommand{\exemple}{\textbf{Exemple~:}\ }
% \square utilisable aussi hors mode math (cases à cocher)
\AtBeginDocument{\let\squareMath\square\renewcommand{\square}{\ensuremath{\squareMath}}}
% Mot ou phrase en chinois à l'intérieur d'un paragraphe français
\newcommand{\zh}[1]{\ifmmode\text{#1}\else{#1}\fi}
\let\eng\zh \let\esp\zh \let\all\zh % alias tolérés
% flèches d'intonation inventées par les modèles
\providecommand{\textsearrow}{}\renewcommand{\textsearrow}{\ensuremath{\searrow}}\providecommand{\textnearrow}{}\renewcommand{\textnearrow}{\ensuremath{\nearrow}}
% \fr{..} : mot français (inventé par les modèles)
\providecommand{\fr}[1]{\textit{#1}}
% zhbox : encadré de texte chinois inventé par les modèles
\newenvironment{zhbox}{\begin{quote}}{\end{quote}}
% \vs : « versus » inventé par les modèles
\providecommand{\vs}{\textit{vs}\ }
% \blank : case à compléter inventée par les modèles
\providecommand{\blank}[1][]{\underline{\hspace{1.6em}}}
\newtcolorbox{exemplebox}[1][]{popbase,colframe=poporange,before upper={\popboxhead{Exemple}{poporange}{#1}}}
\newtcolorbox{definition}[1][]{popbase,colframe=popblue,before upper={\popboxhead{Définition}{popblue}{#1}}}
\newtcolorbox{propriete}[1][]{popbase,colframe=poppurple,before upper={\popboxhead{Propriété}{poppurple}{#1}}}
\newtcolorbox{methode}[1][]{popbase,colframe=popgold,before upper={\popboxhead{Méthode}{popgold}{#1}}}
\newtcolorbox{attention}[1][]{popbase,colframe=poppink,before upper={\popboxhead{Attention}{poppink}{#1}}}
\newtcolorbox{experience}[1][]{popbase,colframe=popblue,colback=popblueL,before upper={\popboxhead{Expérience}{popblue}{#1}}}
% « Le savais-tu ? » : carte violette pleine (contexte, culture, Cameroun)
\newtcolorbox{savaistu}[1][]{popbase,colback=poppurple,colframe=popink,
  fontupper=\popsemi\fontsize{10.4}{14.2}\selectfont\color{white},
  before upper={\poppill{Le savais-tu\,?}{white}{poppurple}\ifx\relax#1\relax\else\hspace{6pt}{\popxbold\color{white}#1}\fi\par\vspace{7pt}}}
% Exercice résolu : énoncé en haut, \tcblower puis la solution sur fond crème
\newcounter{popexo}\newcounter{popexr}
\newtcolorbox{exoresolu}[1][]{popbase,colframe=poporange,colbacklower=popcream,
  segmentation style={popink,line width=1.2pt,dashed},
  before upper={\stepcounter{popexr}\popboxhead{Exercice résolu \thepopexr}{poporange}{#1}},
  before lower={\poppill{Solution}{popink}{popcream}\par\vspace{6pt}}}
% Exercice (non résolu en place) — la correction est dans la partie Corrigés
\newtcolorbox{exercice}[1][]{popbase,colframe=popink,
  before upper={\stepcounter{popexo}\popboxhead{Exercice \thepopexo}{popink}{#1}}}
% Corrigé
\newtcolorbox{corrige}[1][]{popbase,colframe=popgreen,colback=popgreenL,
  before upper={\popboxhead{Corrigé}{popgreen}{#1}}}
% Objectifs / prérequis / bilan
\newtcolorbox{objectifs}{popbase,colframe=popink,colback=poporangeL,
  before upper={\popboxhead{Objectifs de la leçon}{popink}{}}}
\newtcolorbox{prerequis}{popbase,colframe=popink,colback=popblueL,
  before upper={\popboxhead{Prérequis}{popblue}{}}}
\newtcolorbox{bilan}{popbase,colframe=popink,colback=white,boxrule=2.8pt,
  before upper={\popboxhead{Fiche bilan}{poporange}{L'essentiel à connaître pour l'examen}}}
% Figure encadrée avec légende
\newtcolorbox{popfigure}[1][]{enhanced,breakable=false,colback=white,colframe=popline,boxrule=1.4pt,arc=8pt,
  left=10pt,right=10pt,top=10pt,bottom=8pt,before skip=10pt,after skip=12pt,halign=center,
  lower separated=false,halign lower=center,
  fontlower=\popsemi\fontsize{9}{11.5}\selectfont\color{popmuted},
  #1}
\newcommand{\legende}[1]{\par\vspace{6pt}{\popsemi\fontsize{9}{11.5}\selectfont\color{popmuted}#1}}

% ============================================================
% Signature OnBuch+
% ============================================================
\newcommand{\poplogoinline}[1]{{\poplogo\fontsize{#1}{#1}\selectfont\color{popink}OnBuch\color{poporange}+}}
\newcommand{\signatureonbuch}{%
  \begin{tcolorbox}[enhanced,colback=popink,colframe=popink,boxrule=2.2pt,arc=10pt,
      drop shadow=poporange,left=14pt,right=14pt,top=10pt,bottom=10pt,before skip=0pt,after skip=0pt]
    \tikz[baseline=-0.7ex]\node[circle,fill=popgreen,draw=popcream,line width=1.3pt,minimum size=22pt,inner sep=0]{\popcheck{white}};%
    \hspace{9pt}\begin{minipage}[c]{0.62\linewidth}
      {\popxbold\fontsize{12}{13}\selectfont\color{popcream}Signé\ \textbullet\ {\poplogo OnBuch\color{poporange}+}}\par\vspace{2pt}
      {\raggedright\popsemi\fontsize{8}{10.5}\selectfont\color{popline}\MakeUppercase{Équipe pédagogique OnBuch+}\\\MakeUppercase{Conforme au programme officiel MINESEC}\par}
    \end{minipage}\hfill
    {\poplogo\fontsize{22}{22}\selectfont\color{popcream}OnBuch\color{poporange}+}
  \end{tcolorbox}%
}

% ============================================================
% Page de garde
% #1 titre · #2 matière · #3 niveau · #4 module · #5 n° leçon · #6 durée ·
% #7 description / objectifs (texte ou itemize)
% ============================================================
\newcommand{\pagegarde}[7]{%
  \thispagestyle{empty}
  \newgeometry{a4paper,margin=1.7cm,top=1.6cm,bottom=1.4cm}%
  \begin{tikzpicture}[remember picture,overlay]
    \fill[poporange!18] ([xshift=-3.2cm,yshift=-3.6cm]current page.north east) circle (4.6cm);
    \fill[poppurple!14] ([xshift=2.4cm,yshift=7.6cm]current page.south west) circle (3.8cm);
  \end{tikzpicture}%
  {\poplogo\fontsize{30}{30}\selectfont\color{popink}OnBuch\color{poporange}+}\hfill
  \raisebox{0.6ex}{\poppill{Cours signé \textbullet\ Premium}{popink}{popcream}}\par
  \vspace{1pt}\popsquiggle{poppurple}\par
  \vspace{8pt}
  \begin{tcolorbox}[enhanced,colback=poporange,colframe=popink,boxrule=2.8pt,arc=13pt,
      drop shadow=popink,left=18pt,right=18pt,top=12pt,bottom=13pt,after skip=10pt]
    \poppill{#2 \textbullet\ #3}{popink}{popcream}\hspace{5pt}\poppill{#4}{white}{popink}\par\vspace{8pt}
    {\popdisplay\fontsize{14}{14}\selectfont\color{popink}LEÇON #5}\par\vspace{4pt}
    {\raggedright\hyphenpenalty=10000\exhyphenpenalty=10000\popdisplay\fontsize{25}{27}\selectfont\color{popcream}\MakeUppercase{#1}\par}
    \vspace{8pt}
    {\popxbold\fontsize{9.5}{9.5}\selectfont\color{popink}\MakeUppercase{Durée conseillée : #6}}
  \end{tcolorbox}
  \begin{tcolorbox}[enhanced,colback=white,colframe=popink,boxrule=2.2pt,arc=10pt,
      drop shadow=popink,left=16pt,right=16pt,top=10pt,bottom=10pt,after skip=8pt]
    \poppill{Dans cette leçon}{poppurple}{white}\par\vspace{5pt}
    \popsemi\fontsize{9.3}{12.6}\selectfont\color{popink}#7
  \end{tcolorbox}
  \noindent\begin{minipage}[t]{0.487\linewidth}
    \begin{tcolorbox}[enhanced,colback=popgreen,colframe=popink,boxrule=2.2pt,arc=10pt,
        drop shadow=popink,left=11pt,right=11pt,top=10pt,bottom=10pt,height=3.1cm,valign=center]
      \tcbox[colback=white,colframe=popink,boxrule=1.3pt,arc=5pt,boxsep=0pt,nobeforeafter,
        left=3pt,right=3pt,top=3pt,bottom=3pt,valign=center]{\qrcode[height=1.9cm]{https://play.google.com/store/apps/details?id=com.luvvix.onbuch&hl=fr}}%
      \hspace{8pt}\begin{minipage}[c]{0.5\linewidth}
        {\popdisplay\fontsize{11}{12.5}\selectfont\color{white}\MakeUppercase{Télécharge OnBuch}}\par\vspace{4pt}
        {\raggedright\popsemi\fontsize{8.4}{11}\selectfont\color{white}Gratuit \textbullet\ Google Play\\Tuteur IA, annales, concours, résultats\par}
      \end{minipage}
    \end{tcolorbox}
  \end{minipage}\hfill
  \begin{minipage}[t]{0.487\linewidth}
    \begin{tcolorbox}[enhanced,colback=poppurple,colframe=popink,boxrule=2.2pt,arc=10pt,
        drop shadow=popink,left=11pt,right=11pt,top=10pt,bottom=10pt,height=3.1cm,valign=center]
      \tikz[baseline=-0.7ex]\node[circle,fill=white,draw=popink,line width=1.3pt,minimum size=1.6cm,inner sep=0]{\popdisplay\fontsize{24}{24}\selectfont\color{poppurple}+};%
      \hspace{8pt}\begin{minipage}[c]{0.6\linewidth}
        {\popdisplay\fontsize{11}{12.5}\selectfont\color{white}\MakeUppercase{Passe à OnBuch+}}\par\vspace{4pt}
        {\raggedright\popsemi\fontsize{8.4}{11}\selectfont\color{white}Tous les cours signés, Tuteur IA illimité, fiches PDF — onglet OnBuch+ de l'app\par}
      \end{minipage}
    \end{tcolorbox}
  \end{minipage}
  \vspace{8pt}
  \popcontenu
  \vfill
  \signatureonbuch
  \restoregeometry
  \newpage
}

% Rangée « Ce cours contient » (page de garde)
\newcommand{\poptile}[3]{% #1 couleur #2 gros texte #3 légende
  \begin{tcolorbox}[enhanced,colback=#1,colframe=popink,boxrule=2pt,arc=9pt,shadow={2.5pt}{-2.5pt}{0pt}{popink},
      left=6pt,right=6pt,top=9pt,bottom=9pt,halign=center,valign=center,height=1.75cm,width=\linewidth,nobeforeafter]
    {\popdisplay\fontsize{12.5}{13}\selectfont\color{white}\MakeUppercase{#2}}\par\vspace{3pt}
    {\centering\popsemi\fontsize{7.8}{9.5}\selectfont\color{white}#3\par}
  \end{tcolorbox}}
\newcommand{\popcontenu}{%
  \par\noindent{\popxboldls\fontsize{7.5}{7.5}\selectfont\color{popmuted}CE COURS SIGNÉ CONTIENT}\par\vspace{7pt}
  \noindent\begin{minipage}[t]{0.235\linewidth}\poptile{popblue}{Cours}{complet, illustré, pas à pas}\end{minipage}\hfill
  \begin{minipage}[t]{0.235\linewidth}\poptile{poporange}{Méthodes}{et exercices résolus}\end{minipage}\hfill
  \begin{minipage}[t]{0.235\linewidth}\poptile{poppink}{Exercices}{type examen + corrigés}\end{minipage}\hfill
  \begin{minipage}[t]{0.235\linewidth}\poptile{popgreen}{Fiche bilan}{l'essentiel à retenir}\end{minipage}\par}

% Sommaire
\newtcolorbox{sommaire}{enhanced,colback=white,colframe=popink,boxrule=2.2pt,arc=10pt,
  drop shadow=popink,left=18pt,right=18pt,top=13pt,bottom=13pt,before skip=6pt,after skip=20pt,
  before upper={\poppill{Sommaire}{poppurple}{white}\par\vspace{9pt}\popsemi\fontsize{10.6}{14.5}\selectfont\color{popink}}}

% Signature de fin de cours — poussée tout en bas de la dernière page
\newcommand{\findecours}{%
  \par\vfill\nopagebreak
  \begin{center}\popsquiggle{poporange}\end{center}\vspace{4pt}
  \signatureonbuch
}
\AtBeginDocument{\def\llbracket{\mathopen{[\mkern-3.5mu[}}\def\rrbracket{\mathclose{]\mkern-3.5mu]}}} % intervalles d'entiers (glyphe ⟦ absent de la police)
% Glyphes absents de Fira Math : redéfinis à partir de glyphes disponibles
\AtBeginDocument{%
  \def\setminus{\mathbin{\backslash}}\let\smallsetminus\setminus
  \def\vdots{\mathord{\vbox{\baselineskip3.2pt\lineskiplimit0pt\kern4pt\hbox{.}\hbox{.}\hbox{.}}}}%
  \def\ddots{\mathinner{\mkern1mu\raise6.5pt\hbox{.}\mkern2mu\raise3.6pt\hbox{.}\mkern2mu\raise0.7pt\hbox{.}\mkern1mu}}%
  \def\triangle{\mathord{\scalebox{0.9}{$\symup{\Delta}$}}}%
  \def\nabla{\mathord{\raisebox{\depth}{\scalebox{1}[-1]{$\symup{\Delta}$}}}}%
  \def\checkmark{\popcheck{popdark}}%
  \def\star{\mathbin{\ast}}\def\bigstar{\mathord{\ast}}%
  \def\diamond{\mathbin{\scalebox{0.6}{$\Diamond$}}}%
  \def\bigcup{\mathop{\mathchoice{\raisebox{-0.3em}{\scalebox{1.7}{$\cup$}}}{\scalebox{1.25}{$\cup$}}{\cup}{\cup}}\displaylimits}%
  \def\bigcap{\mathop{\mathchoice{\raisebox{-0.3em}{\scalebox{1.7}{$\cap$}}}{\scalebox{1.25}{$\cap$}}{\cap}{\cap}}\displaylimits}%
  \def\lesseqqgtr{\mathrel{\lessgtr}}\def\bigoplus{\mathop{\oplus}\displaylimits}%
}
\providecommand{\ssi}{si et seulement si} % abréviation fréquente
\AtBeginDocument{\providecommand{\wideparen}{\overparen}} % arc de cercle

%% FIGFIX-BEGIN v3 — correctifs globaux des figures (docs/onbuch-generation/tools/figfix)
\usepackage{adjustbox}\usepackage{etoolbox}
% toute figure TikZ/pgfplots plus large que la ligne est réduite à la largeur disponible
\BeforeBeginEnvironment{tikzpicture}{\begin{adjustbox}{max width=\linewidth}}
\AfterEndEnvironment{tikzpicture}{\end{adjustbox}}
% légendes pgfplots lisibles par-dessus les courbes
\pgfplotsset{every axis/.append style={legend style={fill=white,fill opacity=0.92,text opacity=1,draw=black!15,inner sep=3pt}}}
% étiquettes d'arêtes (syntaxe edge["..."]) avec fond blanc
\tikzset{every edge quotes/.append style={fill=white,inner sep=1.2pt}}
% bulles de cartes mentales : texte à la ligne dans la bulle, bulle un peu plus grande
\tikzset{every concept/.append style={text width=2.0cm,align=center,minimum size=2.3cm,inner sep=2pt}}
%% FIGFIX-END
\begin{document}

\pagegarde{Leçon 7 — Expression orale : délinquance juvénile}{Chinois}{1ère A}{Module 1 : Vie familiale et intégration sociale}{ZHA07}{2 h}{Leçon d'expression orale consacrée au thème de la délinquance juvénile. À travers des dialogues modèles, des formules fonctionnelles et des débats guidés, les élèves apprennent à présenter un exposé simple, à donner leur avis et à échanger en chinois avec une prononciation correcte des tons.
\vspace{4pt}
\begin{multicols}{2}\small
\begin{itemize}
\item À la fin de cette leçon, je sais utiliser le vocabulaire de base lié à la délinquance juvénile (犯罪, 青少年, 法律, 教育...)
\item je sais présenter oralement un exposé simple et structuré en chinois sur ce thème
\item je sais donner mon avis avec les formules 我认为…, 我觉得…, 在我看来…
\item je sais exprimer l'accord et le désaccord poliment (我同意/我不同意, 你说得对…)
\item je sais relancer un débat en posant des questions (你觉得呢？为什么呢？)
\item je sais prononcer correctement les tons des mots clés de la leçon
\end{itemize}\end{multicols}}

\begin{sommaire}
\begin{multicols}{2}
\begin{enumerate}
\item Vocabulaire du thème : la délinquance juvénile
\item Dialogue modèle 1 : parler du problème
\item Formules fonctionnelles : présenter, opiner, relancer
\item Dialogue modèle 2 : le débat — punir ou éduquer ?
\item Phonétique : les tons à travailler
\item Jeux de rôle et débat guidé
\item Activité d'intégration
\item Méthodes et exercices résolus
\item Exercices
\item Corrigés des exercices
\item Fiche bilan
\end{enumerate}
\end{multicols}
\end{sommaire}

\begin{objectifs}
\begin{itemize}
\item À la fin de cette leçon, je sais utiliser le vocabulaire de base lié à la délinquance juvénile (犯罪, 青少年, 法律, 教育...)
\item je sais présenter oralement un exposé simple et structuré en chinois sur ce thème
\item je sais donner mon avis avec les formules 我认为…, 我觉得…, 在我看来…
\item je sais exprimer l'accord et le désaccord poliment (我同意/我不同意, 你说得对…)
\item je sais relancer un débat en posant des questions (你觉得呢？为什么呢？)
\item je sais prononcer correctement les tons des mots clés de la leçon
\item je sais participer à un jeu de rôle et à un débat guidé en classe
\item je sais comparer brièvement la situation au Cameroun et en Chine
\end{itemize}
\end{objectifs}

\begin{prerequis}
\begin{itemize}
\item Se présenter et parler de sa famille (Module 1, leçons précédentes)
\item Les quatre tons et le pinyin (acquis de Seconde)
\item Les structures de phrase simples : sujet + verbe + complément
\item Les pronoms personnels et les verbes modaux 应该, 可以, 会
\item Le vocabulaire de l'école et de la société vu en début d'année
\end{itemize}
\end{prerequis}

\newpage

\coursec{Vocabulaire du thème : la délinquance juvénile}

\courssub{Situation d'accroche et problématique}

À Yaoundé comme à Douala, les journaux parlent souvent des \emph{microbes}, ces bandes de jeunes qui sèment l'insécurité dans certains quartiers. Au Royaume-Uni, on débat du \emph{youth crime} ; en Chine, la loi sur la prévention de la délinquance juvénile (\zh{预防未成年人犯罪法}) encadre strictement les mineurs. Comment parler de ce phénomène en chinois ? Comment donner son avis, réagir, relancer un débat ? Cette leçon vous donne les mots et les formules pour en discuter à l'oral.

\courssub{Texte d'observation : la délinquance juvénile, un problème social}

\begin{textebox}[Texte d'observation]
青少年犯罪是一个严重的社会问题。\\
Qīngshàonián fànzuì shì yí ge yánzhòng de shèhuì wèntí.\\
La délinquance juvénile est un grave problème social.\\
\\
很多青少年因为家庭教育缺失，走上了犯罪的道路。\\
Hěnduō qīngshàonián yīnwèi jiātíng jiàoyù quēshī, zǒushàng le fànzuì de dàolù.\\
Beaucoup d'adolescents, à cause d'un manque d'éducation familiale, s'engagent sur la voie de la délinquance.\\
\\
他们有的偷东西，有的抢劫，还有的打架斗殴。\\
Tāmen yǒu de tōu dōngxi, yǒu de qiǎngjié, hái yǒu de dǎjià dòu'ōu.\\
Certains volent, d'autres commettent des vols avec violence, d'autres encore se battent.\\
\\
警察和法院依法惩罚罪犯，但更重要的是预防和教育。\\
Jǐngchá hé fǎyuàn yīfǎ chéngfá zuìfàn, dàn zhòngyào de shì yùfáng hé jiàoyù.\\
La police et les tribunaux sanctionnent les criminels selon la loi, mais la prévention et l'éducation sont plus importantes.
\end{textebox}

\begin{center}
\begin{tabularx}{\linewidth}{|X|X|X|X|}
\hline
\rowcolor{popblueL} Caractères & Pinyin & Nature & Traduction \\
\hline
青少年犯罪 & qīngshàonián fànzuì & NM & délinquance juvénile \\
严重 & yánzhòng & ADJ & grave, sérieux \\
社会问题 & shèhuì wèntí & NM & problème social \\
家庭教育 & jiātíng jiàoyù & NM & éducation familiale \\
缺失 & quēshī & V & manquer, faire défaut \\
走上 & zǒushàng & V & s'engager dans, emprunter \\
道路 & dàolù & N & route, voie, chemin \\
偷东西 & tōu dōngxi & VO & voler (des choses) \\
抢劫 & qiǎngjié & V/N & voler avec violence, braquer / braquage \\
打架斗殴 & dǎjià dòu'ōu & VO & se battre, bagarre \\
警察 & jǐngchá & N & police, policier \\
法院 & fǎyuàn & N & tribunal \\
依法 & yīfǎ & ADV & conformément à la loi \\
惩罚 & chéngfá & V & punir, sanctionner \\
罪犯 & zuìfàn & N & criminel, délinquant \\
预防 & yùfáng & V & prévenir \\
\hline
\end{tabularx}
\end{center}

\courssub{Mots-clés du thème : analyse morphologique et emplois}

\begin{definition}[青少年犯罪  ---  mot central de la leçon]
\zh{青少年犯罪} (qīngshàonián fànzuì) est un nom composé de trois morphèmes :
\begin{itemize}
\item \zh{青少年} (qīngshàonián) = \zh{青} (jeune, vert) + \zh{少年} (adolescent) $\rightarrow$ « adolescents, jeunes » (terme formel, administratif, médiatique).
\item \zh{犯} (fàn) = « commettre (une faute, un crime), violer ».
\item \zh{罪} (zuì) = « crime, faute, péché ».
\end{itemize}
\zh{犯罪} (fànzuì) fonctionne comme un verbe intransitif (« commettre un crime ») ou un nom (« la criminalité, le crime »).
\emph{Attention :} \zh{罪犯} (zuìfàn, N) = « le criminel, le délinquant » (ordre inversé).
\end{definition}

\begin{propriete}[Structure de \zh{青少年} et dérivés]
\begin{center}
\begin{tabularx}{\linewidth}{|X|X|X|}
\hline
\rowcolor{popblueL} Terme & Pinyin & Emploi / Nuance \\
\hline
青少年 & qīngshàonián & Nom collectif, formel : « les adolescents (12--18 ans) » \\
少年 & shàonián & « jeune, adolescent » (littéraire, ou dans \zh{少年宫} « palais de la jeunesse ») \\
青年 & qīngnián & « jeune adulte, jeunesse » (18--35 ans, ex. \zh{青年节} Fête de la Jeunesse) \\
未成年人 & wèichéngniánrén & Juridique : « mineur (moins de 18 ans) » \\
\hline
\end{tabularx}
\end{center}
\end{propriete}

\begin{exemplebox}[Emplois de \zh{犯罪} / \zh{罪犯}]
\begin{itemize}
\item \zh{他犯了罪。} Tā fàn le zuì. --- Il a commis un crime.
\item \zh{青少年犯罪率上升了。} Qīngshàonián fànzuìlǜ shàngshēng le. --- Le taux de délinquance juvénile a augmenté.
\item \zh{警察抓住了罪犯。} Jǐngchá zhuāzhù le zuìfàn. --- La police a attrapé le criminel.
\item \zh{这是一起严重的犯罪案件。} Zhè shì yī qǐ yánzhòng de fànzuì ànjiàn. --- C'est une affaire criminelle grave.
\end{itemize}
\end{exemplebox}

\courssub{Les actes délinquants : distinguer \zh{偷}, \zh{抢}, \zh{打架}}

\begin{attention}[Erreur fréquente : confondre \zh{偷} et \zh{抢}]
En français, « voler » couvre les deux sens. En chinois, le choix du verbe dépend de la \textbf{modalité} :
\begin{itemize}
\item \zh{偷} (tōu, V) = « voler \emph{discrètement}, à l'insu de la victime » (pickpocket, vol à la tire, cambriolage sans effraction violente). Objet souvent : \zh{东西} (dōngxi), \zh{钱} (qián), \zh{手机} (shǒujī).
\item \zh{抢} (qiǎng, V) = « arracher, voler \emph{avec violence ou menace}, en présence de la victime » (braquage, vol à l'arraché). Souvent composé : \zh{抢劫} (qiǎngjié, V/N « vol à main armée, braquage »), \zh{抢夺} (qiǎngduó « arracher »).
\item \zh{打架} (dǎjià, VO) = « se battre, se bagarrer » (souvent entre jeunes, sans arme à feu). \zh{斗殴} (dòu'ōu) est plus formel/juridique (« rixe, bagarre collective »).
\end{itemize}
\emph{Astuce mnémotechnique :} \zh{偷} a le radical \zh{人} (homme qui se cache) ; \zh{抢} a le radical \zh{手} (main qui arrache).
\end{attention}

\begin{center}
\begin{tabularx}{\linewidth}{|X|X|X|X|}
\hline
\rowcolor{popblueL} Verbe & Pinyin & Sens précis & Exemple typique \\
\hline
偷 & tōu & voler discrètement & \zh{小偷偷了我的钱包。} (Le pickpocket a volé mon portefeuille.) \\
抢 & qiǎng & arracher / voler avec violence & \zh{歹徒抢了她的手机。} (Le bandit lui a arraché son téléphone.) \\
抢劫 & qiǎngjié & braquer / vol à main armée & \zh{银行被抢劫了。} (La banque a été braquée.) \\
打架 & dǎjià & se battre (entre individus) & \zh{他们在学校门口打架。} (Ils se sont battus devant l'école.) \\
斗殴 & dòu'ōu & rixe, bagarre collective & \zh{群架斗殴事件。} (Affaire de rixe collective.) \\
\hline
\end{tabularx}
\end{center}

\courssub{Acteurs, institutions et solutions : vocabulaire institutionnel}

\begin{definition}[Acteurs clés : \zh{警察}, \zh{家长}, \zh{老师}, \zh{法院}]
\begin{itemize}
\item \zh{警察} (jǐngchá) : « police, policier » (terme générique). \zh{民警} (mínjǐng) = « policier de métier » (plus formel).
\item \zh{家长} (jiāzhǎng) : « parent d'élève, tuteur légal » (contexte scolaire/éducatif). \zh{父母} (fùmǔ) = « parents » (familial, affectif).
\item \zh{老师} (lǎoshī) : « professeur, maître » (terme de respect, s'applique aussi aux éducateurs).
\item \zh{法院} (fǎyuàn) : « tribunal, cour ». \zh{法官} (fǎguān) = « juge ». \zh{检察院} (jiǎncháyuàn) = « parquet, procurature ».
\item \zh{监狱} (jiānyù) : « prison ». \emph{Prononciation :} \zh{jiān} (1er ton) + \zh{yù} (4e ton). Ne pas confondre avec \zh{简} (jiǎn, 3e ton).
\end{itemize}
\end{definition}

\begin{propriete}[Solutions : \zh{教育}, \zh{法律}, \zh{预防}, \zh{解决}]
\begin{center}
\begin{tabularx}{\linewidth}{|X|X|X|}
\hline
\rowcolor{popblueL} Mot & Pinyin & Collocations fréquentes \\
\hline
教育 & jiàoyù & \zh{家庭教育} (éducation familiale), \zh{法治教育} (éducation à l'État de droit), \zh{接受教育} (recevoir une éducation) \\
法律 & fǎlǜ & \zh{遵守法律} (respecter la loi), \zh{违法} (wéifǎ, enfreindre la loi), \zh{法律意识} (conscience juridique) \\
预防 & yùfáng & \zh{预防为主} (la prévention avant tout), \zh{预防犯罪} (prévenir la criminalité) \\
解决 & jiějué & \zh{解决问题} (résoudre le problème), \zh{解决办法} (solution, moyen de résoudre) \\
矫正 & jiǎozhèng & \zh{矫正措施} (mesures de rééducation), \zh{劳动教养} (rééducation par le travail --- ancien terme) \\
\hline
\end{tabularx}
\end{center}
\end{propriete}

\courssub{Phonétique ciblée : tons à travailler}

\begin{aretenir}[Paires tonales et pièges de prononciation]
\begin{itemize}
\item \zh{法律} (fǎlǜ) : 3e ton + 4e ton. Le \zh{ǜ} (ü, 4e ton) s'arrondit fortement ; ne pas dire *fǎlù.
\item \zh{监狱} (jiānyù) : 1er ton + 4e ton. \zh{jiān} (haut, plat) $\neq$ \zh{jiǎn} (3e ton, « simple »). \emph{Erreur fréquente :} *jiǎnyù.
\item \zh{原因} (yuányīn) : 2e ton (montant) + 1er ton (haut). \zh{yuán} monte comme une question ; \zh{yīn} reste haut.
\item \zh{青少年} (qīngshàonián) : 1 + 4 + 2. \zh{shào} (4e ton, chute nette) ; \zh{nián} (2e ton, montant).
\item \zh{犯罪} (fànzuì) : 4 + 4. Deux chutes nettes ; ne pas « traîner » le 4e ton.
\item \zh{偷} (tōu) 1er ton vs \zh{透} (tòu) 4e ton (traverser) --- homophones segmentaux, tons distinctifs.
\item \zh{抢} (qiǎng) 3e ton vs \zh{强} (qiáng) 2e ton (fort) / (qiǎng) 3e ton (forcer, faire contre son gré).
\end{itemize}
\end{aretenir}

\begin{experience}[Entraînement phonétique : « Tons en miroir »]
En binômes. L'élève A lit une phrase ; l'élève B la répète en inversant l'ordre des mots-clés (ex. A : \zh{警察抓小偷} $\rightarrow$ B : \zh{小偷被警察抓}). Focus sur \zh{fǎlǜ, jiānyù, yuányīn, qīngshàonián}. L'enseignant circule et corrige les tons 3+4 et 2+1.
\end{experience}

\courssub{Carte mentale lexicale : \zh{青少年犯罪}}

\begin{popfigure}
\begin{tikzpicture}[
  mindmap,
  grow cyclic,
  every node/.style={concept, execute at begin node=\hskip0pt, align=center},
  root concept/.append style={concept color=popblue, minimum size=2.8cm, text=white, font=\bfseries\large},
  level 1 concept/.append style={sibling angle=90, level distance=4.2cm, concept color=poporange, minimum size=2.2cm, text=white, font=\bfseries\normalsize},
  level 2 concept/.append style={sibling angle=45, level distance=3.2cm, concept color=popgreen, minimum size=1.8cm, text=white, font=\small},
  level 3 concept/.append style={sibling angle=30, level distance=2.5cm, concept color=poppurple, minimum size=1.5cm, text=white, font=\footnotesize},
  edge from parent/.style={draw=popline, thick}
]
\node[align=center]{青少年犯罪\\qīngshàonián fànzuì}
  child [concept color=poporange] {node[align=center]{Causes\\原因}
    child {node[align=center]{家庭\\jiātíng}
      child {node[align=center]{缺失教育\\quēshī jiàoyù}}
      child {node[align=center]{离异\\líyì}}
    }
    child {node[align=center]{学校\\xuéxiào}
      child {node[align=center]{辍学\\chuòxué}}
      child {node[align=center]{霸凌\\bàlíng}}
    }
    child {node[align=center]{社会\\shèhuì}
      child {node[align=center]{不良风气\\bùliáng fēngqì}}
      child {node[align=center]{网络诱惑\\wǎngluò yòuhuò}}
    }
  }
  child [concept color=poppink] {node[align=center]{Actes\\行为}
    child {node[align=center]{偷\\tōu}
      child {node[align=center]{扒手\\bāshǒu}}
    }
    child {node[align=center]{抢\\qiǎng}
      child {node[align=center]{抢劫\\qiǎngjié}}
    }
    child {node[align=center]{打架\\dǎjià}
      child {node[align=center]{斗殴\\dòu'ōu}}
    }
  }
  child [concept color=popgold] {node[align=center]{Acteurs\\角色}
    child {node[align=center]{警察\\jǐngchá}
      child {node[align=center]{民警\\mínjǐng}}
    }
    child {node[align=center]{家长\\jiāzhǎng}
      child {node[align=center]{监护人\\jiānhùrén}}
    }
    child {node[align=center]{老师\\lǎoshī}
      child {node[align=center]{班主任\\bānzhǔrèn}}
    }
    child {node[align=center]{法院\\fǎyuàn}
      child {node[align=center]{法官\\fǎguān}}
    }
  }
  child [concept color=popgreen] {node[align=center]{Solutions\\对策}
    child {node[align=center]{教育\\jiàoyù}
      child {node[align=center]{法治教育\\fǎzhì jiàoyù}}
    }
    child {node[align=center]{法律\\fǎlǜ}
      child {node[align=center]{预防法\\yùfáng fǎ}}
    }
    child {node[align=center]{预防\\yùfáng}
      child {node[align=center]{社区矫正\\shèqū jiǎozhèng}}
    }
    child {node[align=center]{解决\\jiějué}
      child {node[align=center]{综合治理\\zōnghé zhìlǐ}}
    }
  };
\end{tikzpicture}
\legende{Carte mentale du lexique « délinquance juvénile » : causes, actes, acteurs, solutions.}
\end{popfigure}

\courssub{Exemples traduits : mise en contexte camerounaise et chinoise}

\begin{exemplebox}[Phrases modèles pour l'oral]
\begin{enumerate}
\item \zh{青少年犯罪在喀麦隆和中国都是一个热门话题。}\\
Qīngshàonián fànzuì zài Kāmàilóng hé Zhōngguó dōu shì yí ge rèmén huàtí.\\
La délinquance juvénile est un sujet brûlant au Cameroun comme en Chine.
\item \zh{很多「微生物」青少年在杜阿拉抢劫行人。}\\
Hěnduō ``Wēishēngwù'' qīngshàonián zài Dù'ālā qiǎngjié xíngrén.\\
Beaucoup de jeunes « microbes » braquent des passants à Douala.
\item \zh{家长应该加强家庭教育，预防孩子犯罪。}\\
Jiāzhǎng yīnggāi jiāqiáng jiātíng jiàoyù, yùfáng háizi fànzuì.\\
Les parents devraient renforcer l'éducation familiale pour prévenir la délinquance de leurs enfants.
\item \zh{警察依法惩罚罪犯，同时也要保护未成年人权利。}\\
Jǐngchá yīfǎ chéngfá zuìfàn, tóngshí yě yào bǎohù wèichéngniánrén quánlì.\\
La police sanctionne les délinquants selon la loi, tout en protégeant les droits des mineurs.
\item \zh{我们应该用教育代替单纯的惩罚，帮助他们改过自新。}\\
Wǒmen yīnggāi yòng jiàoyù dàitì dānchún de chéngfá, bāngzhù tāmen gǎoguòzìxīn.\\
Nous devrions remplacer la simple punition par l'éducation, pour les aider à se réinsérer.
\end{enumerate}
\end{exemplebox}

\courssub{Le saviez-vous ? : Cadre légal chinois}

\begin{savaistu}[La loi chinoise de prévention de la délinquance juvénile]
La \zh{预防未成年人犯罪法} (Yùfáng Wèichéngniánrén Fànzuì Fǎ), promulguée en 1999 et révisée en 2012 et 2020, pose le principe \zh{教育、感化、挽救} (éduquer, réformer, sauver). Elle interdit aux mineurs :
\begin{itemize}
\item l'accès aux cybercafés (\zh{网吧}) la nuit (22h--8h) ;
\item l'entrée dans les salles de jeux d'argent, bars, karaokés ;
\item la vente d'alcool et de tabac aux moins de 18 ans.
\end{itemize}
Depuis 2021, la responsabilité pénale peut être engagée dès \textbf{12 ans} pour homicide ou blessure grave (auparavant 14 ans). Des \zh{专门学校} (écoles spécialisées) et \zh{未成年人保护中心} (centres de protection des mineurs) assurent la rééducation plutôt que l'incarcération classique.
\end{savaistu}

\courssub{Exercice résolu : choisir le bon verbe et le bon mot}

\begin{exoresolu}[Complétez avec \zh{偷 / 抢 / 抢劫 / 打架 / 斗殴 / 犯罪 / 罪犯 / 警察 / 法律 / 预防}]
\textbf{Énoncé :}
1. 昨晚，两个年轻人 \_\_\_\_\_ 了一位老人的手提包，老人摔倒了。\\
2. 学校门口发生了一起 \_\_\_\_\_，有五个学生受伤。\\
3. \_\_\_\_\_ 及时赶到，制服了持刀歹徒。\\
4. 我们要加强法治 \_\_\_\_\_，从源头减少青少年 \_\_\_\_\_。\\
5. 那个 \_\_\_\_\_ 已经被法院判刑三年。\\

\textbf{Solution détaillée :}
1. \zh{抢} (qiǎng) --- vol avec violence/arraché sur une personne âgée (présence de la victime, violence). \emph{Phrase :} 昨晚，两个年轻人 \textbf{抢} 了一位老人的手提包，老人摔倒了。 (Zuówǎn, liǎng ge niánqīngrén \textbf{qiǎng} le yí wèi lǎorén de shǒutíbāo, lǎorén shuāidǎo le.)
2. \zh{斗殴} (dòu'ōu) --- bagarre collective, plusieurs blessés, contexte scolaire. \emph{Phrase :} 学校门口发生了一起 \textbf{斗殴}，有五个学生受伤。 (Xuéxiào ménkǒu fāshēng le yī qǐ \textbf{dòu'ōu}, yǒu wǔ ge xuéshēng shòushāng.)
3. \zh{警察} (jǐngchá) --- sujet de l'action « arriver, maîtriser ». \emph{Phrase :} \textbf{警察} 及时赶到，制服了持刀歹徒。 (\textbf{Jǐngchá} jíshí gǎndào, zhìfú le chídāo dǎitú.)
4. \zh{预防} (yùfáng) + \zh{犯罪} (fànzuì) --- collocations standards : « prévention » (nom/verbe) + « criminalité/délinquance ». \emph{Phrase :} 我们要加强法治 \textbf{预防}，从源头减少青少年 \textbf{犯罪}。 (Wǒmen yào jiāqiáng fǎzhì \textbf{yùfáng}, cóng yuántóu jiǎnshǎo qīngshàonián \textbf{fànzuì}.)
5. \zh{罪犯} (zuìfàn) --- « criminel » (nom d'agent), passé devant le tribunal. \emph{Phrase :} 那个 \textbf{罪犯} 已经被法院判刑三年。 (Nà ge \textbf{zuìfàn} yǐjīng bèi fǎyuàn pànxíng sān nián.)
\end{exoresolu}

\courssub{Exercices d'appropriation}

\begin{exercice}[Entraînement 1 : Association mot -- définition]
Reliez chaque terme chinois à sa définition française.
\begin{enumerate}
\item \zh{青少年犯罪} \hfill A. Vol avec violence, menace
\item \zh{偷} \hfill B. Mineur, personne de moins de 18 ans
\item \zh{抢劫} \hfill C. Délinquance juvénile
\item \zh{未成年人} \hfill D. Vol discret, à l'insu de la victime
\item \zh{预防} \hfill E. Prévenir, agir en amont
\item \zh{罪犯} \hfill F. Criminel, délinquant condamné
\item \zh{法治教育} \hfill G. Éducation à l'État de droit
\item \zh{监狱} \hfill H. Prison
\end{enumerate}
\end{exercice}

\begin{exercice}[Entraînement 2 : Production de phrases guidées]
À l'aide des mots donnés, écrivez une phrase chinoise (caractères + pinyin) et sa traduction française.
\begin{enumerate}
\item \zh{青少年 / 犯罪 / 原因 / 复杂} (les causes de la délinquance juvénile sont complexes)
\item \zh{家长 / 应该 / 加强 / 家庭教育} (les parents devraient renforcer l'éducation familiale)
\item \zh{警察 / 依法 / 惩罚 / 罪犯} (la police sanctionne les criminels selon la loi)
\item \zh{我们 / 预防 / 为主 / 解决 / 问题} (nous privilégions la prévention pour résoudre le problème)
\item \zh{喀麦隆 / 也 / 有 / 青少年 / 打架 / 现象} (le Cameroun connaît aussi des phénomènes de bagarres entre jeunes)
\end{enumerate}
\end{exercice}

\begin{corrige}[Exercices 1 et 2]
\textbf{Exercice 1 :} 1--C, 2--D, 3--A, 4--B, 5--E, 6--F, 7--G, 8--H.

\textbf{Exercice 2 (propositions) :}
1. \zh{青少年犯罪的原因很复杂。} Qīngshàonián fànzuì de yuányīn hěn fùzá.
2. \zh{家长应该加强家庭教育。} Jiāzhǎng yīnggāi jiāqiáng jiātíng jiàoyù.
3. \zh{警察依法惩罚罪犯。} Jǐngchá yīfǎ chéngfá zuìfàn.
4. \zh{我们以预防为主，解决问题。} Wǒmen yǐ yùfáng wéi zhǔ, jiějué wèntí.
5. \zh{喀麦隆也有青少年打架的现象。} Kāmàilóng yě yǒu qīngshàonián dǎjià de xiànxiàng.
\end{corrige}

\courssub{Tableau récapitulatif du vocabulaire de la section}

\begin{center}
\begin{tabularx}{\linewidth}{|X|X|X|X|}
\hline
\rowcolor{popblueL} Caractères & Pinyin & Nature & Traduction \\
\hline
青少年犯罪 & qīngshàonián fànzuì & NM & délinquance juvénile \\
青少年 & qīngshàonián & N & adolescents, jeunes \\
少年 & shàonián & N & jeune, adolescent \\
未成年人 & wèichéngniánrén & N & mineur (< 18 ans) \\
犯罪 & fànzuì & V/N & commettre un crime / criminalité \\
罪犯 & zuìfàn & N & criminel, délinquant \\
偷 & tōu & V & voler (discrètement) \\
抢 & qiǎng & V & arracher, voler (avec violence) \\
抢劫 & qiǎngjié & V/N & braquer / braquage \\
打架 & dǎjià & VO & se battre \\
斗殴 & dòu'ōu & V/N & rixe, bagarre collective \\
法律 & fǎlǜ & N & loi \\
警察 & jǐngchá & N & police, policier \\
法院 & fǎyuàn & N & tribunal \\
监狱 & jiānyù & N & prison \\
家庭教育 & jiātíng jiàoyù & NM & éducation familiale \\
预防 & yùfáng & V & prévenir \\
解决 & jiějué & V & résoudre \\
原因 & yuányīn & N & cause \\
影响 & yǐngxiǎng & V/N & influencer / influence \\
家长 & jiāzhǎng & N & parent d'élève, tuteur \\
老师 & lǎoshī & N & professeur, éducateur \\
\hline
\end{tabularx}
\end{center}

\begin{methode}[Comment réviser ce vocabulaire pour l'oral ?]
\begin{enumerate}
\item \textbf{Écoute active :} relisez le texte de la \texttt{textebox} en suivant le pinyin ; répétez chaque phrase 3 fois en imitant les tons (enregistrez-vous).
\item \textbf{Cartes mentales :} reproduisez la carte mentale TikZ sur papier ; ajoutez 2 mots par branche (ex. \zh{网络暴力} sous « Causes / Société »).
\item \textbf{Paires minimales :} entraînez-vous sur \zh{偷/透}, \zh{抢/强}, \zh{监/简}, \zh{法/发} pour verrouiller les tons.
\item \textbf{Micro-débat :} en binôme, l'un défend « punir », l'autre « éduquer » en utilisant au moins 5 mots du tableau.
\item \textbf{Contexte Cameroun :} préparez 3 phrases sur les « microbes » de Yaoundé/Douala avec \zh{青少年犯罪, 抢劫, 家庭教育, 预防, 法律}.
\end{enumerate}
\end{methode}

\coursec{Formules fonctionnelles : présenter, opiner, relancer}

\courssub{Structures pour un exposé oral structuré}

\begin{propriete}[Organiser son propos : connecteurs logiques]
\begin{center}
\begin{tabularx}{\linewidth}{|X|X|X|}
\hline
\rowcolor{popblueL} Fonction & Structure chinoise (Pinyin) & Équivalent français \\
\hline
Introduire le sujet & \zh{今天我想谈谈……} (Jīntiān wǒ xiǎng tántan \dots) & Aujourd'hui, je veux parler de \dots \\
& \zh{关于……这个问题} (Guānyú \dots zhège wèntí) & Concernant le problème de \dots \\
Énumérer (1) & \zh{首先……} (Shǒuxiān \dots) & Premièrement \dots / Tout d'abord \dots \\
Énumérer (2) & \zh{其次……} (Qícì \dots) & Deuxièmement \dots / Ensuite \dots \\
Énumérer (3) & \zh{最后……} (Zuìhòu \dots) & Enfin \dots / Pour finir \dots \\
Ajouter un argument & \zh{而且……} (Érqiě \dots) / \zh{此外……} (Cǐwài \dots) & De plus \dots / En outre \dots \\
Conclure & \zh{总之……} (Zǒngzhī \dots) / \zh{综上所述……} (Zōngshàng shùshù \dots) & En résumé \dots / En conclusion \dots \\
\hline
\end{tabularx}
\end{center}
\end{propriete}

\courssub{Exprimer son opinion et nuancer}

\begin{propriete}[Donner son avis : échelles de certitude]
\begin{center}
\begin{tabularx}{\linewidth}{|X|X|X|}
\hline
\rowcolor{popblueL} Niveau & Structure & Traduction \\
\hline
Opinion personnelle & \zh{我认为……} (Wǒ rènwéi \dots) & Je pense que \dots (standard) \\
& \zh{依我看来……} (Yī wǒ kànlái \dots) & À mon avis \dots / Selon moi \dots \\
Accord & \zh{我同意……} (Wǒ tóngyì \dots) & Je suis d'accord \dots \\
& \zh{有道理。} (Yǒu dàoli.) & C'est vrai / Il y a du bon sens. \\
Désaccord poli & \zh{我不太同意，因为……} (Wǒ bù tài tóngyì, yīnwèi \dots) & Je ne suis pas tout à fait d'accord, car \dots \\
& \zh{恐怕不太对吧？} (Kǒngpà bù tài duì ba?) & Je crains que ce ne soit pas tout à fait exact ? \\
Nuance / Concession & \zh{虽然……，但是……} (Suīrán \dots, dànshì \dots) & Bien que \dots, cependant \dots \\
& \zh{一方面……，另一方面……} (Yī fāngmiàn \dots, lìng yī fāngmiàn \dots) & D'un côté \dots, de l'autre \dots \\
\hline
\end{tabularx}
\end{center}
\end{propriete}

\courssub{Relancer le débat et interagir}

\begin{definition}[Formules de relance et de gestion de tour de parole]
\begin{itemize}
\item \zh{你怎么看？} (Nǐ zěnme kàn?) --- « Qu'en penses-tu ? / Quel est ton avis ? »
\item \zh{你同意吗？} (Nǐ tóngyì ma?) --- « Es-tu d'accord ? »
\item \zh{还有什么要补充的吗？} (Hái yǒu shénme yào bǔchōng de ma?) --- « Y a-t-il quelque chose à ajouter ? »
\item \zh{请说说你的理由。} (Qǐng shuōshuō nǐ de lǐyóu.) --- « Explique ta raison / Justifie-toi. »
\item \zh{比如说呢？} (Bǐrú shuō ne?) --- « Par exemple ? »
\item \zh{换个角度想……} (Huàn gè jiǎodu xiǎng \dots) --- « Si on change d'angle / Si on voit ça autrement \dots »
\item \zh{我们来总结一下。} (Wǒmen lái zǒngjié yíxià.) --- « Faisons un résumé. »
\end{itemize}
\end{definition}

\begin{exemplebox}[Mini-échanges types]
\begin{enumerate}
\item A: \zh{我认为预防比惩罚重要。} (Je pense que la prévention est plus importante que la punition.)\\
B: \zh{有道理。但是法律也必须严厉。} (C'est vrai. Mais la loi doit aussi être sévère.)
\item A: \zh{青少年犯罪的原因很复杂。} (Les causes de la délinquance juvénile sont complexes.)\\
B: \zh{比如说呢？} (Par exemple ?)\\
A: \zh{比如家庭教育缺失、网络诱惑等。} (Par exemple le manque d'éducation familiale, les pièges d'Internet, etc.)
\end{enumerate}
\end{exemplebox}

\coursec{Dialogue modèle 2 : Le débat — Punir ou éduquer ?}

\courssub{Contexte : Deux lycéens, \zh{小李} (Xiǎo Lǐ) et \zh{小王} (Xiǎo Wáng), discutent après avoir lu un article sur un vol à l'arraché commis par un mineur.}

\begin{textebox}[Dialogue : 惩罚还是教育？]
小李：你看这则新闻，一个未成年人抢劫了老人，太可恶了！\\
Xiǎo Lǐ: Nǐ kàn zhè zé xīnwén, yí ge wèichéngniánrén qiǎngjié le lǎorén, tài kěwù le!\\
Petit Li : Regarde cette info, un mineur a braqué une personne âgée, c'est odieux !

小王：确实可恶。但是，你觉得单纯判刑能解决问题吗？\\
Xiǎo Wáng: Quèshí kěwù. Dànshì, nǐ juéde dānchún pànxíng néng jiějué wèntí ma?\\
Petit Wang : C'est odieux en effet. Mais penses-tu qu'une simple peine de prison résout le problème ?

小李：法律必须严厉，否则没人怕。惩罚是为了威慑。\\
Xiǎo Lǐ: Fǎlǜ bìxū yánlì, fǒuzé méi rén pà. Chéngfá shì wèile wēishè.\\
Petit Li : La loi doit être sévère, sinon personne n'a peur. La punition sert à dissuader.

小王：我不同意。很多青少年犯罪是因为家庭教育缺失。\\
Xiǎo Wáng: Wǒ bù tóngyì. Hěnduō qīngshàonián fànzuì shì yīnwèi jiātíng jiàoyù quēshī.\\
Petit Wang : Je ne suis pas d'accord. Beaucoup de délinquances juvéniles sont dues au manque d'éducation familiale.

小李：那你有什么办法？\\
Xiǎo Lǐ: Nà nǐ yǒu shénme bànfǎ?\\
Petit Li : Alors quelle est ta solution ?

小王：应该以预防为主，加强法治教育，帮助他们改过自新。\\
Xiǎo Wáng: Yīnggāi yǐ yùfáng wéi zhǔ, jiāqiáng fǎzhì jiàoyù, bāngzhù tāmen gǎiguòzìxīn.\\
Petit Wang : Il faut privilégier la prévention, renforcer l'éducation à l'État de droit, les aider à se réinsérer.

小李：有道理。预防和惩罚要结合。\\
Xiǎo Lǐ: Yǒu dàoli. Yùfáng hé chéngfá yào jiéhé.\\
Petit Li : C'est juste. Prévention et punition doivent être combinées.
\end{textebox}

\courssub{Analyse linguistique du dialogue}

\begin{propriete}[Points de grammaire clés observés]
\begin{itemize}
\item \zh{单纯} (dānchún) : « simple, pur, uniquement » (adj). \zh{单纯判刑} = « seulement prononcer une peine ».
\item \zh{威慑} (wēishè) : « dissuasion, intimidation » (V/N). \zh{起威慑作用} (qǐ wēishè zuòyòng) = « jouer un rôle dissuasif ».
\item \zh{否则} (fǒuzé) : « sinon, autrement » (conjonction, registre formel/écrit). À l'oral : \zh{要不} (yàobù).
\item \zh{改过自新} (gǎiguòzìxīn) : chengyu (成语) « corriger ses fautes et recommencer à neuf » = « se réinsérer, se repentir sincèrement ».
\item \zh{结合} (jiéhé) : « combiner, associer ». Structure : A \zh{和/与} B \zh{结合}.
\end{itemize}
\end{propriete}

\courssub{Phonétique du dialogue : focus sur l'intonation interrogative et l'insistance}

\begin{aretenir}[Intonation et accent de phrase]
\begin{itemize}
\item Questions ouvertes (\zh{怎么}, \zh{什么}) : ton descendant à la fin (↘). Ex: \zh{你怎么看？} (Nǐ zěnme kàn↘)
\item Questions fermées (\zh{吗}) : ton montant à la fin (↗). Ex: \zh{你同意吗？} (Nǐ tóngyì ma↗)
\item Insistance sur \zh{必须} (bìxū), \zh{应该} (yīnggāi) : accentuer le premier caractère, ton plein.
\item \zh{可恶} (kěwù) : 3e ton + 4e ton. \zh{kě} fait le "creux" (214) avant \zh{wù} (51).
\end{itemize}
\end{aretenir}

\begin{experience}[Lecture expressive du dialogue]
En groupes de 3 (Li, Wang, Narrateur). Travailler l'attitude : Li = indigné/ferme ; Wang = calme/réfléchi ; Narrateur = neutre. Enregistrer et comparer avec le modèle audio de l'application OnBuch+.
\end{experience}

\coursec{Jeux de rôle et débat guidé}

\courssub{Activité 1 : Jeu de rôle « Au commissariat » (3--4 personnes)}

\begin{experience}[Scénario : Interrogatoire / Médiation]
\textbf{Rôles :}
\begin{itemize}
\item \textbf{Le policier} (\zh{警察}) : ferme, procédural, utilise \zh{依法}, \zh{供认}, \zh{配合调查}.
\item \textbf{Le jeune délinquant} (\zh{未成年人}) : défensif ou repentant, utilise \zh{我没想到}, \zh{一时冲动}, \zh{后悔}.
\item \textbf{Le parent / tuteur} (\zh{家长 / 监护人}) : inquiet, responsable, utilise \zh{加强管教}, \zh{配合学校}, \zh{担保}.
\item \textbf{L'éducateur / Juge pour enfants} (optionnel, \zh{法官 / 社工}) : orienté solution, utilise \zh{矫正}, \zh{心理疏导}, \zh{社区服务}.
\end{itemize}
\textbf{Consignes :} 5 min de préparation (notes en pinyin/mots-clés autorisés), 3--4 min de jeu. Le policier doit obtenir les faits ; le parent doit protéger l'intérêt de l'enfant ; l'éducateur propose une mesure \zh{非监禁化} (non carcérale).
\textbf{Lexique utile :} \zh{供认} (gòngrèn, avouer), \zh{一时冲动} (yíshí chōngdòng, coup de tête), \zh{悔过} (huǐguò, regretter), \zh{取保候审} (qǔbǎo hòushěn, liberté sous caution), \zh{社区矫正} (shèqū jiǎozhèng, rééducation communautaire).
\end{experience}

\courssub{Activité 2 : Débat guidé « La responsabilité pénale à 12 ans : pour ou contre ? »}

\begin{methode}[Déroulement du débat structuré (20 min)]
\begin{enumerate}
\item \textbf{Préparation (5 min) :} En deux groupes (Pour / Contre). Chaque groupe note 3 arguments en chinois (mots-clés + structures \zh{首先/其次/而且}).
\item \textbf{Arguments « Pour » (ex.) :}
      \begin{itemize}
      \item \zh{严重犯罪必须严惩，维护法律尊严。} (Les crimes graves doivent être sévèrement punis pour défendre la dignité de la loi.)
      \item \zh{威慑作用，预防恶性案件低龄化。} (Effet dissuasif, prévenir la baisse d'âge des affaires graves.)
      \item \zh{受害者权利得到保障。} (Les droits des victimes sont garantis.)
      \end{itemize}
\item \textbf{Arguments « Contre » (ex.) :}
      \begin{itemize}
      \item \zh{未成年人心智不成熟，应以教育矫正为主。} (Les mineurs sont immatures, il faut privilégier l'éducation/rééducation.)
      \item \zh{监狱环境不利于改过自新，易再犯。} (L'environnement carcéral nuit à la réinsertion, risque de récidive.)
      \item \zh{家庭、学校、社会责任更大，不能只惩罚孩子。} (Famille, école, société portent une plus grande responsabilité, on ne peut pas seulement punir l'enfant.)
      \end{itemize}
\item \textbf{Débat (10 min) :} Tour de parole (1 min/groupe). Utiliser formules de relance : \zh{你怎么看？}, \zh{我不太同意，因为……}, \zh{有道理，但是……}.
\item \textbf{Synthèse (5 min) :} Un élève par groupe fait une conclusion orale de 30 sec avec \zh{总之}, \zh{建议}.
\end{enumerate}
\end{methode}

\courssub{Grille d'auto-évaluation orale (à remplir après l'activité)}

\begin{center}
\begin{tabularx}{\linewidth}{|X|X|X|X|}
\hline
\rowcolor{popblueL} Critère & \zh{优秀} (A) & \zh{良好} (B) & \zh{需努力} (C) \\
\hline
Prononciation (Tons) & Tons 1-4 stables, \zh{ü} correct, pas de *jiǎnyù & Quelques hésitations sur tons 3+4 & Tons plats, confusion \zh{tōu/tòu}, \zh{qiǎng/qiáng} \\
Vocabulaire thème & Utilise 8+ mots du tableau (ex. \zh{预防}, \zh{矫正}, \zh{法治教育}) & Utilise 5--7 mots de base & Vocabulaire pauvre, répétitif (\zh{好}, \zh{坏}) \\
Structures & \zh{首先/其次/最后}, \zh{虽然/但是}, \zh{以……为主} & \zh{因为/所以}, \zh{我觉得} simples & Phrases hachées, pas de connecteurs \\
Interaction & Relance (\zh{你怎么看？}), réagit (\zh{有道理}) & Répond aux questions, peu d'initiative & Monologue, pas d'échange \\
\hline
\end{tabularx}
\end{center}

\courssub{Tableau récapitulatif du vocabulaire de la section (complété)}

\begin{center}
\begin{tabularx}{\linewidth}{|X|X|X|X|}
\hline
\rowcolor{popblueL} Caractères & Pinyin & Nature & Traduction \\
\hline
青少年犯罪 & qīngshàonián fànzuì & NM & délinquance juvénile \\
青少年 & qīngshàonián & N & adolescents, jeunes \\
少年 & shàonián & N & jeune, adolescent \\
未成年人 & wèichéngniánrén & N & mineur (< 18 ans) \\
犯罪 & fànzuì & V/N & commettre un crime / criminalité \\
罪犯 & zuìfàn & N & criminel, délinquant \\
偷 & tōu & V & voler (discrètement) \\
抢 & qiǎng & V & arracher, voler (avec violence) \\
抢劫 & qiǎngjié & V/N & braquer / braquage \\
打架 & dǎjià & VO & se battre \\
斗殴 & dòu'ōu & V/N & rixe, bagarre collective \\
法律 & fǎlǜ & N & loi \\
警察 & jǐngchá & N & police, policier \\
法院 & fǎyuàn & N & tribunal \\
监狱 & jiānyù & N & prison \\
家庭教育 & jiātíng jiàoyù & NM & éducation familiale \\
预防 & yùfáng & V & prévenir \\
解决 & jiějué & V & résoudre \\
原因 & yuányīn & N & cause \\
影响 & yǐngxiǎng & V/N & influencer / influence \\
家长 & jiāzhǎng & N & parent d'élève, tuteur \\
老师 & lǎoshī & N & professeur, éducateur \\
\hline
\multicolumn{4}{|c|}{\textbf{Supplément : Formules fonctionnelles \& Débat}} \\
\hline
首先 / 其次 / 最后 & shǒuxiān / qícì / zuìhòu & Adv & Premièrement / Deuxièmement / Enfin \\
我认为 / 依我看来 & wǒ rènwéi / yī wǒ kànlái & V / Adv & Je pense / À mon avis \\
我同意 / 我不太同意 & wǒ tóngyì / wǒ bù tài tóngyì & V & Je suis d'accord / Pas tout à fait d'accord \\
虽然……但是…… & suīrán \dots dànshì \dots & Conj & Bien que \dots cependant \dots \\
你怎么看？ & nǐ zěnme kàn? & Q & Qu'en penses-tu ? \\
有道理 & yǒu dàoli & V & C'est raisonnable / Il y a du sens \\
改过自新 & gǎiguòzìxīn & Idiome & Se réinsérer, se corriger \\
威慑 & wēishè & V/N & Dissuader / Dissuasion \\
矫正 & jiǎozhèng & V/N & Rééduquer / Rééducation \\
社区矫正 & shèqū jiǎozhèng & NM & Rééducation communautaire \\
\hline
\end{tabularx}
\end{center}

\coursec{Leçon 7 — Expression orale : la délinquance juvénile}

\courssub{Vocabulaire du thème : la délinquance juvénile}

Avant d'entrer dans le dialogue, fixons le lexique de base. Ce vocabulaire (niveau HSK 3-4) est indispensable pour nommer les acteurs, les actes, les institutions et les concepts abstraits liés au sujet. Apprenez les caractères, le pinyin \emph{avec les tons} et le sens ; entraînez-vous à écrire les caractères les plus fréquents (\zh{偷}, \zh{抓}, \zh{责}, \zh{教}, \zh{法}) en respectant l'ordre des traits.

\begin{center}
\begin{tabularx}{\linewidth}{|X|X|X|X|}
\hline
\rowcolor{popblueL} Caractères & Pinyin & Nature & Traduction \\
\hline
偷 & tōu & verbe & voler \\
\hline
偷窃 & tōuqiè & verbe/nom & vol (terme plus formel) \\
\hline
抢 & qiǎng & verbe & arracher, voler avec violence \\
\hline
打架 & dǎjià & verbe & se battre, bagarrer \\
\hline
吸毒 & xīdú & verbe & prendre de la drogue \\
\hline
警察 & jǐngchá & nom & police, policier \\
\hline
派出所 & pàichūsuǒ & nom & commissariat de quartier \\
\hline
抓 & zhuā & verbe & attraper, arrêter (par la police) \\
\hline
判 & pàn & verbe & juger, condamner (tribunal) \\
\hline
坐牢 & zuòláo & verbe & faire de la prison \\
\hline
责任 & zérèn & nom & responsabilité \\
\hline
原因 & yuányīn & nom & cause, raison \\
\hline
后果 & hòuguǒ & nom & conséquence \\
\hline
教育 & jiàoyù & nom/verbe & éducation ; éduquer \\
\hline
惩罚 & chéngfá & nom/verbe & punition ; punir \\
\hline
改造 & gǎizào & verbe & rééduquer, réformer (un délinquant) \\
\hline
未成年人 & wèichéngniánrén & nom & mineur (moins de 18 ans) \\
\hline
监护人 & jiānhùrén & nom & tuteur, responsable légal \\
\hline
社会 & shèhuì & nom & société \\
\hline
环境 & huánjìng & nom & environnement \\
\hline
失业 & shīyè & verbe/adj & chômage ; au chômage \\
\hline
辍学 & chuòxué & verbe & abandonner l'école \\
\hline
不良 & bùliáng & adjectif & mauvais, déviant (ex: \zh{不良少年}) \\
\hline
严重 & yánzhòng & adjectif & grave, sérieux \\
\hline
预防 & yùfáng & verbe & prévenir \\
\hline
\end{tabularx}
\end{center}

\begin{aretenir}[Clés d'écriture (ordre des traits)]
\begin{itemize}
\item \zh{偷} (tōu) : clé \zh{人} (rén, homme) à gauche → haut → bas ; partie phonétique \zh{俞} à droite.
\item \zh{抓} (zhuā) : clé \zh{扌} (main) à gauche → horizontal, crochet vertical, point ; phonétique \zh{爪} (griffe) à droite.
\item \zh{责} (zé) : clé \zh{贝} (coquillage/argent) en bas → haut ; haut \zh{朿} (épine).
\item \zh{教} (jiào) : clé \zh{攵} (pū, action de taper) en bas à droite ; haut \zh{孝} (piété filiale).
\item \zh{法} (fǎ) : clé \zh{氵} (eau) à gauche → point, point, trait montant ; droite \zh{去}.
\end{itemize}
\end{aretenir}

\courssub{Dialogue modèle 1 : parler d'un fait divers — observation et analyse}

Dans cette section, nous mettons le vocabulaire en situation. Écoutez (ou lisez) un échange authentique entre deux élèves de Terminale à Yaoundé, Mballa et 李娜 (Lǐ Nà), qui commentent un fait divers survenu dans leur quartier. Ce dialogue illustre la progression logique \textbf{constat → réaction → causes → responsabilités} et introduit la structure passive \zh{被}.

\begin{textebox}[Dialogue 1 — 你听说了吗？]
A（Mballa）：你听说了吗？我们区有几个年轻人偷东西被警察抓了。\\
Nǐ tīngshuō le ma? Wǒmen qū yǒu jǐ ge niánqīngrén tōu dōngxi bèi jǐngchá zhuā le.\\
« Tu as entendu la nouvelle ? Il y a quelques jeunes de notre quartier qui ont volé des choses et ont été arrêtés par la police. »\\[4pt]
B（李娜）：真的吗？太可怕了！他们为什么这样做？\\
Zhēn de ma? Tài kěpà le! Tāmen wèishénme zhèyàng zuò?\\
« Vraiment ? C'est terrible ! Pourquoi ont-ils fait cela ? »\\[4pt]
A（Mballa）：我觉得因为他们没有工作，也没有好好上学。\\
Wǒ juéde yīnwèi tāmen méiyǒu gōngzuò, yě méiyǒu hǎohāo shàngxué.\\
« Je pense que c'est parce qu'ils n'ont pas de travail et qu'ils n'ont pas non plus étudié sérieusement. »\\[4pt]
B（李娜）：家庭也有责任。父母应该关心孩子。\\
Jiātíng yě yǒu zérèn. Fùmǔ yīnggāi guānxīn háizi.\\
« La famille a aussi sa responsabilité. Les parents devraient prendre soin de leurs enfants. »
\end{textebox}

\begin{popfigure}
\begin{tikzpicture}[
  bulle/.style={rounded corners=8pt, draw=#1, fill=#1!15, text width=3.1cm, align=center, minimum height=1.7cm, font=\small},
  fleche/.style={-{Stealth[length=3mm]}, thick, popdark}
]
\node[bulle=popblue, align=center] (c1) at (0,0){\textbf{1. Constat}\\ \zh{你听说了吗？…被警察抓了。}\\ \emph{annoncer le fait}};
\node[bulle=poporange, right=1.1cm of c1, align=center] (c2){\textbf{2. Réaction}\\ \zh{真的吗？太可怕了！}\\ \emph{réagir, s'étonner}};
\node[bulle=popgreen, right=1.1cm of c2, align=center] (c3){\textbf{3. Causes}\\ \zh{我觉得因为…}\\ \emph{expliquer pourquoi}};
\node[bulle=poppurple, right=1.1cm of c3, align=center] (c4){\textbf{4. Responsabilités}\\ \zh{家庭也有责任。…应该…}\\ \emph{proposer, juger}};
\draw[fleche] (c1) -- (c2);
\draw[fleche] (c2) -- (c3);
\draw[fleche] (c3) -- (c4);
\end{tikzpicture}
\legende{La progression du dialogue : constat, réaction, causes, responsabilités.}
\end{popfigure}

\begin{aretenir}[Un plan à réutiliser : les 4 étapes de l'échange sur un fait de société]
\textbf{1. Constat} : \zh{你听说了吗？} + fait (\zh{…被…了}).\\
\textbf{2. Réaction} : \zh{真的吗？} + adjectif évaluatif (\zh{太…了！}).\\
\textbf{3. Causes} : \zh{我觉得因为…} + phrase complète (sujet + verbe).\\
\textbf{4. Responsabilités / Solutions} : \zh{…也有责任。…应该/不应该…}.
\end{aretenir}

\courssub{Point de langue 1 : le passif avec 被 (bèi) — structure et emploi}

\textbf{Observation.} Dans la phrase de Mballa \zh{年轻人偷东西被警察抓了}, le sujet \zh{年轻人} (les jeunes) \emph{subit} l'action d'être arrêtés. En français, on utilise « être + participe passé » (voix passive). En chinois, la marque \zh{被} (bèi) placée \emph{devant l'agent} signale la passivité.

\begin{propriete}[Structure du passif avec 被 bèi]
\textbf{Schéma :} \textbf{Patient (Sujet)} + \zh{被} + \textbf{Agent} + \textbf{Verbe} + \textbf{\zh{了} / Complément de résultat / Complément de degré}.
\begin{itemize}
\item \textbf{Patient} : celui qui subit l'action (devient le sujet grammatical).
\item \zh{被} (bèi) : préposition marquant la passivité (équivaut à « par »).
\item \textbf{Agent} : celui qui fait l'action (peut être omis si inconnu ou évident).
\item \textbf{Verbe + suite} : \textbf{interdit de laisser le verbe nu} après \zh{被}. Il faut \zh{了} (accompli), un complément de résultat (\zh{打破}, \zh{吃完}) ou un complément de degré (\zh{批评得很惨}).
\end{itemize}
\end{propriete}

\begin{exemplebox}[Le passif 被 en contexte]
\zh{我的自行车被偷了。} Wǒ de zìxíngchē bèi tōu le.\\
« Mon vélo a été volé. » (Agent omis : on ne sait pas par qui.)\\[4pt]
\zh{他被老师批评了。} Tā bèi lǎoshī pīpíng le.\\
« Il a été critiqué par le professeur. »\\[4pt]
\zh{手机被弟弟打破了。} Shǒujī bèi dìdi dǎpò le.\\
« Le téléphone a été cassé par le petit frère. » (Résultat : \zh{打破})\\[4pt]
\zh{小偷被当场抓获。} Xiǎotōu bèi dāngchǎng zhuāhuò.\\
« Le voleur a été arrêté sur le champ. » (Registre plus formel : \zh{抓获}).
\end{exemplebox}

\begin{attention}[Pièges majeurs pour francophones]
\begin{enumerate}
\item \textbf{Ne pas utiliser \zh{是} (shì) pour le passif.} \zh{*他是警察抓的} (sauf construction \zh{是…的} pour l'emphase sur le passé, hors programme ici). La marque passive est \zh{被} (ou \zh{叫}, \zh{让} à l'oral familier).
\item \textbf{Verbe nu interdit.} \zh{*小偷被警察抓} $\rightarrow$ Faux. Dire \zh{小偷被警察抓了} ou \zh{小偷被警察抓住了} (résultat).
\item \textbf{Agent inanimé.} \zh{被} s'utilise surtout avec un agent animé (personne, animal). Pour une cause naturelle/inanimée, on préfère \zh{由于} (yóuyú) + cause + \zh{导致}… ou structure active : \zh{大雨淹了房子} (La pluie a inondé la maison) plutôt que \zh{房子被大雨淹了} (moins naturel).
\end{enumerate}
\end{attention}

\begin{exoresolu}[Transformation active → passif]
Énoncé : Transformez \zh{警察抓了小偷。} (Jǐngchá zhuā le xiǎotōu.) en phrase passive avec \zh{被}.
\tcblower
\textbf{Solution détaillée :}
\begin{enumerate}
\item Identifier le patient (objet) : \zh{小偷} (le voleur). Il devient le sujet.
\item Identifier l'agent (sujet) : \zh{警察} (la police). Il suit \zh{被}.
\item Garder le verbe + \zh{了} : \zh{抓了}.
\item Résultat : \zh{小偷被警察抓了。} Xiǎotōu bèi jǐngchá zhuā le.
\end{enumerate}
\end{exoresolu}

\courssub{Point de langue 2 : exprimer la cause avec 因为 (yīnwèi) et 所以 (suǒyǐ)}

\textbf{Observation.} À la question \zh{他们为什么这样做？}, Mballa répond : \zh{我觉得因为他们没有工作，也没有好好上学}. La cause est introduite par \zh{因为}.

\begin{propriete}[Structure causale 因为… (所以…)]
\textbf{Schéma :} \zh{因为} + \textbf{Cause (phrase complète)}, (\zh{所以}) + \textbf{Conséquence (phrase complète)}.
\begin{itemize}
\item \zh{因为} (yīnwèi) = « parce que ». Introduit une \textbf{proposition entière} (sujet + verbe), jamais un nom seul.
\item \zh{所以} (suǒyǐ) = « donc, par conséquent ». Souvent corrélé à \zh{因为} dans la même phrase (correct en chinois, fautif en français : *\emph{Parce que… donc…}*).
\item On peut n'utiliser que \zh{因为} si la conséquence est implicite ou déjà énoncée (cas du dialogue).
\end{itemize}
\end{propriete}

\begin{exemplebox}[Parce que / Donc]
\zh{因为他没钱，所以他去偷东西。} Yīnwèi tā méi qián, suǒyǐ tā qù tōu dōngxi.\\
« Comme il n'avait pas d'argent, il est allé voler. »\\[4pt]
\zh{因为家庭环境不好，所以他辍学了。} Yīnwèi jiātíng huánjìng bù hǎo, suǒyǐ tā chuòxué le.\\
« À cause d'un mauvais environnement familial, il a décroché. »\\[4pt]
\zh{我觉得因为缺乏教育。} Wǒ juéde yīnwèi quēfá jiàoyù.\\
« Je pense que c'est à cause du manque d'éducation. » (Conséquence implicite).
\end{exemplebox}

\begin{attention}[Erreur fréquente : \zh{因为} + nom seul]
\textbf{Faux :} \zh{*因为贫穷，他偷东西。} (Parce que la pauvreté, il vole… — \zh{贫穷} est un nom/adjectif, pas une phrase).
\textbf{Correct :} \zh{因为他很穷，所以他偷东西。} (Parce qu'\textbf{il est pauvre}…).
\textbf{Alternative avec nom :} \zh{因为贫穷的关系，他偷东西。} Yīnwèi pínqióng de guānxi… (« À cause de la pauvreté… » — \zh{关系} nominalise).
\end{attention}

\courssub{Point de langue 3 : le devoir moral et le conseil avec 应该 (yīnggāi)}

\textbf{Observation.} 李娜 conclut : \zh{父母应该关心孩子}. \zh{应该} exprime une norme sociale, un devoir moral, un conseil avisé — l'équivalent du conditionnel français « devrait / devraient ».

\begin{propriete}[应该 yīnggāi — devoir moral / conseil]
\textbf{Structure affirmative :} Sujet + \zh{应该} + Verbe (+ Complément).
\textbf{Structure négative :} Sujet + \zh{不应该} + Verbe (+ Complément).
\begin{itemize}
\item \zh{应该} = « devoir » (opinion, morale, bon sens). Plus doux que \zh{必须} (bìxū, obligation stricte, contrainte externe).
\item \zh{不应该} = « ne pas devoir / ne pas avoir à » (désapprobation morale).
\end{itemize}
\end{propriete}

\begin{center}
\begin{tabularx}{\linewidth}{|X|X|X|}
\hline
\rowcolor{popblueL} Structure & Exemple (caractères + pinyin) & Traduction \\
\hline
被 + Agent + V + 了 & \zh{小偷被警察抓了。} Xiǎotōu bèi jǐngchá zhuā le. & Le voleur a été arrêté par la police. \\
\hline
因为 + Cause, 所以 + Cons. & \zh{因为没上学，所以找不到工作。} Yīnwèi méi shàngxué, suǒyǐ zhǎo bu dào gōngzuò. & Comme il n'a pas étudié, il ne trouve pas de travail. \\
\hline
应该 + V & \zh{父母应该关心孩子。} Fùmǔ yīnggāi guānxīn háizi. & Les parents devraient s'occuper de leurs enfants. \\
\hline
不应该 + V & \zh{年轻人不应该打架。} Niánqīngrén bù yīnggāi dǎjià. & Les jeunes ne devraient pas se battre. \\
\hline
必须 + V & \zh{罪犯必须坐牢。} Zuìfàn bìxū zuòláo. & Les criminels doivent (obligation légale) aller en prison. \\
\hline
\end{tabularx}
\end{center}

\courssub{Point de langue 4 : la question « pourquoi » — 为什么 (wèishénme)}

\begin{propriete}[Position de 为什么]
En chinois, \zh{为什么} reste \textbf{à l'intérieur de la phrase}, après le sujet, avant le verbe. Il ne se met jamais en tête de phrase comme « Pourquoi » en français.
\textbf{Schéma :} Sujet + \zh{为什么} + Verbe (+ Complément) + ?\\
Exemple : \zh{他们为什么偷东西？} Tāmen wèishénme tōu dōngxi? (Eux pourquoi voler choses ?)
\end{propriete}

\begin{attention}[Ordre des mots : Question « Pourquoi »]
\textbf{Faux (calque français) :} \zh{*为什么他们偷东西？}
\textbf{Correct :} \zh{他们为什么偷东西？}
\end{attention}

\courssub{Phonétique : les tons à travailler — exercice ciblé}

La prononciation correcte des tons est notée à l'oral. Voici les 4 mots « pièges » du dialogue 1, avec leur profil tonal et un exercice de discrimination.

\begin{center}
\begin{tabularx}{\linewidth}{|X|X|X|X|}
\hline
\rowcolor{popblueL} Mot & Pinyin & Profil des tons & Exercice de discrimination (lire à voix haute) \\
\hline
听说 & tīngshuō & 1 + 1 (haut plat ×2) & \zh{听说} — \zh{同事} (tóngshì) — \zh{天书} (tiānshū) \\
\hline
可怕 & kěpà & 3 + 4 (bas descendant + chute brutale) & \zh{可怕} — \zh{可怜} (kělián) — \zh{客帕} (kèpà, inexistant) \\
\hline
责任 & zérèn & 2 + 4 (montant + chute) & \zh{责任} — \zh{则人} (zérén) — \zh{贼人} (zéirén) \\
\hline
关心 & guānxīn & 1 + 1 (haut plat ×2) & \zh{关心} — \zh{观音} (Guānyīn) — \zh{关新} (guānxīn, autre sens) \\
\hline
\end{tabularx}
\end{center}

\begin{methode}[Entraînement tonal en 3 passes (à faire en binôme)]
\begin{enumerate}
\item \textbf{Isolation :} Répétez chaque mot du tableau 3 fois de suite en exagérant le contour tonal (main qui monte/descend).
\item \textbf{En phrase :} Relisez le dialogue 1 \emph{uniquement en pinyin}, en marquant les 4 mots cibles. Enregistrez-vous si possible.
\item \textbf{Échange :} Jouez le dialogue 1 avec un camarade. L'écouteur note les erreurs de tons sur une grille (1=haut, 2=montant, 3=bas, 4=descendant). Changez de rôle.
\end{enumerate}
\end{methode}

\courssub{Formules fonctionnelles pour le débat : présenter, opiner, relancer}

Pour passer du dialogue simple au débat structuré (section suivante), vous devez maîtriser ces « briques » orales. Classez-les par fonction ; mémorisez les 2-3 phrases par catégorie.

\begin{definition}[Boîte à outils orale — Débat « Délinquance : punir ou éduquer ? »]
\begin{center}
\begin{tabularx}{\linewidth}{|X|X|X|}
\hline
\rowcolor{popblueL} Fonction & Chinois (Caractères + Pinyin) & Français \\
\hline
\multicolumn{3}{|c|}{\cellcolor{popblueL!50}\textbf{1. Ouvrir / Annoncer le sujet}} \\
\hline
Lancer le sujet & \zh{今天我们讨论一下未成年人犯罪。} Jīntiān wǒmen tǎolùn yíxià wèichéngniánrén fànzuì. & Aujourd'hui, discutons de la délinquance des mineurs. \\
Citer un fait & \zh{最近我们这里发生了一件事……} Zuìjìn wǒmen zhèli fāshēng le yí jiàn shì… & Récemment, il s'est passé un fait ici… \\
\hline
\multicolumn{3}{|c|}{\cellcolor{poporangeL!50}\textbf{2. Donner son avis / Prendre position}} \\
\hline
Opinion personnelle & \zh{我认为…… / 我觉得……} Wǒ rènwéi… / Wǒ juéde… & Je pense que… / Je trouve que… \\
Insister & \zh{在我看来，……} Zài wǒ kàn lái… & À mon avis… \\
Nuancer & \zh{一方面……，另一方面……} Yīfāngmiàn…, lìngyīfāngmiàn… & D'un côté…, de l'autre… \\
\hline
\multicolumn{3}{|c|}{\cellcolor{popgreenL!50}\textbf{3. Argumenter (Cause / Conséquence / But)}} \\
\hline
Cause & \zh{因为……所以……} Yīnwèi… suǒyǐ… & Parce que… donc… \\
But & \zh{为了……，我们应该……} Wèile…, wǒmen yīnggāi… & Pour…, nous devrions… \\
Exemple & \zh{比如说……} Bǐrú shuō… & Par exemple… \\
\hline
\multicolumn{3}{|c|}{\cellcolor{poppurpleL!50}\textbf{4. Réagir / Approuver / Contester}} \\
\hline
Approuver & \zh{我赞同。/ 说得有道理。} Wǒ zàntóng. / Shuō de yǒu dàolǐ. & Je suis d'accord. / Ce que tu dis a du sens. \\
Contester (doux) & \zh{我不太同意。/ 这不一定吧？} Wǒ bù tài tóngyì. / Zhè bù yídìng ba? & Je ne suis pas tout à fait d'accord. / Pas sûr… \\
Contester (fort) & \zh{这不对。/ 我不认为这样。} Zhè bù duì. / Wǒ bù rènwéi zhèyàng. & C'est faux. / Je ne pense pas ça. \\
Demander précision & \zh{你是什么意思？} Nǐ shì shénme yìsi? & Qu'est-ce que tu veux dire ? \\
\hline
\multicolumn{3}{|c|}{\cellcolor{poppinkL!50}\textbf{5. Relancer / Passer la parole}} \\
\hline
Inviter & \zh{你怎么看？/ 你有什么建议？} Nǐ zěnme kàn? / Nǐ yǒu shénme jiànyì? & Qu'en penses-tu ? / Quelle est ta proposition ? \\
Rebondir & \zh{除此之外，……} Chúle zhīwài… & En plus de cela… \\
Conclure & \zh{总之，……} Zǒngzhī… & En résumé… / Bref… \\
\hline
\end{tabularx}
\end{center}
\end{definition}

\courssub{Dialogue modèle 2 : débat — punir ou éduquer ?}

Ce dialogue met en scène trois élèves (Mballa, 李娜, 阿明 Āmíng) et utilise les formules ci-dessus. Observez comment l'échange rebondit.

\begin{textebox}[Dialogue 2 — 惩罚还是教育？]
A（Mballa）：今天我们讨论一下：未成年人偷东西，应该判刑还是教育？\\
Jīntiān wǒmen tǎolùn yíxià: wèichéngniánrén tōu dōngxi, yīnggāi pànxíng háishi jiàoyù?\\
« Aujourd'hui, discutons : pour un mineur qui vole, faut-il condamner ou éduquer ? »\\[4pt]
B（李娜）：我认为应该以教育为主。因为他们还小，不懂法。\\
Wǒ rènwéi yīnggāi yǐ jiàoyù wéi zhǔ. Yīnwèi tāmen hái xiǎo, bù dǒng fǎ.\\
« Je pense qu'il faut privilégier l'éducation. Parce qu'ils sont encore jeunes, ils ne connaissent pas la loi. »\\[4pt]
C（阿明）：我不太同意。如果不惩罚，他们会再犯。法律面前人人平等。\\
Wǒ bù tài tóngyì. Rúguǒ bù chéngfá, tāmen huì zàifàn. Fǎlǜ miànqián rénrén píngděng.\\
« Je ne suis pas tout à fait d'accord. Si on ne punit pas, ils récidiveront. Devant la loi, tous sont égaux. »\\[4pt]
B（李娜）：你是什么意思？把孩子送进监狱吗？那样后果更严重。\\
Nǐ shì shénme yìsi? Bǎ háizi sòng jìn jiānyù ma? Nà yàng hòuguǒ gèng yánzhòng.\\
« Qu'est-ce que tu veux dire ? Envoyer les enfants en prison ? Les conséquences seraient encore plus graves. »\\[4pt]
A（Mballa）：一方面要依法处理，另一方面要找原因。家庭、学校、社会都有责任。\\
Yīfāngmiàn yào yīfǎ chǔlǐ, lìngyīfāngmiàn yào zhǎo yuányīn. Jiātíng, xuéxiào, shèhuì dōu yǒu zérèn.\\
« D'un côté, il faut traiter selon la loi ; de l'autre, chercher les causes. Famille, école, société : toutes ont leur responsabilité. »\\[4pt]
C（阿明）：说得有道理。总之，预防比惩罚更重要。\\
Shuō de yǒu dàolǐ. Zǒngzhī, yùfáng bǐ chéngfá gèng zhòngyào.\\
« Ce que tu dis a du sens. En résumé, la prévention est plus importante que la punition. »
\end{textebox}

\begin{aretenir}[Vocabulaire clé du Dialogue 2]
\begin{center}
\begin{tabularx}{\linewidth}{|X|X|X|X|}
\hline
\rowcolor{popblueL} Caractères & Pinyin & Nature & Traduction \\
\hline
讨论 & tǎolùn & verbe & discuter, débattre \\
\hline
判刑 & pànxíng & verbe/nom & condamner à une peine / condamnation \\
\hline
以…为主 & yǐ… wéi zhǔ & structure & privilégier…, faire de… la priorité \\
\hline
不懂法 & bù dǒng fǎ & verbe & ne pas connaître la loi / l'ignorer \\
\hline
再犯 & zàifàn & verbe & récidiver \\
\hline
法律面前人人平等 & fǎlǜ miànqián rénrén píngděng & chengyu & devant la loi, tous sont égaux \\
\hline
依法处理 & yīfǎ chǔlǐ & verbe & traiter selon la loi / conformément à la loi \\
\hline
预防 & yùfáng & verbe/nom & prévenir / prévention \\
\hline
\end{tabularx}
\end{center}
\end{aretenir}

\courssub{Jeux de rôle et débat guidé — mise en situation}

\courssub{Activité 1 : Jeu de rôle « Au commissariat » (3 personnes)}

\begin{experience}[Scénario : Vol à l'étalage au Marché Central de Yaoundé]
\textbf{Rôles :}
\begin{itemize}
\item \textbf{Élève A (Le jeune / Le suspect)} : \zh{王小明} (Wáng Xiǎomíng), 16 ans, élève en 1ère, a volé un téléphone portable. Il est nerveux, cherche des excuses (pauvreté, pression copains).
\item \textbf{Élève B (Le policier)} : \zh{张警官} (Zhāng Jǐnguān), calme, professionnel. Utilise le passif \zh{被} pour constater les faits, pose des questions \zh{为什么}, informe des conséquences (\zh{判刑}, \zh{坐牢}, \zh{通知监护人}).
\item \textbf{Élève C (Le parent / Tuteur)} : \zh{王父} (Wáng Fù) ou \zh{王母} (Wáng Mǔ), arrive plus tard. Inquiet, colère, puis responsable. Utilise \zh{应该} pour parler de l'éducation future (\zh{以后应该好好管教}).
\end{itemize}
\textbf{Contraintes linguistiques obligatoires :}
\begin{enumerate}
\item Au moins 2 phrases au passif \zh{被} (ex: \zh{手机被你偷了}, \zh{你被抓住了}).
\item Au moins 1 structure \zh{因为…所以…}.
\item Au moins 2 \zh{应该 / 不应该}.
\item Utiliser le vocabulaire : \zh{派出所}, \zh{监护人}, \zh{未成年人}, \zh{后果}.
\end{enumerate}
\textbf{Déroulement :} 5 min préparation (notes en pinyin autorisées) → 3-4 min jeu devant la classe → 2 min feedback pair (tons, structures, vocabulaire).
\end{experience}

\courssub{Activité 2 : Débat structuré « Prison ou Centre de rééducation ? » (Groupe de 4-6)}

\begin{experience}[Débat guidé — Sujet : \zh{未成年人犯罪，应该坐牢还是去少管所？}]
\textbf{Préparation (10 min) :} Chaque groupe tire un rôle \textbf{POUR la prison} ou \textbf{POUR le centre de rééducation (少管所 shàoguǎnsuǒ)}. Préparez 3 arguments avec \zh{因为…所以…} et 1 contre-argument anticipé.
\textbf{Vocabulaire d'appoint :}
\begin{center}
\begin{tabularx}{\linewidth}{|X|X|X|}
\hline
\rowcolor{popblueL} Camp & Argument type (chinois) & Traduction \\
\hline
\zh{坐牢} & \zh{因为犯罪就要受惩罚，所以应该坐牢。} & Crime $\rightarrow$ punition $\rightarrow$ prison. \\
& \zh{不坐牢，别人会效仿，社会不安全。} & Pas de prison $\rightarrow$ imitation $\rightarrow$ insécurité. \\
\hline
\zh{少管所 / 教育} & \zh{因为未成年人心理不成熟，所以应该教育改造。} & Mineurs immatures $\rightarrow$ éducation/réforme. \\
& \zh{坐牢后果太严重，毁了前程。} & Prison $\rightarrow$ conséquences graves, avenir gâché. \\
\hline
\end{tabularx}
\end{center}
\textbf{Déroulement du débat (15 min) :}
\begin{enumerate}
\item \textbf{Tour de table (2 min) :} Chaque membre donne son avis principal (\zh{我认为…因为…所以…}).
\item \textbf{Échange libre (10 min) :} Utilisez les formules de relance (\zh{你怎么看？}, \zh{我不太同意…}, \zh{除此之外…}). Le professeur note : justesse des tons, structures \zh{被}/\zh{因为}/\zh{应该}, richesse vocabulaire.
\item \textbf{Synthèse orale (3 min) :} Un rapporteur par groupe résume en 30 sec : \zh{总之，我们组认为……} (En résumé, notre groupe pense que…).
\end{enumerate}
\end{experience}

\courssub{Expression écrite guidée : compte-rendu de débat}

\begin{exercice}[Rédaction : Compte-rendu du débat (80-100 caractères)]
Rédigez un court paragraphe en chinois (caractères) résumant la position de votre groupe lors de l'Activité 2. Structure imposée :
\begin{enumerate}
\item Sujet : \zh{关于未成年人犯罪……} (Guānyú wèichéngniánrén fànzuì…)
\item Position : \zh{我们组认为应该……} (Wǒmen zǔ rènwéi yīnggāi…)
\item Argument principal : \zh{因为……所以……} (Yīnwèi… suǒyǐ…)
\item Mesure concrète : \zh{建议……} (Jiànyì…) ou \zh{应该……} (Yīnggāi…)
\end{enumerate}
\textbf{Exemple modèle :}
\zh{关于未成年人犯罪，我们组认为应该以教育为主。因为未成年人心理不成熟，所以坐牢后果太严重。建议加强家庭教育和学校预防，少管所应该多安排心理辅导和技能培训。}
\end{exercice}

\begin{corrige}[Exercice « Compte-rendu de débat » — Modèle corrigé]
\zh{关于未成年人犯罪，我们组认为应该以教育为主，辅以适当惩罚。因为未成年人心理不成熟，不懂法，所以单纯坐牢后果太严重，容易毁了前程。建议家庭、学校、社会共同负责，加强预防，少管所应重点进行心理辅导和职业技能培训，帮助他们改过自新。}\\
\textbf{Points vérifiés :} \zh{以…为主，辅以…} (structure avancée), \zh{因为…所以…}, \zh{应该} x2, \zh{建议} + VV, vocabulaire thème (\zh{心理辅导} xīnlǐ fúndǎo, \zh{职业技能} zhíyè jìnéng, \zh{改过自新} gǎiguòzìxīn).
\end{corrige}

\courssub{Exercices de consolidation (écrit / traduction)}

\begin{exercice}[Entraînement écrit — Structures ciblées]
\begin{enumerate}
\item \textbf{Passif \zh{被} :} Mettez au passif (avec agent ou sans agent).
  \begin{enumerate}
  \item \zh{老师批评了那个学生。} (Le prof a critiqué cet élève.)
  \item \zh{小偷偷了我的钱包。} (Le voleur a volé mon portefeuille.)
  \item \zh{大雨淹了农田。} (La forte pluie a inondé les champs — \emph{choisir la structure la plus naturelle}.)
  \end{enumerate}
\item \textbf{Cause \zh{因为…所以…} :} Reliez les deux phrases.
  \begin{enumerate}
  \item Il a abandonné l'école. Il ne trouve pas de travail. (\zh{辍学}, \zh{找不到工作})
  \item La famille ne s'occupe pas de lui. Il est devenu délinquant. (\zh{关心}, \zh{变成}, \zh{不良少年})
  \end{enumerate}
\item \textbf{Devoir \zh{应该} :} Traduisez en chinois.
  \begin{enumerate}
  \item Les parents devraient passer plus de temps avec leurs enfants. (\zh{陪伴} péibàn)
  \item Les jeunes ne devraient pas imiter les mauvais comportements. (\zh{模仿} mófǎng, \zh{行为} xíngwéi)
  \item La société doit (obligation forte) protéger les mineurs. (\zh{保护} bǎohù, \zh{必须})
  \end{enumerate}
\item \textbf{Traduction version (Chinois $\rightarrow$ Français) :}
  \zh{最近我们社区发生了一起未成年人打架事件。几个年轻人因为琐事在球场打架，被派出所警察带走了。我觉得家庭应该负主要责任，父母应该好好管教孩子。学校也应该加强法制教育。总之，预防比惩罚重要。}
\end{enumerate}
\end{exercice}

\begin{corrige}[Exercices de consolidation]
\begin{enumerate}
\item \textbf{Passif :}
  \begin{enumerate}
  \item \zh{那个学生被老师批评了。} Nà ge xuéshēng bèi lǎoshī pīpíng le.
  \item \zh{我的钱包被小偷偷了。} Wǒ de qiánbāo bèi xiǎotōu tōu le. (Agent gardé) ou \zh{我的钱包被偷了。} (Agent omis).
  \item \textbf{Attention :} Agent inanimé (\zh{大雨}). Préférer l'actif : \zh{大雨淹了农田。} Ou \zh{由于大雨，农田被淹了。} (Yóuyú dàyǔ, nóngtián bèi yān le — \zh{由于} + cause, passif possible mais moins spontané).
  \end{enumerate}
\item \textbf{Cause :}
  \begin{enumerate}
  \item \zh{因为他辍学了，所以他找不到工作。} Yīnwèi tā chuòxué le, suǒyǐ tā zhǎo bu dào gōngzuò.
  \item \zh{因为家庭不关心他，所以他变成了不良少年。} Yīnwèi jiātíng bù guānxīn tā, suǒyǐ tā biànchéng le bùliáng shàonián.
  \end{enumerate}
\item \textbf{Devoir :}
  \begin{enumerate}
  \item \zh{父母应该多陪伴孩子。} Fùmǔ yīnggāi duō péibàn háizi.
  \item \zh{年轻人不应该模仿不良行为。} Niánqīngrén bù yīnggāi mófǎng bùliáng xíngwéi.
  \item \zh{社会必须保护未成年人。} Shèhuì bìxū bǎohù wèichéngniánrén.
  \end{enumerate}
\item \textbf{Version :}
« Récemment, une affaire de bagarre entre mineurs s'est produite dans notre quartier. Quelques jeunes se sont battus sur le terrain de sport pour une futilité et ont été emmenés au commissariat par la police. Je pense que la famille doit assumer la responsabilité principale, les parents devraient bien éduquer leurs enfants. L'école aussi devrait renforcer l'éducation juridique. En résumé, la prévention est plus importante que la punition. »
\end{enumerate}
\end{corrige}

\courssub{Bilan de la leçon — Compétences validées}

\begin{aretenir}[Check-list de fin de leçon — Leçon 7]
\begin{itemize}
\item \textbf{Vocabulaire :} Je connais 20+ mots du thème (actes, acteurs, institutions, concepts abstraits) et j'en écris 10 correctement.
\item \textbf{Grammaire :} Je construis le passif \zh{被} (patient + \zh{被} + agent + V + \zh{了}/résultat), la cause \zh{因为…所以…}, le devoir \zh{应该}/\zh{不应该}/\zh{必须}, la question \zh{为什么}.
\item \textbf{Oral :} Je suis capable de suivre le plan \textbf{Constat → Réaction → Causes → Responsabilités} pour commenter un fait divers, et d'utiliser les formules fonctionnelles pour débattre (approuver, contester, relancer, conclure).
\item \textbf{Phonétique :} Je prononce correctement les tons 1+1 (\zh{听说}, \zh{关心}), 3+4 (\zh{可怕}), 2+4 (\zh{责任}) en phrase.
\item \textbf{Écrit :} Je peux rédiger un compte-rendu de débat structuré (80-100 caractères) et traduire un paragraphe sur le thème.
\item \textbf{Socioculturel :} Je peux comparer les approches camerounaise et chinoise (rééducation \zh{少管所} vs prison, rôle de la famille \zh{监护人}, éducation morale \zh{德育}).
\end{itemize}
\end{aretenir}

\begin{savaistu}[Regard croisé Cameroun — Chine : la justice des mineurs]
\begin{itemize}
\item \textbf{Cameroun :} Majorité pénale à 18 ans (Code pénal, art. 80). Mesures de protection, d'assistance, de surveillance et d'éducation (placement en centre d'observation, liberté surveillée). Peine privative de liberté possible mais exceptionnelle pour les 13-18 ans.
\item \textbf{Chine :} Majorité pénale à 16 ans (14 ans pour crimes graves : homicide, vol à main armée, viol, trafic de drogue — réforme 2021). Système distinct : \zh{少年法庭} (tribunaux pour mineurs), \zh{少管所} (centres de rééducation par le travail/études), \zh{封存犯罪记录} (scellement du casier judiciaire à 18 ans sous conditions). Principe : \zh{教育为主、惩罚为辅} (éducation prioritaire, punition accessoire).
\item \textbf{Point commun :} Reconnaissance de l'immaturité, primat de la réinsertion sur la répression, rôle clé de la famille (\zh{监护人} / tuteur) et de l'école.
\end{itemize}
\end{savaistu}

\coursec{Formules fonctionnelles : présenter, opiner, relancer}

Après avoir étudié le dialogue modèle 1 où deux élèves évoquaient le phénomène de la délinquance juvénile, nous disposons maintenant du vocabulaire de base. Pour passer de la simple évocation à un véritable exposé structuré ou à un débat argumenté, il faut maîtriser les \cle{formules fonctionnelles} : ces « briques » du discours qui permettent d'organiser sa pensée, d'exprimer son opinion avec nuance, d'interagir avec ses interlocuteurs et de conclure efficacement. En classe de Première, l'examen oral (épreuve du Bac) attend de vous que vous sachiez \emph{prendre la parole} de façon cohérente pendant 2 à 3 minutes, puis \emph{réagir} aux questions ou aux arguments des autres. Cette section vous donne la boîte à outils complète.

\courssub{Structurer un exposé oral : salutation, annonce du sujet et plan}

\begin{textebox}{Mini-exposé modèle : structure formelle}
大家好，今天我要谈的问题是青少年犯罪。\\
Dàjiā hǎo, jīntiān wǒ yào tán de wèntí shì qīngshàonián fànzuì.\\
Bonjour à tous, aujourd'hui je vais parler du problème de la délinquance juvénile.\\

首先，我要分析原因。家庭教育缺失、网络不良影响和学校压力太大是主要原因。\\
Shǒuxiān, wǒ yào fēnxī yuányīn. Jiātíng jiàoyù quēshī, wǎngluò bùliáng yǐngxiǎng hé xuéxiào yālì tài dà shì zhǔyào yuányīn.\\
D'abord, je vais analyser les causes. Le manque d'éducation familiale, la mauvaise influence d'Internet et la pression scolaire trop forte sont les causes principales.\\

然后，我要提出建议。家长应该多陪伴孩子，学校要加强心理辅导。\\
Ránhòu, wǒ yào tí chū jiànyì. Jiāzhǎng yīnggāi duō péibàn háizi, xuéxiào yào jiāqiáng xīnlǐ fǔdǎo.\\
Ensuite, je vais proposer des suggestions. Les parents devraient davantage accompagner leurs enfants, l'école doit renforcer le soutien psychologique.\\

最后，我想呼吁全社会共同关注青少年成长。\\
Zuìhòu, wǒ xiǎng hūyù quán shèhuì gòngtóng guānzhù qīngshàonián chéngzhǎng.\\
Enfin, je souhaite lancer un appel à toute la société pour qu'elle porte ensemble attention à la croissance des adolescents.\\

总之，预防青少年犯罪，教育比惩罚更重要。\\
Zǒngzhī, yùfáng qīngshàonián fànzuì, jiàoyù bǐ chéngfá gèng zhòngyào.\\
En résumé, pour prévenir la délinquance juvénile, l'éducation est plus importante que la punition.
\end{textebox}

\begin{propriete}[Structure type d'un mini-exposé scolaire (4 temps)]
\textbf{1. Salutation + Annonce du sujet} \\
\zh{大家好，今天我要谈的问题是……} (Dàjiā hǎo, jīntiān wǒ yào tán de wèntí shì……) — « Bonjour, aujourd'hui je vais parler de… » \\
Autres variantes : \zh{大家好，我的发言主题是……} (Dàjiā hǎo, wǒ de fāyán zhǔtí shì……) — « Bonjour, mon sujet d'intervention est… »

\textbf{2. Annonce du plan (marqueurs séquentiels obligatoires)} \\
\zh{首先……} (Shǒuxiān……) — « D'abord / En premier lieu… » \\
\zh{然后……} (Ránhòu……) — « Ensuite / Puis… » \\
\zh{最后……} (Zuìhòu……) — « Enfin / Pour terminer… » \\
\emph{Note :} En chinois, l'énoncé du plan précède souvent le développement de chaque partie (méthode déductive), contrairement au français où l'on peut dire « Je vais faire trois parties : d'abord…, ensuite…, enfin… » avant de développer. Ici, on dit \zh{首先，我要分析原因。} puis on développe l'analyse.

\textbf{3. Développement + Avis personnel} \\
Utilisation des formules d'opinion (voir section 2) et d'argumentation (\zh{因为……所以……}, \zh{比如说}, \zh{此外}).

\textbf{4. Conclusion formelle} \\
\zh{总之，我认为……} (Zǒngzhī, wǒ rènwéi……) — « En résumé, je pense que… » \\
\zh{最后，我想说……} (Zuìhòu, wǒ xiǎng shuō……) — « Pour finir, je veux dire… »
\end{propriete}

\begin{methode}[Construire son exposé en 4 étapes (Bac Oral)]
\begin{enumerate}
    \item \textbf{Accroche \& Sujet} : Écrivez votre phrase d'ouverture avec \zh{今天我要谈的问题是} + [Nom du thème]. Ex: \zh{青少年犯罪}, \zh{家庭教育}, \zh{网络成瘾}.
    \item \textbf{Plan en 3 points} : Identifiez 3 idées forces (Causes, Conséquences, Solutions / Famille, École, Société). Rédigez les 3 phrases commençant par \zh{首先}, \zh{然后}, \zh{最后}.
    \item \textbf{Argumentaire} : Pour chaque point, préparez 2 phrases : une idée + un exemple/une justification (\zh{因为……所以……}).
    \item \textbf{Chute} : Rédigez la conclusion avec \zh{总之} + votre avis tranché (\zh{我认为/我觉得}).
\end{enumerate}
\end{methode}

\begin{center}
\begin{tabularx}{\linewidth}{|X|X|X|X|}
\hline
\rowcolor{popblueL} \textbf{Marqueur} & \textbf{Pinyin} & \textbf{Fonction} & \textbf{Équivalent français (registre soutenu)} \\
\hline
首先 & Shǒuxiān & Premier argument / Première étape & En premier lieu, Primo, D'abord \\
\hline
然后 & Ránhòu & Deuxième argument / Suite & En second lieu, Deuxièmement, Ensuite \\
\hline
最后 & Zuìhòu & Dernier argument / Conclusion du développement & Enfin, Pour terminer, Ultimement \\
\hline
此外 & Cǐwài & Ajout d'un argument (bonus) & Par ailleurs, De plus, En outre \\
\hline
总之 & Zǒngzhī & Conclusion générale & En résumé, En conclusion, Bref \\
\hline
\end{tabularx}
\end{center}

\courssub{Exprimer son opinion : nuances entre \cle{认为}, \cle{觉得} et \cle{在……看来}}

L'erreur classique des francophones est d'utiliser uniquement \zh{我觉得} (Wǒ juéde) pour tout traduire « je pense que / je crois que / à mon avis ». En chinois, le choix du verbe signale le \textbf{degré de certitude} et le \textbf{registre de langue}.

\begin{definition}[Trois verbes d'opinion, trois registres]
\begin{itemize}
    \item \textbf{\zh{认为} (rènwéi)} : Opinion \textbf{réfléchie, rationnelle, argumentée}. Registre \emph{formel / écrit / soutenu}. « Je considère que / J'estime que / Mon avis (fondé) est que… »
    \item \textbf{\zh{觉得} (juéde)} : Impression, \textbf{ressenti subjectif}, intuition. Registre \emph{courant / oral}. « J'ai l'impression que / Je trouve que / Il me semble que… »
    \item \textbf{\zh{在……看来} (zài … kànlái)} : Structure prépositionnelle = « Du point de vue de… / Selon… / À l'avis de… ». Permet d'attribuer l'opinion à soi (\zh{在我看来}) ou à un tiers (\zh{在专家看来}, \zh{在老师看来}).
\end{itemize}
\end{definition}

\begin{exemplebox}[Comparaison d'emploi dans le débat sur la délinquance]
\begin{itemize}
    \item \zh{我认为预防犯罪应该从家庭教育抓起。} \\
    Wǒ rènwéi yùfáng fànzuì yīnggāi cóng jiātíng jiàoyù zhuā qǐ. \\
    \emph{« J'estime que la prévention de la délinquance doit partir de l'éducation familiale. »} (Avis argumenté, position de principe).
    
    \item \zh{我觉得现在的孩子压力太大了。} \\
    Wǒ juéde xiànzài de háizi yālì tài dà le. \\
    \emph{« Je trouve que les enfants d'aujourd'hui ont trop de pression. »} (Ressenti, observation empathique).
    
    \item \zh{在心理学家看来，缺爱是核心原因。} \\
    Zài xīnlǐxuéjiā kànlái, quē'ài shì héxīn yuányīn. \\
    \emph{« Du point de vue des psychologues, le manque d'amour est la cause centrale. »} (Citation d'autorité, distanciation).
    
    \item \zh{在我看来，单纯惩罚解决不了问题。} \\
    Zài wǒ kànlái, dānchún chéngfá jiějué bu liǎo wèntí. \\
    \emph{« À mon avis, la simple punition ne résout pas le problème. »} (Position personnelle affirmée, style exposé).
\end{itemize}
\end{exemplebox}

\begin{attention}[Piège « Bac » : Varier les verbes d'opinion]
\begin{itemize}
    \item \textbf{Interdit de dire :} \zh{我觉得…… 我觉得…… 我觉得……} (3 fois de suite). Cela montre une pauvreté lexicale (niveau HSK 1/2).
    \item \textbf{Attendu (niveau HSK 3/4 / Bac) :} \zh{我认为…… 此外，在专家看来…… 就我个人而言，我觉得……}
    \item \textbf{Distinction clé :} \zh{认为} s'emploie souvent avec des compléments d'objet complexes (phrases entières). \zh{觉得} s'emploie aussi pour des sensations physiques (\zh{我觉得累} « je me sens fatigué ») ou des adjectifs simples (\zh{我觉得好} « je trouve ça bien »). \zh{在……看来} ne prend \textbf{jamais} de complément d'objet direct : on ne dit pas \zh{*我在看这个问题} mais \zh{在这个问题上，我看法是……} ou \zh{在我看来……}.
\end{itemize}
\end{attention}

\courssub{Interagir : Accord, Désaccord poli et Relance}

Un débat n'est pas une suite de monologues. Les correcteurs du Bac évaluent la \cle{compétence interactionnelle} : savez-vous écouter, valider, contredire sans agresser, et relancer votre partenaire ?

\begin{propriete}[Échelle de l'accord et du désaccord (Registre scolaire respectueux)]
\textbf{A. Exprimer l'accord (valider l'interlocuteur)}
\begin{itemize}
    \item \zh{我同意你的看法。} (Wǒ tóngyì nǐ de kànfǎ.) — « Je suis d'accord avec ton avis. » (Standard, clair).
    \item \zh{你说得对。} (Nǐ shuō de duì.) — « Tu as raison / Ce que tu dis est juste. » (Très courant à l'oral, valorise l'autre).
    \item \zh{完全赞同。} (Wánquán zàntóng.) — « Tout à fait d'accord / Entièrement d'accord. » (Plus fort, engagement total).
    \item \zh{我也有同感。} (Wǒ yě yǒu tónggǎn.) — « Je partage ce sentiment. » (Accord sur un ressenti/émotion).
\end{itemize}

\textbf{B. Exprimer le désaccord poli (contredire sans heurter)}
\textit{Règle d'or :} \textbf{Marqueur de politesse} + \textbf{Marqueur de désaccord} + \textbf{Proposition alternative}.
\begin{itemize}
    \item \zh{对不起，我不同意你的看法。我觉得……} (Duìbuqǐ, wǒ bù tóngyì nǐ de kànfǎ. Wǒ juéde……) — « Excuse-moi, je ne suis pas d'accord. Je pense que… » (Formule passe-partout, sûre).
    \item \zh{你的观点有道理，但是……} (Nǐ de guāndiǎn yǒu dàoli, dànshì……) — « Ton point de vue a du sens, mais… » (Validation partielle \zh{有道理} + pivot \zh{但是}). \textbf{Très valorisé à l'examen.}
    \item \zh{恐怕不太对吧？} (Kǒngpà bù tài duì ba?) — « Je crains que ce ne soit pas tout à fait exact ? » (Désaccord atténué, interrogatif, très poli).
    \item \zh{我不太赞成这个说法。} (Wǒ bù tài zànchéng zhège shuōfǎ.) — « Je n'approuve pas trop cette façon de le dire. » (Désaccord sur la formulation, moins personnel).
\end{itemize}

\textbf{C. Relancer le débat (passer le relais, approfondir)}
\begin{itemize}
    \item \zh{你觉得呢？} (Nǐ juéde ne?) — « Et toi, qu'en penses-tu ? » (Ouvert, standard).
    \item \zh{你怎么看？} (Nǐ zěnme kàn?) — « Comment tu vois les choses ? / Quel est ton point de vue ? » (Plus large, invite à une analyse).
    \item \zh{为什么呢？} (Wèishénme ne?) — « Pourquoi ? » (Demande de justification, force l'argumentation).
    \item \zh{你能举个例子吗？} (Nǐ néng jǔ ge lìzi ma?) — « Peux-tu donner un exemple ? » (Exigence de concrétude, niveau HSK 3+).
    \item \zh{还有别的原因吗？} (Hái yǒu bié de yuányīn ma?) — « Y a-t-il d'autres causes ? » (Élargissement du champ).
\end{itemize}
\end{propriete}

\begin{exemplebox}[Échanges types : Accord / Désaccord / Relance]
\textbf{Scénario :} Élève A propose la sévérité. Élève B répond.

\textit{Accord + Relance :}
\begin{itemize}
    \item A: \zh{我觉得应该加重处罚。} (Wǒ juéde yīnggāi jiāzhòng chǔfá.) — « Je pense qu'il faut alourdir les peines. »
    \item B: \zh{我同意你的看法，法律必须有威慑力。你觉得呢？} (Wǒ tóngyì nǐ de kànfǎ, fǎlǜ bìxū yǒu wēishèlì. Nǐ juéde ne?) — « Je suis d'accord, la loi doit avoir un pouvoir de dissuasion. Et toi ? »
\end{itemize}

\textit{Désaccord poli + Argument + Relance :}
\begin{itemize}
    \item A: \zh{打工比上学有用。} (Dǎgōng bǐ shàngxué yǒuyòng.) — « Travailler est plus utile qu'étudier. »
    \item B: \zh{你的观点有道理，但是没有文凭将来很难找好工作。你能举个成功的例子吗？} (Nǐ de guāndiǎn yǒu dàoli, dànshì méiyǒu wénpíng jiānglái hěn nán zhǎo hǎo gōngzuò. Nǐ néng jǔ ge chénggōng de lìzi ma?) — « Ton avis a du sens, mais sans diplôme, il sera très dur de trouver un bon travail plus tard. Peux-tu citer un exemple de réussite ? »
\end{itemize}
\end{exemplebox}

\courssub{Organigramme du débat : la boucle interactionnelle}

\begin{popfigure}
\begin{tikzpicture}[
    node distance=1.2cm and 1.5cm,
    box/.style={draw=popblue, thick, rounded corners, fill=popblueL, text width=3.2cm, align=center, font=\small, minimum height=1.8cm},
    arrow/.style={-Stealth, thick, popdark},
    looparrow/.style={-Stealth, thick, poporange, dashed}
]
% Nœuds
\node[box, align=center] (avis){\textbf{Donner son avis}\\\zh{我认为 / 觉得 / 在我看来}};
\node[box, right=of avis, align=center] (react){\textbf{Réagir}\\\zh{同意 / 不同意 / 有道理但…}};
\node[box, below=of react, align=center] (relance){\textbf{Relancer}\\\zh{你觉得呢？ / 为什么？ / 举个例子？}};
\node[box, left=of relance, align=center] (concl){\textbf{Conclure / Synthétiser}\\\zh{总之 / 综上所述}};

% Flèches principales
\draw[arrow] (avis) -- node[above, font=\tiny, popdark] {Écoute} (react);
\draw[arrow] (react) -- node[right, font=\tiny, popdark] {Approfondir} (relance);
\draw[arrow] (relance) -- node[below, font=\tiny, popdark] {Nouvel avis} (concl);
\draw[arrow] (concl) -- node[left, font=\tiny, popdark] {Clôture / Nouveau tour} (avis);

% Boucle de relance vers avis (nouveau tour de parole)
\draw[looparrow] (relance) to[out=30, in=150, looseness=1.5] node[above, font=\tiny, poporange] {Tour suivant} (avis);

% Flèche désaccord -> nouvel avis direct
\draw[arrow, popgreen] (react) to[out=-120, in=-60, looseness=1.2] node[below, font=\tiny, popgreen] {Contre-argument} (avis);
\end{tikzpicture}
\legende{Organigramme de l'interaction orale : la boucle Avis $\rightarrow$ Réaction $\rightarrow$ Relance $\rightarrow$ Conclusion (et retour). La flèche verte montre le contre-argument direct ; la flèche orange en pointillés montre le passage de témoin au partenaire.}
\end{popfigure}

\courssub{Exercice résolu : Transformer un monologue en débat interactif}

\begin{exoresolu}[De l'exposé au dialogue : restructuration]
\textbf{Consigne :} Voici un paragraphe unique d'un élève. Transformez-le en un échange de 4 répliques entre \textbf{Élève A} (qui présente) et \textbf{Élève B} (qui réagit, relance, conclut). Utilisez les formules de cette leçon.

\textit{Paragraphe source :} « Je pense que la délinquance vient de la famille. D'abord les parents ne surveillent pas. Ensuite les parents donnent de l'argent mais pas d'amour. Donc il faut punir les parents. En résumé, la famille est la clé. »

\textbf{Solution détaillée :}

\begin{center}
\begin{tabularx}{\linewidth}{|X|X|}
\hline
\rowcolor{popgreenL} \textbf{Élève A (Exposé structuré)} & \textbf{Élève B (Interaction)} \\
\hline
\zh{大家好，今天我要谈的问题是青少年犯罪。首先，我认为家庭是核心原因。} & \\
\hline
& \zh{你觉得呢？能具体说说吗？} \textit{(Relance : ouverture + demande de précision)} \\
\hline
\zh{然后，很多父母只给钱，不给爱，缺乏陪伴。} & \zh{你的观点有道理，但是只惩罚父母公平吗？} \textit{(Accord partiel + Désaccord poli + Pivot)} \\
\hline
\zh{最后，我觉得应该加强家庭教育指导，而不只是惩罚。总之，预防犯罪要从源头抓起。} & \zh{我同意你的看法，教育比惩罚更重要。} \textit{(Accord final + Reformulation de la conclusion)} \\
\hline
\end{tabularx}
\end{center}

\textbf{Analyse pédagogique :}
\begin{enumerate}
    \item A utilise \zh{首先/然后/最后} pour structurer (critère « Cohérence »).
    \item A varie \zh{我认为} (avis principal) et \zh{我觉得} (impression sur la solution).
    \item B utilise \zh{你觉得呢？} pour passer la main (critère « Interaction »).
    \item B utilise \zh{你的观点有道理，但是…} pour désaccorder poliment (critère « Registre / Sociolinguistique »).
    \item B conclut par \zh{我同意…} validant la synthèse de A.
\end{enumerate}
\end{exoresolu}

\courssub{Entraînement guidé : « Le tribunal de la classe »}

\begin{exercice}[Mise en situation : Débat « Punir ou Éduquer ? »]
\textbf{Contexte :} Le proviseur organise une table ronde d'élèves sur : \zh{「青少年犯罪：该重惩罚还是重教育？」} (La délinquance juvénile : faut-il privilégier la punition ou l'éducation ?).

\textbf{Rôles :}
\begin{itemize}
    \item \textbf{Groupe 1 (Partisans « Éducation ») :} Préparez un exposé de 4 phrases (Salut+Sujet, \zh{首先}, \zh{然后}, \zh{最后}) + Conclusion \zh{总之}. Utilisez \zh{我认为}, \zh{在我们看来}.
    \item \textbf{Groupe 2 (Partisans « Punition ») :} Même structure. Utilisez \zh{我觉得}, \zh{我们认为}.
    \item \textbf{Groupe 3 (Journalistes / Modérateurs) :} Préparez 5 questions de relance variées (\zh{为什么？}, \zh{举个例子？}, \zh{怎么看？}, \zh{有证据吗？}, \zh{如果……怎么办？}).
\end{itemize}

\textbf{Déroulement :}
1. Un élève du Groupe 1 présente (1 min).
2. Un élève du Groupe 2 réagit avec \zh{有道理，但是…} + contre-argument (30 s).
3. Un Journaliste relance l'un ou l'autre (\zh{你能举个例子吗？}).
4. Changement de rôles. Le professeur note l'usage des marqueurs (\zh{首先/然后/最后}, \zh{我认为/觉得}, \zh{有道理/不同意}, \zh{总之}).
\end{exercice}

\begin{savaistu}[Culture de l'oral en Chine : \zh{演讲比赛} (Yǎnjiǎng bǐsài)]
En Chine, les \textbf{concours d'éloquence} (\zh{演讲比赛}) sont une institution scolaire majeure, du primaire à l'université. Ils sont souvent diffusés à la télévision (ex: \zh{《超级演说家》} « Super Speaker »).
\begin{itemize}
    \item \textbf{Structure imposée :} Les candidats \emph{doivent} utiliser \zh{首先、其次、最后} (ou \zh{第一、第二、第三}) pour structurer leur discours de 3 à 5 minutes. L'absence de ces marqueurs est pénalisée comme un manque de logique (\zh{逻辑不清}).
    \item \textbf{Rhétorique :} L'usage de \zh{在……看来} pour citer des autorités (Confucius, leaders, scientifiques) et de \zh{总之} / \zh{综上所述} pour conclure est la norme du « bon chinois standard » (\zh{标准普通话}).
    \item \textbf{Parallèle Cameroun :} Au Cameroun, les concours de \emph{discours} (Concours d'éloquence du MINESEC, Jeux de la Francophonie) valorisent aussi la structure \emph{Introduction / Développement en 3 parties / Conclusion}. La différence culturelle majeure : en Chine, l'émotion (\zh{感人}) est souvent recherchée dans la conclusion (récit personnel, appel aux valeurs collectives), alors qu'en contexte francophone camerounais, l'argumentation rationnelle et la maîtrise de la langue française priment souvent. Maîtriser les deux codes est un atout pour nos bacheliers !
\end{itemize}
\end{savaistu}

\begin{aretenir}[Check-list « Bac Oral » — Formules fonctionnelles (Section 3)]
\begin{itemize}
    \item \textbf{Structure :} \zh{大家好，今天我要谈的问题是…} $\rightarrow$ \zh{首先… 然后… 最后…} $\rightarrow$ \zh{总之，我认为…}
    \item \textbf{Opinion :} \zh{认为} (réfléchi/formel) \vs \zh{觉得} (ressenti/oral) \vs \zh{在…看来} (attribution).
    \item \textbf{Accord :} \zh{我同意…} / \zh{你说得对。}
    \item \textbf{Désaccord poli :} \zh{你的观点有道理，但是…} (Le \textbf{Saint Graal} de l'interaction !).
    \item \textbf{Relance :} \zh{你觉得呢？} / \zh{你能举个例子吗？} / \zh{为什么呢？}
    \item \textbf{Tons :} Attention aux tons 3 (\zh{认为 rènwéi}, \zh{觉得 juéde}, \zh{看法 kànfǎ}) et ton neutre (\zh{呢 ne}, \zh{吗 ma}).
\end{itemize}
\end{aretenir}

\coursec{Dialogue modèle 2 : le débat — punir ou éduquer ?}

Dans la section précédente, nous avons appris les \cle{formules fonctionnelles} qui permettent de présenter un sujet (\zh{今天我谈的是…}), de donner son avis (\zh{我认为…}, \zh{在我看来…}) et de relancer la parole (\zh{你怎么看？}). Il est temps de les mettre en action dans une situation authentique : un \cle{débat} entre deux lycéens. Le sujet est classique et très riche : face à la \cle{délinquance juvénile}, faut-il \cle{punir} ou faut-il \cle{éduquer} ?

\courssub{Observation : le dialogue du débat}

Lisez d'abord le dialogue en silence, repérez les formules déjà connues, puis écoutez la lecture du professeur en marquant les tons au crayon.

\begin{textebox}[Dialogue 2 : 惩罚还是教育？]
A：我认为年轻的罪犯应该受到严厉的惩罚。\\
\emph{Wǒ rènwéi niánqīng de zuìfàn yīnggāi shòudào yánlì de chéngfá.}\\
« Je pense que les jeunes délinquants devraient recevoir une punition sévère. »\\[4pt]
B：对不起，我不同意。在我看来，他们还年轻，教育比惩罚更重要。\\
\emph{Duìbuqǐ, wǒ bù tóngyì. Zài wǒ kànlái, tāmen hái niánqīng, jiàoyù bǐ chéngfá gèng zhòngyào.}\\
« Pardon, je ne suis pas d'accord. À mon avis, ils sont encore jeunes : l'éducation est plus importante que la punition. »\\[4pt]
A：可是如果没有惩罚，他们会再犯罪。\\
\emph{Kěshì rúguǒ méiyǒu chéngfá, tāmen huì zài fànzuì.}\\
« Mais s'il n'y a pas de punition, ils recommenceront à commettre des délits. »\\[4pt]
B：你说得有道理，但是监狱不能解决所有问题。学校和家庭应该帮助他们。\\
\emph{Nǐ shuō de yǒu dàolǐ, dànshì jiānyù bù néng jiějué suǒyǒu de wèntí. Xuéxiào hé jiātíng yīnggāi bāngzhù tāmen.}\\
« Ce que tu dis est sensé, mais la prison ne peut pas résoudre tous les problèmes. L'école et la famille devraient les aider. »\\[4pt]
A：嗯，也许你说得对。惩罚和教育都重要。\\
\emph{Ng, yěxǔ nǐ shuō de duì. Chéngfá hé jiàoyù dōu zhòngyào.}\\
« Hum, peut-être que tu as raison. La punition et l'éducation sont toutes les deux importantes. »
\end{textebox}

\textbf{Glossaire du dialogue et expressions utiles}

\begin{center}
\begin{tabularx}{\linewidth}{|X|X|X|X|}
\hline
\rowcolor{popblueL} Caractères & Pinyin & Nature & Traduction \\
\hline
罪犯 & zuìfàn & nom & criminel, délinquant \\
\hline
受到 & shòudào & verbe & recevoir, subir \\
\hline
严厉 & yánlì & adjectif & sévère, strict \\
\hline
惩罚 & chéngfá & nom / verbe & punition ; punir \\
\hline
同意 & tóngyì & verbe & être d'accord \\
\hline
教育 & jiàoyù & nom / verbe & éducation ; éduquer \\
\hline
犯罪 & fànzuì & verbe & commettre un crime \\
\hline
道理 & dàolǐ & nom & raison, bon sens \\
\hline
监狱 & jiānyù & nom & prison \\
\hline
解决 & jiějué & verbe & résoudre \\
\hline
责任 & zérèn & nom & responsabilité \\
\hline
也许 & yěxǔ & adverbe & peut-être \\
\hline
\end{tabularx}
\end{center}

\textbf{Questions de compréhension}
\begin{enumerate}
\item Que pense A au début du débat ? \emph{(Que pense A ?)}
\item Pourquoi B préfère-t-il l'éducation ?
\item Quel argument A utilise-t-il pour défendre la punition ?
\item Comment le débat se termine-t-il ? Qui « gagne » ?
\end{enumerate}

\courssub{Point de langue 1 : le comparatif avec 比}

\textbf{Observation.} Dans le dialogue, B dit : \zh{教育比惩罚更重要。} \emph{Jiàoyù bǐ chéngfá gèng zhòngyào.} « L'éducation est plus importante que la punition. » Le mot \zh{比} \emph{bǐ} correspond au français « que » dans « plus important \emph{que} ». Mais attention à l'ordre : en chinois, on dit littéralement « éducation — comparé-à — punition — plus importante ».

\begin{propriete}[Le comparatif avec 比]
Structure : \cle{A + 比 + B + adjectif}.\\
\zh{教育比惩罚重要。} \emph{Jiàoyù bǐ chéngfá zhòngyào.} « L'éducation est plus importante que la punition. »
\begin{itemize}
\item Pour renforcer : ajouter \zh{更} \emph{gèng} « encore plus » ou \zh{得多} \emph{de duō} « beaucoup plus » après l'adjectif : \zh{教育比惩罚重要得多。} « L'éducation est \emph{beaucoup} plus importante. »
\item Pour une petite différence : \zh{一点儿} \emph{yìdiǎnr} après l'adjectif : \zh{教育比惩罚重要一点儿。} « …un peu plus importante. »
\item \textbf{Négation} : on n'utilise pas \zh{不比} pour « moins… que », mais \zh{没有} : \cle{A + 没有 + B + adjectif}. \zh{惩罚没有教育重要。} \emph{Chéngfá méiyǒu jiàoyù zhòngyào.} « La punition n'est pas aussi importante que l'éducation. »
\item Jamais \zh{很} directement après \zh{比} : on dit \zh{更}, pas \zh{很重要}.
\end{itemize}
\end{propriete}

\begin{popfigure}
\begin{tikzpicture}[node distance=0.5cm and 0.9cm]
\node[draw=popblue, fill=popblueL, rounded corners, minimum width=2.2cm, minimum height=0.9cm, align=center] (A) {A\\ \zh{教育}};
\node[draw=poporange, fill=poporangeL, rounded corners, minimum width=1.2cm, minimum height=0.9cm, align=center, right=of A] (bi) {\zh{比}\\ bǐ};
\node[draw=popblue, fill=popblueL, rounded corners, minimum width=2.2cm, minimum height=0.9cm, align=center, right=of bi] (B) {B\\ \zh{惩罚}};
\node[draw=popgreen, fill=popgreenL, rounded corners, minimum width=2.6cm, minimum height=0.9cm, align=center, right=of B] (adj) {adjectif\\ (\zh{更})\zh{重要}};
\node[draw=poppurple, fill=poppinkL, rounded corners, minimum width=2.2cm, minimum height=0.9cm, align=center, right=of adj] (deg) {degré\\ \zh{得多}/\zh{一点儿}};
\draw[-{Stealth}, thick, popdark] (A) -- (bi);
\draw[-{Stealth}, thick, popdark] (bi) -- (B);
\draw[-{Stealth}, thick, popdark] (B) -- (adj);
\draw[-{Stealth}, thick, popdark, dashed] (adj) -- (deg);
\node[below=0.35cm of bi, align=center, text=popmuted] {\zh{教育比惩罚更重要。} « L'éducation est plus importante que la punition. »};
\end{tikzpicture}
\legende{Schéma du comparatif : A + 比 + B + adjectif (+ degré optionnel)}
\end{popfigure}

\begin{center}
\begin{tabularx}{\linewidth}{|X|X|X|}
\hline
\rowcolor{popblueL} Structure & Exemple & Traduction \\
\hline
A + 比 + B + adj. & \zh{教育比惩罚重要。} Jiàoyù bǐ chéngfá zhòngyào. & L'éducation est plus importante que la punition. \\
\hline
A + 比 + B + 更 + adj. & \zh{学校比监狱更有用。} Xuéxiào bǐ jiānyù gèng yǒuyòng. & L'école est encore plus utile que la prison. \\
\hline
A + 比 + B + adj. + 得多 & \zh{雅温得比我的小村大得多。} Yǎwēndé bǐ wǒ de xiǎo cūn dà de duō. & Yaoundé est beaucoup plus grande que mon village. \\
\hline
A + 没有 + B + adj. & \zh{惩罚没有教育重要。} Chéngfá méiyǒu jiàoyù zhòngyào. & La punition est moins importante que l'éducation. \\
\hline
\end{tabularx}
\end{center}

\courssub{Point de langue 2 : l'hypothèse avec 如果…就/会…}

\textbf{Observation.} A objecte : \zh{如果没有惩罚，他们会再犯罪。} \emph{Rúguǒ méiyǒu chéngfá, tāmen huì zài fànzuì.} « S'il n'y a pas de punition, ils récidiveront. » Le mot \zh{如果} \emph{rúguǒ} = « si », et \zh{会} \emph{huì} marque la conséquence probable.

\begin{propriete}[L'hypothèse avec 如果…(就/会)…]
Structure : \cle{如果 + condition，sujet + 就/会 + conséquence}.\\
\zh{如果你帮助他们，他们就会改变。} \emph{Rúguǒ nǐ bāngzhù tāmen, tāmen jiù huì gǎibiàn.} « Si tu les aides, ils changeront. »
\begin{itemize}
\item \zh{就} \emph{jiù} insiste sur la conséquence logique directe ; \zh{会} \emph{huì} marque la probabilité future. Les deux peuvent se combiner : \zh{就会}.
\item \zh{如果} peut se placer avant ou après le sujet de la condition : \zh{如果没有惩罚…} = \zh{惩罚如果没有…} (la première forme est la plus courante).
\item En français, « si » + présent ; en chinois, pas de temps : le contexte suffit.
\end{itemize}
\end{propriete}

\begin{exemplebox}[Exemples d'hypothèses pour le débat]
\zh{如果学校不关心学生，问题会更多。} \emph{Rúguǒ xuéxiào bù guānxīn xuésheng, wèntí huì gèng duō.} « Si l'école ne se préoccupe pas des élèves, il y aura plus de problèmes. »\\
\zh{如果家庭也有责任，父母就应该多陪孩子。} \emph{Rúguǒ jiātíng yě yǒu zérèn, fùmǔ jiù yīnggāi duō péi háizi.} « Si la famille a aussi une responsabilité, les parents devraient passer plus de temps avec leurs enfants. »\\
\zh{如果只惩罚不教育，年轻人不会进步。} \emph{Rúguǒ zhǐ chéngfá bù jiàoyù, niánqīngrén bú huì jìnbù.} « Si l'on punit sans éduquer, les jeunes ne progresseront pas. »
\end{exemplebox}

\courssub{Point de langue 3 : l'opposition avec 可是 / 但是}

\zh{可是} \emph{kěshì} et \zh{但是} \emph{dànshì} signifient tous deux « mais ». \zh{但是} est un peu plus formel, \zh{可是} plus oral. Ils introduisent la réfutation : \zh{可是如果没有惩罚…} « Mais s'il n'y a pas de punition… ». On peut les combiner avec \zh{虽然} « bien que » : \zh{虽然你说得有道理，但是我不同意。} « Bien que tu aies raison, je ne suis pas d'accord. »

\begin{methode}[Nuancer au lieu de contredire brutalement]
Dans un débat noté, évitez de répondre seulement \zh{我不同意！} « Je ne suis pas d'accord ! », trop sec. Procédez en trois étapes :
\begin{enumerate}
\item \textbf{Validez} d'abord l'argument adverse : \zh{你说得有道理。} \emph{Nǐ shuō de yǒu dàolǐ.} « Ce que tu dis est sensé. »
\item \textbf{Introduisez} l'opposition : \zh{但是…} / \zh{可是…} « mais… ».
\item \textbf{Donnez} votre contre-argument avec une structure grammaticale forte (comparatif ou hypothèse) : \zh{但是监狱不能解决所有问题。} « Mais la prison ne peut pas résoudre tous les problèmes. »
\end{enumerate}
Cette stratégie de \cle{concession} est très valorisée à l'oral : elle montre que vous écoutez avant de répondre.
\end{methode}

\begin{exemplebox}[Phrases clés à réutiliser]
\zh{你说得有道理。} \emph{Nǐ shuō de yǒu dàolǐ.} « Ce que tu dis est sensé. » (structure : verbe + \zh{得} + adjectif, « parler de façon sensée »)\\
\zh{监狱不能解决所有问题。} \emph{Jiānyù bù néng jiějué suǒyǒu de wèntí.} « La prison ne peut pas résoudre tous les problèmes. »\\
\zh{学校和家庭应该帮助他们。} \emph{Xuéxiào hé jiātíng yīnggāi bāngzhù tāmen.} « L'école et la famille devraient les aider. »
\end{exemplebox}

\begin{attention}[La place de 也 et 都]
Les francophones placent souvent mal \zh{也} \emph{yě} « aussi » et \zh{都} \emph{dōu} « tous/toutes les deux ». Règle absolue : ils se placent \textbf{avant le verbe}, jamais en fin de phrase comme « aussi » en français.\\
Correct : \zh{家庭也有责任。} \emph{Jiātíng yě yǒu zérèn.} « La famille aussi a une responsabilité. »\\
Correct : \zh{惩罚和教育都重要。} \emph{Chéngfá hé jiàoyù dōu zhòngyào.} « La punition et l'éducation sont toutes les deux importantes. »\\
Faux : \zh{家庭有责任也。} et \zh{惩罚和教育重要都。}
\end{attention}

\courssub{Travail des tons du dialogue}

Quatre mots-clés à entraîner à voix haute, en exagérant les tons :
\begin{itemize}
\item \zh{惩罚} \emph{chéngfá} : 2\textsuperscript{e} + 2\textsuperscript{e} ton — deux montées successives, comme deux questions qui s'enchaînent.
\item \zh{严厉} \emph{yánlì} : 2\textsuperscript{e} + 4\textsuperscript{e} ton — monter puis tomber sèchement, ton « sévère » adapté au sens !
\item \zh{解决} \emph{jiějué} : 3\textsuperscript{e} + 2\textsuperscript{e} ton — descendre-remonter, puis monter.
\item \zh{道理} \emph{dàolǐ} : 4\textsuperscript{e} + 3\textsuperscript{e} ton — tomber, puis descendre-remonter.
\end{itemize}

\begin{exoresolu}[Construisez votre argument]
Énoncé : vous êtes « pour l'éducation ». Complétez en chinois : (1) un comparatif avec \zh{比} pour dire que l'aide est plus utile que la prison ; (2) une hypothèse avec \zh{如果…就…} pour dire que si on les aide, ils ne recommenceront pas.
\tcblower
Solution détaillée :\\
(1) \zh{帮助比监狱更有用。} \emph{Bāngzhù bǐ jiānyù gèng yǒuyòng.} « L'aide est plus utile que la prison. » Structure : A (\zh{帮助}) + \zh{比} + B (\zh{监狱}) + \zh{更} + adjectif (\zh{有用}).\\
(2) \zh{如果我们帮助他们，他们就不会再犯罪。} \emph{Rúguǒ wǒmen bāngzhù tāmen, tāmen jiù bú huì zài fànzuì.} « Si nous les aidons, ils ne récidiveront pas. » Structure : \zh{如果} + condition，sujet + \zh{就} + \zh{不会} + \zh{再} + verbe. Notez \zh{再} « de nouveau » placé avant le verbe, et la mutation de ton : \zh{不} devient 2\textsuperscript{e} ton (\emph{bú}) devant \zh{会} (4\textsuperscript{e} ton).
\end{exoresolu}

\begin{exercice}[À vous de débattre]
Par binômes, rejouez le dialogue en remplaçant deux arguments : A doit utiliser une hypothèse avec \zh{如果…会…}, B doit utiliser un comparatif avec \zh{比…更…} et terminer par la formule de concession \zh{你说得有道理，但是…}. Le professeur écoutera surtout les tons de \zh{惩罚}, \zh{严厉}, \zh{解决} et \zh{道理}.
\end{exercice}

\begin{aretenir}[L'essentiel de la section]
\begin{itemize}
\item Comparatif : \cle{A + 比 + B + adjectif} ; négation avec \zh{没有} ; renforcement avec \zh{更} ou \zh{得多}.
\item Hypothèse : \cle{如果…，…就/会…} pour annoncer une conséquence.
\item Opposition : \zh{可是} / \zh{但是} ; stratégie de concession : \zh{你说得有道理，但是…}.
\item \zh{也} et \zh{都} se placent toujours avant le verbe.
\end{itemize}
\end{aretenir}

\coursec{Phonétique : les tons à travailler}

Après avoir mis en œuvre le dialogue modèle 2 et exploré les arguments du débat « punir ou éduquer ? », il est indispensable de fixer la prononciation du lexique spécifique à ce thème. En chinois, un ton mal prononcé change le sens du mot — et dans un exposé noté, cela peut nuire à la compréhension comme à la note. Cette section révise les quatre tons et le ton neutre sur le vocabulaire de la leçon, traite les paires minimales pièges, explicite la règle de sandhi du troisième ton, travaille l'intonation des questions, et propose des exercices de discrimination auditive ainsi qu'un chant choral des formules fonctionnelles avec gestuelle.

\courssub{Révision des quatre tons et du ton neutre sur le vocabulaire de la leçon}

\begin{textebox}[Vocabulaire thématique — tons et sens]
青少年 qīngshàonián \quad 犯罪 fànzuì \quad 原因 yuányīn \quad 影响 yǐngxiǎng \quad 教育 jiàoyù \quad 重要 zhòngyào \quad 觉得 juéde \quad 问题 wèntí \quad 社会 shèhuì \quad 家庭 jiātíng \quad 学校 xuéxiào \quad 法律 fǎlǜ \quad 偷窃 tōuqiè \quad 打架 dǎjià \quad 吸毒 xīdú \quad 网瘾 wǎngyǐn \quad 监护人 jiānhùrén \quad 心理 xīnlǐ \quad 辅导 fǔdǎo \quad 预防 yùfáng
\end{textebox}

Observez la répartition des tons dans ce lexique (se reporter au tableau ci-dessous pour le ton exact de chaque caractère) :
\begin{itemize}
\item \textbf{1er ton (plat haut 55)} : \zh{青} qīng, \zh{偷} tōu, \zh{家} jiā, \zh{心} xīn, \zh{监} jiān, \zh{吸} xī.
\item \textbf{2e ton (montant 35)} : \zh{年} nián, \zh{原} yuán, \zh{人} rén, \zh{防} fáng.
\item \textbf{3e ton (descendant-montant 214)} : \zh{少} shǎo, \zh{影} yǐng, \zh{打} dǎ, \zh{网} wǎng, \zh{理} lǐ, \zh{辅} fǔ, \zh{法} fǎ, \zh{瘾} yǐn, \zh{导} dǎo.
\item \textbf{4e ton (descendant franc 51)} : \zh{犯} fàn, \zh{罪} zuì, \zh{响} xiǎng, \zh{育} yù, \zh{要} yào, \zh{题} tí, \zh{会} huì, \zh{社} shè, \zh{校} xiào, \zh{问} wèn, \zh{窃} qiè, \zh{架} jià, \zh{毒} dú, \zh{护} hù, \zh{预} yù, \zh{教} jiào, \zh{重} zhòng.
\item \textbf{Ton neutre} : \zh{的} de (dans \zh{觉得} juéde), \zh{人} ren (dans \zh{监护人} jiānhùrén), et les particules \zh{了} le, \zh{吗} ma, \zh{的} de.
\end{itemize}

\begin{center}
\begin{tabularx}{\linewidth}{|X|X|X|X|}
\hline
\rowcolor{popblueL} \textbf{Caractères} & \textbf{Pinyin} & \textbf{Nature} & \textbf{Traduction} \\ \hline
青少年 & qīngshàonián & nom & adolescent(s) \\ \hline
犯罪 & fànzuì & verbe / nom & commettre un crime / criminalité \\ \hline
原因 & yuányīn & nom & cause, raison \\ \hline
影响 & yǐngxiǎng & verbe / nom & influencer / influence \\ \hline
教育 & jiàoyù & verbe / nom & éduquer / éducation \\ \hline
重要 & zhòngyào & adjectif & important \\ \hline
觉得 & juéde & verbe & penser, trouver que \\ \hline
问题 & wèntí & nom & problème, question \\ \hline
社会 & shèhuì & nom & société \\ \hline
家庭 & jiātíng & nom & famille, foyer \\ \hline
学校 & xuéxiào & nom & école \\ \hline
法律 & fǎlǜ & nom & loi, législation \\ \hline
偷窃 & tōuqiè & verbe & voler (à la tire) \\ \hline
打架 & dǎjià & verbe & se battre \\ \hline
吸毒 & xīdú & verbe & prendre de la drogue \\ \hline
网瘾 & wǎngyǐn & nom & cyberdépendance \\ \hline
监护人 & jiānhùrén & nom & tuteur légal \\ \hline
心理 & xīnlǐ & nom & psychologie, mental \\ \hline
辅导 & fǔdǎo & verbe / nom & encadrer, tutorer / soutien scolaire \\ \hline
预防 & yùfáng & verbe & prévenir \\ \hline
\end{tabularx}
\end{center}

\begin{popfigure}
\begin{tikzpicture}[scale=1.1, every node/.style={font=\small}]
% Axes
\draw[popline, thick, ->] (-0.5,0) -- (5.5,0) node[right] {Temps};
\draw[popline, thick, ->] (0,-0.5) -- (0,4.5) node[above] {Hauteur};

% Ligne de référence milieu
\draw[popmuted, dashed] (0,2) -- (5,2);

% 1er ton : plat haut (55)
\draw[popblue, very thick] (0.2,3.8) -- (1.3,3.8);
\node[popblue, align=center] at (0.75,4.2) {1er ton \\ \textbf{55} \\ \zh{青} qīng};

% 2e ton : montant (35)
\draw[poporange, very thick] (1.7,2) .. controls (2.2,3) and (2.8,3.5) .. (3.2,3.8);
\node[poporange, align=center] at (2.45,4.2) {2e ton \\ \textbf{35} \\ \zh{年} nián};

% 3e ton : descendant-montant (214)
\draw[popgreen, very thick] (3.6,2) .. controls (4.0,1) and (4.4,1) .. (4.8,2.5);
\node[popgreen, align=center] at (4.2,0.4) {3e ton \\ \textbf{214} \\ \zh{影} yǐng};

% 4e ton : descendant franc (51)
\draw[poppurple, very thick] (5.2,3.8) -- (5.8,0.5);
\node[poppurple, align=center] at (5.5,4.2) {4e ton \\ \textbf{51} \\ \zh{犯} fàn};

% Marques de repère
\foreach \y/\label in {3.8/5, 3/4, 2/3, 1/2, 0.5/1} {
    \draw[popmuted] (-0.2,\y) -- (0,\y);
    \node[left, popmuted] at (-0.2,\y) {\tiny \label};
}
\end{tikzpicture}
\legende{Schéma des quatre tons du mandarin standard (système de Chao) avec un mot-exemple du vocabulaire de la leçon par ton. Le 1er ton reste haut (55), le 2e monte (35), le 3e descend puis remonte (214), le 4e chute franchement (51).}
\end{popfigure}

\begin{methode}[Geste de la main pour ancrer chaque ton]
Associez un mouvement de la main droite à chaque ton pendant la répétition chorale :
\begin{enumerate}
\item \textbf{1er ton (plat haut)} : main horizontale, paume vers le bas, niveau épaule — ne bouge pas.
\item \textbf{2e ton (montant)} : main partie du bas (niveau ventre) vers le haut (niveau visage), mouvement fluide.
\item \textbf{3e ton (creux)} : main descend bas (niveau hanche) en s'arrondissant, puis remonte légèrement.
\item \textbf{4e ton (chute)} : main partie du haut (niveau tête) vers le bas (niveau taille), coup sec, comme un couperet.
\item \textbf{Ton neutre} : petit coup de poignet léger, sans amplitude, main relâchée.
\end{enumerate}
Répétez chaque mot du tableau en synchronisant le geste. Cela ancre la mélodie dans la mémoire corporelle.
\end{methode}

\courssub{Paires minimales à risque}

Certains mots du thème forment des paires minimales où seul le ton distingue le sens. Les francophones, dont la langue n'est pas tonale, les confondent facilement.

\begin{exemplebox}[Paires minimales — fǎ / fà, tōu / tóu, jiào / jiāo / jiào]
\begin{itemize}
\item \zh{法} fǎ (3e ton) = loi, méthode \quad vs \quad \zh{发} fà (4e ton) = cheveux, émettre.
  \begin{itemize}
  \item \zh{法律} fǎlǜ (loi) — \zh{头发} tóufà (cheveux).
  \item Piège : dire \zh{*fàlǜ} pour « loi » → incompréhensible.
  \end{itemize}
\item \zh{偷} tōu (1er ton) = voler \quad vs \quad \zh{头} tóu (2e ton) = tête.
  \begin{itemize}
  \item \zh{偷窃} tōuqiè (voler) — \zh{头痛} tóutòng (mal de tête).
  \item Piège : \zh{*tóuqiè} sonne comme « tête-couper » au lieu de « vol ».
  \end{itemize}
\item \zh{教} jiào (4e ton) dans \zh{教育} jiàoyù = éduquer \quad vs \quad \zh{交} jiāo (1er ton) = remettre, croiser \quad vs \quad \zh{叫} jiào (4e ton) = appeler, crier.
  \begin{itemize}
  \item \zh{教育} jiàoyù (éducation) — \zh{交流} jiāoliú (échanger) — \zh{叫醒} jiàoxǐng (réveiller).
  \item Dans le débat, on dit \zh{加强教育} jiāqiáng jiàoyù « renforcer l'éducation » : attention à ne pas aplatir le 4e ton de \zh{教}.
  \end{itemize}
\end{itemize}
\end{exemplebox}

\begin{attention}[Erreur fréquente des francophones — aplatir le 4e ton]
Le 4e ton (\zh{犯} fàn, \zh{罪} zuì, \zh{教} jiào, \zh{育} yù, \zh{要} yào, \zh{打} dǎ, \zh{吸} xī, \zh{毒} dú, \zh{防} fáng, \zh{预} yù, \zh{社} shè, \zh{校} xiào, \zh{窃} qiè, \zh{架} jià, \zh{护} hù) doit tomber \textbf{franchement}, comme un ordre sec : « Arrête ! » En français, l'intonation descend à la fin de phrase, mais pas sur une seule syllabe avec cette amplitude. Si vous dites \zh{fànzuì} avec un 4e ton mou (mi-descendu), cela ressemble à un 3e ton mal fait et l'auditeur entendra \zh{fǎnzuǐ} « bouche retournée » (absurde). \textbf{Exercice} : placez la main haut, laissez-la tomber d'un coup sec en prononçant \zh{犯} fàn — \zh{罪} zuì — \zh{教} jiào — \zh{育} yù — \zh{预} yù — \zh{防} fáng. Répétez 10 fois.
\end{attention}

\courssub{Séries à répéter et règle de sandhi du 3e ton}

Le vocabulaire de la leçon contient plusieurs mots à deux syllabes où le troisième ton subit une transformation obligatoire (sandhi).

\begin{propriete}[Sandhi du troisième ton — deux 3e tons consécutifs]
Lorsque deux syllabes de 3e ton se suivent, \textbf{la première devient un 2e ton (montant)} à l'oral. L'écriture en pinyin garde le 3e ton (notation lexicale), mais la prononciation réelle change.
\begin{center}
\begin{tabularx}{\linewidth}{|X|X|X|X|}
\hline
\rowcolor{popblueL} \textbf{Mot (écrit)} & \textbf{Pinyin lexical} & \textbf{Pinyin réel (sandhi)} & \textbf{Traduction} \\ \hline
影响 & yǐngxiǎng & \textbf{yíngxiǎng} (2+3) & influence / influencer \\ \hline
你好 & nǐhǎo & \textbf{níhǎo} (2+3) & bonjour \\ \hline
很好 & hěnhǎo & \textbf{hénhǎo} (2+3) & très bien \\ \hline
辅导 & fǔdǎo & \textbf{fúdǎo} (2+3) & tutorat, encadrement \\ \hline
\end{tabularx}
\end{center}
\textbf{Règle} : 3e + 3e $\rightarrow$ 2e + 3e. Cela vaut \textbf{uniquement à l'oral}. En dictionnaire et en écriture pinyin standard, on note toujours le ton lexical (3e). Attention : 3e + 4e (ex. \zh{打架} dǎjià, \zh{想要} xiǎngyào) $\rightarrow$ pas de sandhi (la 1re syllabe reste 3e, souvent réalisée en « demi-3e ton » bas).
\end{propriete}

\begin{exemplebox}[Séries à répéter — application du sandhi]
Répétez chaque série en appliquant le sandhi là où il faut. Le geste de la main (méthode ci-dessus) aide à sentir le passage du creux à la montée.

\begin{enumerate}
\item \zh{青少年} qīngshàonián (1-4-2) — pas de sandhi (shào 4e, nián 2e).
\item \zh{犯罪} fànzuì (4-4) — deux 4e tons : bien faire chuter chaque syllabe.
\item \zh{原因} yuányīn (2-1) — montant puis plat haut.
\item \zh{影响} yǐngxiǎng $\rightarrow$ \textbf{yíngxiǎng} (2-3) — \emph{sandhi actif}.
\item \zh{教育} jiàoyù (4-4) — deux chutes franches.
\item \zh{重要} zhòngyào (4-4) — idem.
\item \zh{觉得} juéde (2 + neutre) — 2e ton monté, puis ton neutre léger.
\item \zh{监护人} jiānhùrén (1-4-neutre) — attention à \zh{护} hù (4e ton).
\item \zh{心理辅导} xīnlǐ fǔdǎo $\rightarrow$ \textbf{xīnlǐ fúdǎo} (1-3, 2-3) — \zh{心理} xīnlǐ (1+3, pas de sandhi) ; \zh{辅导} fǔdǎo (3+3 $\rightarrow$ 2+3) $\rightarrow$ fúdǎo.
\item \zh{预防} yùfáng (4-2) — chute puis montée.
\end{enumerate}
\end{exemplebox}

\begin{exoresolu}[Discrimination de sandhi]
\textbf{Énoncé :} L'enseignant lit les mots suivants. L'élève note le pinyin \textbf{réel (oral)} avec le sandhi appliqué si nécessaire.
\begin{enumerate}
\item 影响
\item 很好
\item 教育
\item 心理辅导
\item 青少年
\end{enumerate}

\tcblower
\textbf{Solution détaillée :}
\begin{enumerate}
\item \zh{影响} $\rightarrow$ \textbf{yíngxiǎng} (2+3) — sandhi : yǐng (3e) + xiǎng (3e) $\rightarrow$ yíng (2e) + xiǎng.
\item \zh{很好} $\rightarrow$ \textbf{hénhǎo} (2+3) — sandhi : hěn (3e) + hǎo (3e) $\rightarrow$ hén (2e) + hǎo.
\item \zh{教育} $\rightarrow$ \textbf{jiàoyù} (4+4) — pas de sandhi (4e + 4e).
\item \zh{心理辅导} $\rightarrow$ \textbf{xīnlǐ fúdǎo} (1-3, 2-3) — \zh{心理} xīnlǐ (1+3, pas de sandhi car 1+3) ; \zh{辅导} fǔdǎo (3+3 $\rightarrow$ 2+3) $\rightarrow$ fúdǎo.
\item \zh{青少年} $\rightarrow$ \textbf{qīngshàonián} (1-4-2) — pas de sandhi.
\end{enumerate}
\end{exoresolu}

\courssub{Intonation des questions}

Dans un débat oral, vous poserez et recevrez des questions. L'intonation finale diffère selon le type de question.

\begin{propriete}[Intonation des questions en chinois]
\begin{itemize}
\item \textbf{Question fermée avec 吗 ma (ton neutre)} : la phrase monte globalement sur la dernière syllabe avant \zh{吗}, puis \zh{吗} se prononce haut et court. L'intonation globale est \textbf{montante}.
  \begin{itemize}
  \item \zh{你觉得教育重要吗？} Nǐ juéde jiàoyù zhòngyào ma? (Tu penses que l'éducation est importante ?) — montée sur \zh{要} yào, puis \zh{ma} haut.
  \item \zh{有监护人吗？} Yǒu jiānhùrén ma? (Y a-t-il un tuteur ?)
  \end{itemize}
\item \textbf{Question ouverte avec 为什么 wèishéme (4e + 2e + neutre)} : l'intonation est \textbf{descendante} sur la fin de phrase, comme une affirmation. Le mot \zh{为什么} porte l'interrogation, pas la mélodie finale.
  \begin{itemize}
  \item \zh{为什么青少年犯罪？} Wèishéme qīngshàonián fànzuì? (Pourquoi les adolescents commettent-ils des crimes ?) — chute sur \zh{罪} zuì.
  \item \zh{怎么预防？} Zěnme yùfáng? (Comment prévenir ?) — \zh{怎么} zěnme (3e+neutre $\rightarrow$ 2e+neutre), chute sur \zh{防} fáng.
  \end{itemize}
\item \textbf{Question alternative (A 还是 B)} : intonation montante sur A, descendante sur B.
\item \textbf{Question rhétorique} : intonation descendante affirmée.
\end{itemize}
\end{propriete}

\begin{exemplebox}[Contraste : 吗 vs 为什么]
\begin{tabularx}{\linewidth}{|X|X|X|}
\hline
\rowcolor{popblueL} \textbf{Chinois} & \textbf{Pinyin} & \textbf{Traduction} \\ \hline
你觉得重要吗？ & Nǐ juéde zhòngyào \textbf{ma?} & Est-ce que tu trouves ça important ? (montée) \\ \hline
为什么重要？ & Wèishéme zhòngyào\textbf{?} & Pourquoi est-ce important ? (descente) \\ \hline
有法律吗？ & Yǒu fǎlǜ \textbf{ma?} & Y a-t-il une loi ? (montée) \\ \hline
为什么没有法律？ & Wèishéme méiyǒu fǎlǜ\textbf{?} & Pourquoi n'y a-t-il pas de loi ? (descente) \\ \hline
\end{tabularx}
\end{exemplebox}

\begin{attention}[Piège : monter sur 为什么]
Ne dites pas \zh{*Wèishéme qīngshàonián fànzuì \uparrow} avec une montée finale. Cela sonne comme une question fermée mal formée. La question en \zh{为什么} se termine \textbf{bas}. Entraînez-vous : \zh{为什么} (4-2-neutre) + phrase affirmative descendante.
\end{attention}

\courssub{Exercices de discrimination auditive et chant choral}

\begin{experience}[Discrimination auditive — « Levez la carte du ton »]
\textbf{Matériel} : chaque élève a 5 cartes (1, 2, 3, 4, 0 pour ton neutre).
\textbf{Déroulement} :
\begin{enumerate}
\item L'énonciateur (enseignant ou élève désigné) lit un mot du vocabulaire de la leçon \textbf{une seule fois}, à vitesse naturelle.
\item Les élèves lèvent \textbf{immédiatement} la carte correspondant au ton de la \textbf{dernière syllabe} du mot (ou de la syllabe cible annoncée).
\item L'enseignant valide visuellement, corrige si besoin, fait répéter collectivement.
\end{enumerate}
\textbf{Exemples de mots à dicter} (cibler la dernière syllabe) :
\begin{itemize}
\item \zh{青少年} qīngshàonián $\rightarrow$ dernier ton = 2 (carte 2)
\item \zh{犯罪} fànzuì $\rightarrow$ 4 (carte 4)
\item \zh{原因} yuányīn $\rightarrow$ 1 (carte 1)
\item \zh{影响} yǐngxiǎng $\rightarrow$ 3 (carte 3) \emph{mais prononcé 2+3, dernière = 3}
\item \zh{教育} jiàoyù $\rightarrow$ 4 (carte 4)
\item \zh{觉得} juéde $\rightarrow$ neutre (carte 0)
\item \zh{监护人} jiānhùrén $\rightarrow$ neutre (carte 0)
\item \zh{预防} yùfáng $\rightarrow$ 2 (carte 2)
\item \zh{偷窃} tōuqiè $\rightarrow$ 4 (carte 4)
\item \zh{吸毒} xīdú $\rightarrow$ 2 (carte 2)
\end{itemize}
\textbf{Variante} : l'enseignant dicte une phrase complète du débat, les élèves lèvent la carte du ton de la syllabe finale de la phrase.
\end{experience}

\begin{methode}[Chant choral des formules fonctionnelles avec gestuelle des tons]
Choisissez 6 à 8 formules fonctionnelles du débat (section « Formules fonctionnelles »). Pour chacune :
\begin{enumerate}
\item L'enseignant modèle la phrase en exagérant légèrement la mélodie tonale, en faisant le geste de la main (méthode p. 1).
\item La classe répète \textbf{en chœur} en synchronisant le geste.
\item On répète 3 fois : (1) lent, gestes amples ; (2) vitesse normale, gestes réduits ; (3) vitesse débat, gestes internalisés (main sur table, mouvement imperceptible).
\item On alterne : moitié classe = proposition, moitié = relance/opinion.
\end{enumerate}
\textbf{Corpus suggéré} :
\begin{enumerate}
\item \zh{我觉得教育很重要。} Wǒ juéde jiàoyù hěn zhòngyào. — Tons : \zh{我} wǒ (3), \zh{觉得} juéde (2+0), \zh{教育} jiàoyù (4+4), \zh{很} hěn (3), \zh{重要} zhòngyào (4+4). Vérifier \zh{juéde} (2+neutre), \zh{zhòngyào} (4+4).
\item \zh{但是法律也必要。} Dànshì fǎlǜ yě bìyào. — Tons : \zh{但是} dànshì (4+4), \zh{法律} fǎlǜ (3+4), \zh{也} yě (3), \zh{必要} bìyào (4+4).
\item \zh{为什么不辅导？} Wèishéme bù fǔdǎo? — Tons : \zh{为什么} wèishéme (4+2+0), \zh{不} bù (4e ton, car suivi de 3e ton \zh{辅} fǔ), \zh{辅导} fǔdǎo $\rightarrow$ \textbf{fúdǎo} (2+3, sandhi). Intonation finale descendante.
\item \zh{我们应该预防。} Wǒmen yīnggāi yùfáng. — Tons : \zh{我们} wǒmen (3+neutre), \zh{应该} yīnggāi (1+1), \zh{预防} yùfáng (4+2).
\item \zh{这是社会问题。} Zhè shì shèhuì wèntí. — Tons : \zh{这} zhè (4), \zh{是} shì (4), \zh{社会} shèhuì (4+4), \zh{问题} wèntí (4+2).
\item \zh{监护人要负责。} Jiānhùrén yào fùzé. — Tons : \zh{监护人} jiānhùrén (1+4+0), \zh{要} yào (4), \zh{负责} fùzé (4+2).
\end{enumerate}
\end{methode}

\begin{savaistu}[Un ton, un sens — l'enjeu dans un débat noté]
En chinois, les tons sont \textbf{phonologiques} : ils distinguent les mots comme les phonèmes /p/ et /b/ en français. Exemples classiques :
\begin{itemize}
\item \zh{问} wèn (4e ton) = demander \quad vs \quad \zh{闻} wén (2e ton) = sentir, humer.
\item \zh{买} mǎi (3e ton) = acheter \quad vs \quad \zh{卖} mài (4e ton) = vendre.
\item \zh{睡觉} shuìjiào (4+4) = dormir \quad vs \zh{水饺} shuǐjiǎo (3+3$\rightarrow$2+3) = raviolis bouillis.
\end{itemize}
Dans un exposé ou un débat noté, dire \zh{*wén} au lieu de \zh{wèn} (« Je veux \textbf{sentir} une question » au lieu de « Je veux \textbf{poser} une question ») crée une confusion comique mais pénalisante. De même, \zh{*fànzuǐ} (boue retournée) au lieu de \zh{fànzuì} (criminalité) fait perdre la crédibilité de l'argumentaire. \textbf{Travailler les tons, c'est travailler le sens.}
\end{savaistu}

\begin{exoresolu}[Mini-dictée tonale — phrases du débat]
\textbf{Énoncé :} L'enseignant dicte les phrases suivantes. L'élève écrit le pinyin \textbf{avec tons réels (sandhi appliqué)} et la traduction.

1. \zh{我觉得预防很重要。}
2. \zh{为什么青少年犯罪？}
3. \zh{监护人应该辅导。}
4. \zh{教育比惩罚好。}

\tcblower
\textbf{Solution détaillée :}
\begin{enumerate}
\item \textbf{Wǒ juéde yùfáng hěn zhòngyào.} (Je pense que la prévention est très importante.)
  \begin{itemize}
  \item \zh{我} wǒ (3e) — \zh{觉得} juéde (2+0) — \zh{预防} yùfáng (4+2) — \zh{很} hěn (3e ton, devant 4e ton souvent réalisé en demi-3e ton bas ou 2e ton en flux rapide ; notation lexicale : hěn) — \zh{重要} zhòngyào (4+4).
  \end{itemize}
\item \textbf{Wèishéme qīngshàonián fànzuì?} (Pourquoi les adolescents commettent-ils des crimes ?)
  \begin{itemize}
  \item \zh{为什么} wèishéme (4+2+0) — \zh{青少年} qīngshàonián (1-4-2) — \zh{犯罪} fànzuì (4+4). Intonation finale descendante.
  \end{itemize}
\item \textbf{Jiānhùrén yīnggāi fúdǎo.} (Les tuteurs devraient encadrer / tutorater.)
  \begin{itemize}
  \item \zh{监护人} jiānhùrén (1+4+0) — \zh{应该} yīnggāi (1+1) — \zh{辅导} fǔdǎo $\rightarrow$ \textbf{fúdǎo} (2+3, sandhi).
  \end{itemize}
\item \textbf{Jiàoyù bǐ chéngfá hǎo.} (L'éducation vaut mieux que la punition.)
  \begin{itemize}
  \item \zh{教育} jiàoyù (4+4) — \zh{比} bǐ (3e ton, devant 2e ton \zh{chéng} $\rightarrow$ reste 3e, souvent demi-3e) — \zh{惩罚} chéngfá (2+2) — \zh{好} hǎo (3e ton final = 3e plein bas 214).
  \end{itemize}
\end{enumerate}
\end{exoresolu}

\begin{exercice}[Entraînement autonome — à faire en binôme]
\begin{enumerate}
\item \textbf{Lecture à vue} : l'un lit le tableau de vocabulaire (section 1) à voix haute, l'autre note les tons entendus sur une fiche (1/2/3/4/0). Inversez les rôles.
\item \textbf{Phrases à transformer} : prenez les 6 formules du chant choral. Changez \zh{我觉得} en \zh{我认为} wǒ rènwéi (3+2), \zh{教育} en \zh{管教} guǎnjiào (3+4), \zh{重要} en \zh{关键} guānjiàn (1+4). Répétez avec gestuelle.
\item \textbf{Questions croisées} : l'élève A pose une question en \zh{吗}, l'élève B répond puis pose une question en \zh{为什么}. Exemple : A : \zh{你觉得预防重要吗？} B : \zh{我觉得很重要。为什么你不觉得重要？} — Attention à l'intonation finale (montée vs descente).
\item \textbf{Enregistrement} : chaque élève s'enregistre (téléphone) en lisant les 10 mots de l'exercice de discrimination. Réécoutez, cochez les tons justes / faux, refaites les fautes 5 fois avec geste.
\end{enumerate}
\end{exercice}

\begin{corrige}[Exercice — Entraînement autonome]
\begin{enumerate}
\item \textbf{Auto-correction par les pairs} : pas de corrigé unique. L'enseignant circule, vérifie les fiches de tons, corrige les confusions 1/4 (plat vs chute) et 2/3 (montée vs creux).
\item \textbf{Transformations} :
  \begin{itemize}
  \item \zh{我认为管教很关键。} Wǒ rènwéi guǎnjiào hěn guānjiàn. (3+2, 3+4, 3, 1+4) — \zh{管教} guǎnjiào (3+4, pas de sandhi car 3+4) ; \zh{关键} guānjiàn (1+4).
  \item Vérifier : \zh{我认为} wǒ rènwéi (3+2) — \zh{管教} guǎnjiào (3+4) — \zh{很} hěn (3) — \zh{关键} guānjiàn (1+4).
  \end{itemize}
\item \textbf{Questions croisées} : valider que la question \zh{吗} monte (dernière syllabe avant \zh{ma} haute, \zh{ma} neutre haut) et que la question \zh{为什么} descend (dernière syllabe de la phrase basse). Modèle :
  \begin{itemize}
  \item A : \zh{你觉得预防重要吗？} Nǐ juéde yùfáng zhòngyào ma? $\nearrow$
  \item B : \zh{我觉得很重要。为什么你不觉得重要？} Wǒ juéde hěn zhòngyào. Wèishéme nǐ bù juéde zhòngyào? $\searrow$
  \end{itemize}
\item \textbf{Enregistrement} : critère de réussite = 9/10 tons justes sur les 10 mots cibles. Si échec, ré-enregistrer après 5 répétitions gestuelles par mot fautif.
\end{enumerate}
\end{corrige}

\coursec{Jeux de rôle et débat guidé}

Après avoir travaillé la prononciation des tons et les formules fonctionnelles pour présenter, donner son avis et relancer, nous passons maintenant à la mise en situation réelle. Cette section propose trois jeux de rôle ancrés dans le contexte camerounais et un débat guidé en classe entière. L'objectif est de mobiliser le vocabulaire, les structures grammaticales (\zh{比}, \zh{应该}, \zh{要}, \zh{因为…所以…}) et les formules de politesse (notamment \cle{您}) dans des échanges oraux fluides et argumentés.

\courssub{Jeu de rôle 1 : Interview radio — les « microbes » à Douala}

\begin{textebox}[Interview radio : un journaliste et un chef de quartier]
记者：王大哥，您好。最近杜拉的“微生物”现象很严重，您怎么看？\\
Jìzhě: Wáng Dàgē, nín hǎo. Zuìjìn Dùlā de ``wēishēngwù'' xiànxiàng hěn yánzhòng, nín zěnme kàn?\\
Journaliste : Bonjour Grand-frère Wang. Récemment, le phénomène des « microbes » à Douala est grave, qu'en pensez-vous ?

王大哥：记者同志，你好。这确实让人担心。这些年轻人大多没上学，没工作，整天在街上游荡。\\
Wáng Dàgē: Jìzhě tóngzhì, nǐ hǎo. Zhè quèshí ràng rén dānxīn. Zhèxiē niánqīngrén dàduō méi shàngxué, méi gōngzuò, zhěngtiān zài jiēshàng yóudàng.\\
Grand-frère Wang : Camarade journaliste, bonjour. C'est effectivement inquiétant. Ces jeunes pour la plupart ne vont pas à l'école, n'ont pas de travail, ils errent dans les rues toute la journée.

记者：那您觉得原因是什么？\\
Jìzhě: Nà nín juéde yuányīn shì shénme?\\
Journaliste : Selon vous, quelles en sont les causes ?

王大哥：首先，家庭教育缺失。很多父母忙着赚钱，管不了孩子。然后，学校离家远，学费贵。最后，社会缺乏对青年的引导。\\
Wáng Dàgē: Shǒuxiān, jiātíng jiàoyù quēshī. Hěnduō fùmǔ mángzhe zhuànqián, guǎn bù liǎo háizi. Ránhòu, xuéxiào lí jiā yuǎn, xuéfèi guì. Zuìhòu, shèhuì quēfá duì qīngnián de yǐndǎo.\\
Grand-frère Wang : D'abord, l'éducation familiale fait défaut. Beaucoup de parents sont occupés à gagner de l'argent, ne peuvent pas s'occuper des enfants. Ensuite, l'école est loin de la maison, les frais de scolarité sont chers. Enfin, la société manque d'encadrement pour la jeunesse.

记者：有什么解决办法吗？\\
Jìzhě: Yǒu shénme jiějué bànfǎ ma?\\
Journaliste : Y a-t-il des solutions ?

王大哥：我们应该加强社区管理，组织职业培训，让他们有事做，有钱赚。总之，预防比惩罚更重要。\\
Wáng Dàgē: Wǒmen yīnggāi jiāqiáng shèqū guǎnlǐ, zǔzhī zhíyè péixùn, ràng tāmen yǒu shì zuò, yǒu qián zhuàn. Zǒngzhī, yùfáng bǐ chéngfá gèng zhòngyào.\\
Grand-frère Wang : Nous devrions renforcer la gestion communautaire, organiser des formations professionnelles, pour qu'ils aient de quoi faire, de quoi gagner de l'argent. En résumé, la prévention est plus importante que la punition.
\end{textebox}

\begin{center}
\begin{tabularx}{\linewidth}{|X|X|X|X|}
\hline
\rowcolor{popblueL} Caractères & Pinyin & Nature & Traduction \\
\hline
微生物 & wēishēngwù & nom & microbe (argot : jeune délinquant de rue) \\
游荡 & yóudàng & verbe & errer, traîner \\
缺失 & quēshī & verbe/nom & faire défaut, absence \\
引导 & yǐndǎo & verbe/nom & guider, encadrement \\
职业培训 & zhíyè péixùn & nom & formation professionnelle \\
社区 & shèqū & nom & quartier, communauté \\
预防 & yùfáng & verbe/nom & prévention \\
惩罚 & chéngfá & verbe/nom & punition \\
\hline
\end{tabularx}
\end{center}

\begin{attention}[Registre de langue : 您 vs 你]
Dans le dialogue, le journaliste s'adresse au chef de quartier avec \cle{您 nín} (forme polie de « vous ») : \zh{您怎么看？} « Qu'en pensez-vous ? ». Le chef de quartier répond avec \cle{你 nǐ} car il parle à un cadet ou un égal. \textbf{Erreur fréquente} : utiliser \zh{你} avec un adulte, un supérieur, un policier ou un inconnu âgé. Règle : \textbf{adulte / supérieur / inconnu âgé → 您} ; pair / cadet / enfant → \textbf{你}.
\end{attention}

\begin{methode}[Préparer sa prise de parole en 3 étapes]
\begin{enumerate}
\item \textbf{Noter 3 idées structurées} avec connecteurs logiques : \zh{首先} (d'abord), \zh{然后} (ensuite), \zh{最后} (enfin) ou \zh{因为…所以…} (parce que… donc…).
\item \textbf{Répéter à voix haute} en marquant les tons sur chaque syllabe (ex. \zh{预防} \emph{yùfáng} 4-2, \zh{惩罚} \emph{chéngfá} 2-2). S'enregistrer si possible.
\item \textbf{Parler sans lire} : ne garder que les mots-clés sur une fiche, regarder son interlocuteur, utiliser les formules de relance (\zh{您觉得呢？} \zh{还有什么要补充的吗？}).
\end{enumerate}
\end{methode}

\courssub{Jeu de rôle 2 : Réunion de parents d'élèves — aider les jeunes en difficulté}

\begin{textebox}[Réunion parents-professeurs : que faire ?]
班主任：各位家长，大家好。今天我们讨论如何帮助有困难的学生。\\
Bānzhǔrèn: Gèwèi jiāzhǎng, dàjiā hǎo. Jīntiān wǒmen tǎolùn rúhé bāngzhù yǒu kùnnán de xuéshēng.\\
Professeur principal : Chers parents, bonjour. Aujourd'hui nous discutons de la façon d'aider les élèves en difficulté.

家长甲：老师，我儿子最近不想上学，天天跟那些“微生物”混在一起。\\
Jiāzhǎng Jiǎ: Lǎoshī, wǒ érzi zuìjìn bù xiǎng shàngxué, tiāntiān gēn nàxiē ``wēishēngwù'' hùn zài yīqǐ.\\
Parent A : Monsieur le prof, mon fils ne veut plus aller à l'école, il traîne tous les jours avec ces « microbes ».

班主任：这很危险。我们应该多和孩子沟通，了解他的想法。\\
Bānzhǔrèn: Zhè hěn wēixiǎn. Wǒmen yīnggāi duō hé háizi gōutōng, liǎojiě tā de xiǎngfǎ.\\
Professeur principal : C'est très dangereux. Nous devrions davantage communiquer avec l'enfant, comprendre sa pensée.

家长乙：我觉得学校也要加强管理，比如课后活动，别让他们闲着。\\
Jiāzhǎng Yǐ: Wǒ juéde xuéxiào yě yào jiāqiáng guǎnlǐ, bǐrú kèhòu huódòng, bié ràng tāmen xiánzhe.\\
Parent B : Je trouve que l'école doit aussi renforcer la gestion, par exemple les activités périscolaires, ne pas les laisser oisifs.

班主任：好建议。家校合作，共同预防。最后，总结一下：首先，家庭要多关心；然后，学校要提供活动；最后，社区要给机会。\\
Bānzhǔrèn: Hǎo jiànyì. Jiāxiào hézuò, gòngtóng yùfáng. Zuìhòu, zǒngjié yīxià: shǒuxiān, jiātíng yào duō guānxīn; ránhòu, xuéxiào yào tígōng huódòng; zuìhòu, shèqū yào gěi jīhuì.\\
Professeur principal : Bonne suggestion. Coopération famille-école, prévention commune. Pour conclure : d'abord, la famille doit plus se soucier ; ensuite, l'école doit proposer des activités ; enfin, le quartier doit donner des opportunités.
\end{textebox}

\begin{exemplebox}[Structures utiles pour proposer une solution]
\begin{itemize}
\item \zh{我们应该……} \emph{Wǒmen yīnggāi…} — Nous devrions…
\item \zh{建议……} \emph{Jiànyì…} — Je suggère que…
\item \zh{最好……} \emph{Zuìhǎo…} — Il vaudrait mieux…
\item \zh{要不……吧？} \emph{Yàobù… ba?} — Et si on… ? / Pourquoi ne pas… ?
\end{itemize}
\end{exemplebox}

\courssub{Jeu de rôle 3 : Au commissariat — la loi expliquée avec politesse}

\begin{textebox}[Commissariat : un policier et un jeune]
警察：年轻人，你今年多大？带身份证了吗？\\
Jǐngchá: Niánqīngrén, nǐ jīnnián duōdà? Dài shēnfènzhèng le ma?\\
Policier : Jeune homme, quel âge as-tu cette année ? As-tu ta carte d'identité ?

年轻人：叔叔，我十七岁。身份证没带。\\
Niánqīngrén: Shūshu, wǒ shíqī suì. Shēnfènzhèng méi dài.\\
Jeune : Monsieur l'agent (oncle), j'ai 17 ans. Je n'ai pas ma carte d'identité.

警察：你是学生吗？为什么不在学校？\\
Jǐngchá: Nǐ shì xuéshēng ma? Wèishénme bù zài xuéxiào?\\
Policier : Tu es élève ? Pourquoi n'es-tu pas à l'école ?

年轻人：我……我不想读书了。\\
Niánqīngrén: Wǒ… wǒ bù xiǎng dúshū le.\\
Jeune : Je… je ne veux plus étudier.

警察：听我说。根据法律，未满十八周岁是未成年人，受法律保护，但也要遵守法律。\\
Jǐngchá: Tīng wǒ shuō. Gēnjù fǎlǜ, wèi mǎn shíbā zhōusuì shì wèichéngniánrén, shòu fǎlǜ bǎohù, dàn yě yào zūnshǒu fǎlǜ.\\
Policier : Écoute-moi. Selon la loi, les moins de 18 ans sont des mineurs, protégés par la loi, mais doivent aussi respecter la loi.

年轻人：我知道了，叔叔。以后我会好好学习，不再惹事。\\
Niánqīngrén: Wǒ zhīdào le, shūshu. Yǐhòu wǒ huì hǎohǎo xuéxí, bù zài rěshì.\\
Jeune : J'ai compris, Monsieur l'agent. À l'avenir j'étudierai sérieusement, je ne ferai plus de bêtises.

警察：很好。记住，预防犯罪从自我约束开始。你回去吧，注意安全。\\
Jǐngchá: Hěn hǎo. Jìzhù, yùfáng fànzuì cóng zìwǒ yuēshù kāishǐ. Nǐ huíqù ba, zhùyì ānquán.\\
Policier : Très bien. Souviens-toi : la prévention du crime commence par l'autodiscipline. Tu peux rentrer, fais attention à ta sécurité.
\end{textebox}

\begin{propriete}[Vocabulaire juridique et civique de base]
\begin{center}
\begin{tabularx}{\linewidth}{|X|X|X|X|}
\hline
\rowcolor{popblueL} Caractères & Pinyin & Nature & Traduction \\
\hline
警察 & jǐngchá & nom & policier \\
身份证 & shēnfènzhèng & nom & carte d'identité \\
未成年人 & wèichéngniánrén & nom & mineur (moins de 18 ans) \\
法律 & fǎlǜ & nom & loi \\
保护 & bǎohù & verbe & protéger \\
遵守 & zūnshǒu & verbe & respecter, observer (loi, règle) \\
自我约束 & zìwǒ yuēshù & nom & autodiscipline \\
惹事 & rěshì & verbe & causer des ennuis, faire des bêtises \\
\hline
\end{tabularx}
\end{center}
\end{propriete}

\begin{attention}[Politesse obligatoire : 您 avec l'autorité]
Dans un commissariat, le jeune doit s'adresser au policier avec \cle{您 nín} ou un titre respectueux (\zh{叔叔 shūshu} « oncle », \zh{警官 jǐngguān} « agent »). Dire \zh{你} à un policier est perçu comme une provocation ou un manque d'éducation. \textbf{Modèle} : \zh{您好，警官。请问我有什么可以帮您的？} (Bonjour Monsieur l'agent. Puis-je vous aider ?).
\end{attention}

\courssub{Débat guidé en classe entière : « Faut-il punir ou éduquer ? »}

\begin{methode}[Organisation du débat — 30 minutes]
\begin{enumerate}
\item \textbf{Constitution des équipes} (5 min) : deux groupes de 6-8 élèves — Équipe A (\zh{主张惩罚} « pour la punition »), Équipe B (\zh{主张教育} « pour l'éducation »). Désigner : \zh{主持人} (animateur), \zh{发言人} (orateurs 1-3), \zh{计时员} (chronométreur), \zh{记录员} (rapporteur).
\item \textbf{Préparation arguments} (5 min) : chaque équipe note 4 arguments avec connecteurs (\zh{首先}、\zh{再者}、\zh{而且}、\zh{因此}) et un exemple concret (Cameroun ou Chine).
\item \textbf{Déroulement} (15 min) : voir organigramme ci-dessous.
\item \textbf{Bilan linguistique} (5 min) : l'enseignant note au tableau 3 phrases bien tournées et 3 erreurs de tons/registre à corriger collectivement.
\end{enumerate}
\end{methode}

\begin{popfigure}
\begin{tikzpicture}[
  node distance=1.2cm and 1.5cm,
  box/.style={rectangle, rounded corners, draw=popblue, thick, fill=popblueL, text width=2.8cm, align=center, minimum height=1.2cm},
  arrow/.style={-Stealth, thick, popblue},
  decision/.style={diamond, draw=poporange, thick, fill=poporangeL, text width=2.5cm, align=center, aspect=2}
]
\node[box, align=center] (start){Ouverture\\主持人\\Présente le thème\\et les règles};
\node[box, below=of start, align=center] (teamA){Tour 1 — Équipe A\\论点 1-2\\Arguments principaux};
\node[box, below=of teamA, align=center] (teamB){Tour 1 — Équipe B\\论点 1-2\\Arguments principaux};
\node[decision, below=of teamB, align=center] (free){Relances libres\\自由辩论\\Questions croisées\\Répliques};
\node[box, below=of free, align=center] (concl){Conclusion\\总之\\Synthèse + vote};
\draw[arrow] (start) -- (teamA);
\draw[arrow] (teamA) -- (teamB);
\draw[arrow] (teamB) -- (free);
\draw[arrow] (free) -- (concl);
% side labels
\node[left=0.3cm of start, text=popdark, font=\small\bfseries] {1 min};
\node[left=0.3cm of teamA, text=popdark, font=\small\bfseries] {3 min};
\node[left=0.3cm of teamB, text=popdark, font=\small\bfseries] {3 min};
\node[left=0.3cm of free, text=popdark, font=\small\bfseries] {6 min};
\node[left=0.3cm of concl, text=popdark, font=\small\bfseries] {2 min};
\end{tikzpicture}
\legende{Organigramme du déroulement du débat guidé}
\end{popfigure}

\begin{textebox}[Animateur : script d'ouverture et de conclusion]
主持人：同学们好。今天的辩题是：“青少年犯罪，应该重惩罚还是重教育？”\\
Zhǔchírén: Tóngxuémen hǎo. Jīntiān de biàntí shì: ``Qīngshàonián fànzuì, yīnggāi zhòng chéngfá háishì zhòng jiàoyù?''\\
Animateur : Bonjour camarades. Le sujet du débat d'aujourd'hui est : « Délinquance juvénile, faut-il privilégier la punition ou l'éducation ? »

主持人：现在请A组发言。\\
Zhǔchírén: Xiànzài qǐng A zǔ fāyán.\\
Animateur : Maintenant, la parole est au groupe A.

主持人：谢谢A组。现在请B组发言。\\
Zhǔchírén: Xièxie A zǔ. Xiànzài qǐng B zǔ fāyán.\\
Animateur : Merci groupe A. Maintenant, la parole est au groupe B.

主持人：进入自由辩论环节，双方可以提问。\\
Zhǔchírén: Jìnrù zìyóu biàolùn huánjié, shuāngfāng kěyǐ tíwèn.\\
Animateur : Place aux relances libres, les deux camps peuvent s'interroger.

主持人：辩论结束。请记录员汇总要点。最后，我们投票。\\
Zhǔchírén: Biàolùn jiéshù. Qǐng jìlùyuán huìzǒng yàodiǎn. Zuìhòu, wǒmen tóupiào.\\
Animateur : Fin du débat. Le rapporteur fait la synthèse. Enfin, nous votons.

主持人：总之，无论惩罚还是教育，目标都是让年轻人走上正路。谢谢大家。\\
Zhǔchírén: Zǒngzhī, wúlùn chéngfá háishì jiàoyù, mùbiāo dōu shì ràng niánqīngrén zǒushàng zhènglù. Xièxie dàjiā.\\
Animateur : En résumé, que ce soit punition ou éducation, le but est de remettre les jeunes sur le droit chemin. Merci à tous.
\end{textebox}

\begin{aretenir}[Phrases d'ouverture et de conclusion d'exposé — À MÉMORISER pour l'examen oral]
\begin{itemize}
\item \textbf{Ouverture} : \zh{大家好，今天我要谈谈……} \emph{Dàjiā hǎo, jīntiān wǒ yào tántan…} (Bonjour à tous, aujourd'hui je vais parler de…)
\item \textbf{Annonce du plan} : \zh{我从三个方面来说：首先……然后……最后……} \emph{Wǒ cóng sān gè fāngmiàn lái shuō: shǒuxiān… ránhòu… zuìhòu…} (J'aborde trois aspects : d'abord… ensuite… enfin…)
\item \textbf{Transition} : \zh{接下来，我们来看看……} \emph{Jiēxiálai, wǒmen lái kànkan…} (Ensuite, regardons…)
\item \textbf{Conclusion} : \zh{总之，我认为……希望大家……} \emph{Zǒngzhī, wǒ rènwéi… xīwàng dàjiā…} (En résumé, je pense que… j'espère que tous…)
\item \textbf{Remerciement} : \zh{谢谢大家的聆听。} \emph{Xièxie dàjiā de língtīng.} (Merci pour votre écoute.)
\end{itemize}
\end{aretenir}

\courssub{Grille d'observation pour l'enseignant (et auto-évaluation)}

\begin{center}
\begin{tabularx}{\linewidth}{|X|X|X|X|}
\hline
\rowcolor{popblueL} Critère & \textbf{4 — Excellent} & \textbf{2-3 — En progrès} & \textbf{1 — À travailler} \\
\hline
Prononciation des tons & Tons justes, naturels, sandhi 3-3 maîtrisé & Quelques erreurs de tons isolées & Tons plats ou confus, incompréhensibilité \\
\hline
Formules fonctionnelles & Utilise 5+ formules (ouvrir, relancer, conclure) & Utilise 2-3 formules, répétitives & Pas de formules, phrases hachées \\
\hline
Fluidité / Débit & Débit naturel, pauses logiques, pas de lecture & Débit lent, hésitations, lit parfois ses notes & Très lent, lit tout, longs silences \\
\hline
Argumentation & Idées claires, exemples précis (CM/CN), connecteurs & Idées présentes mais vagues, peu d'exemples & Idées confuses, pas d'exemple, pas de lien logique \\
\hline
Registre de langue & 您 / 叔叔 / 警官 utilisés correctement & Mélange 你/您, 1-2 maladresses & Tutoiement inapproprié, registre familier \\
\hline
\end{tabularx}
\end{center}

\courssub{Comparaison socioculturelle : Cameroun — Chine}

\begin{savaistu}[Rôle de la famille élargie vs École et Comité de quartier]
\begin{itemize}
\item \textbf{Cameroun} : La gestion des jeunes en difficulté repose largement sur la \textbf{famille élargie} (oncles, tantes, grands-parents) et les \textbf{chefs traditionnels} (chefs de quartier, lamidos, fons). La médiation familiale précède souvent la justice formelle. Le « conseil de famille » peut décider d'un placement chez un oncle au village pour « remettre sur le droit chemin ».
\item \textbf{Chine} : Depuis les années 1980, le \textbf{comité de quartier} (\zh{居委会 jūwěihuì}) et l'\textbf{école} sont les premiers relais. La \zh{预防青少年犯罪法} (Loi sur la prévention de la délinquance juvénile, 1999, révisée 2020) impose aux écoles des \zh{法治副校长} (vice-principaux chargés de l'éducation légale) et aux comités de quartier un suivi des mineurs « à risque ». La famille reste responsable (\zh{监护人责任}), mais l'État encadre fortement.
\item \textbf{Point commun} : Dans les deux pays, la \textbf{prévention} (\zh{预防}) prime officiellement sur la \textbf{répression} (\zh{惩罚}), mais les moyens diffèrent : solidarité lignagère au Cameroun, maillage administratif et éducatif en Chine.
\end{itemize}
\end{savaistu}

\begin{exoresolu}[Entraînement : transformer en chinois]
\textbf{Énoncé} : Traduisez en chinois (caractères + pinyin) les phrases suivantes pour le débat.
\begin{enumerate}
\item « Je pense que l'éducation est plus efficace que la punition. »
\item « Selon la loi chinoise, les mineurs de moins de 18 ans sont protégés. »
\item « Nous devrions organiser plus d'activités après les cours. »
\item « Monsieur le policier, j'ai compris, je ne le referai plus. »
\end{enumerate}

\tcblower
\textbf{Solution détaillée}
\begin{enumerate}
\item \zh{我认为教育比惩罚更有效。} \emph{Wǒ rènwéi jiàoyù bǐ chéngfá gèng yǒuxiào.} (Structure \zh{比} comparatif + \zh{更} « encore plus ».)
\item \zh{根据中国法律，未满十八岁的未成年人受保护。} \emph{Gēnjù Zhōngguó fǎlǜ, wèi mǎn shíbā suì de wèichéngniánrén shòu bǎohù.} (\zh{根据} « selon » + \zh{未满} « n'avoir pas atteint ».)
\item \zh{我们应该多组织一些课后活动。} \emph{Wǒmen yīnggāi duō zǔzhī yīxiē kèhòu huódòng.} (\zh{应该} + \zh{多} « davantage » + verbe.)
\item \zh{警官，我懂了，以后不敢了。} \emph{Jǐngguān, wǒ dǒng le, yǐhòu bù gǎn le.} (Registre poli \zh{警官}, \zh{不敢了} « je n'oserai plus » = engagement formel.)
\end{enumerate}
\end{exoresolu}

\begin{exercice}[Préparation au débat — À faire en binôme avant la séance]
\begin{enumerate}
\item Rédigez 3 arguments POUR la punition (Équipe A) en utilisant \zh{首先/然后/最后} et au moins un exemple camerounais ou chinois.
\item Rédigez 3 arguments POUR l'éducation (Équipe B) avec la même structure.
\item Choisissez 4 formules de relance dans la liste ci-dessous et écrivez une mini-réplique pour chacune :
\begin{itemize}
\item \zh{我不同意，因为……} (Je ne suis pas d'accord, car…)
\item \zh{您能举个例子吗？} (Pouvez-vous donner un exemple ?)
\item \zh{换个角度看，……} (Sous un autre angle, …)
\item \zh{也就是说，您的意思是……} (C'est-à-dire que vous voulez dire…)
\end{itemize}
\item Entraînez-vous à lire vos arguments à voix haute en marquant les tons (enregistrez-vous 1 minute).
\end{enumerate}
\end{exercice}

\begin{corrige}[Exercice — Éléments de réponse attendus]
\begin{enumerate}
\item \textbf{Équipe A (Punition)} : \zh{首先，法律必须有威信，重罚才能震慑。然后，缺乏管教的年轻人会再犯。最后，保护受害者也需要判决。} (D'abord, la loi doit avoir de l'autorité, seule une lourde peine dissuade. Ensuite, les jeunes sans encadrement récidivent. Enfin, protéger les victimes exige un jugement.)
\item \textbf{Équipe B (Éducation)} : \zh{首先，大多数未成年人犯罪因家庭缺位，教育能补救。然后，职业培训给他们出路，降低再犯率。最后，中国《预防未成年人犯罪法》明确“教育为主、惩罚为辅”。} (D'abord, la plupart des mineurs délinquants souffrent d'absence familiale, l'éducation peut combler. Ensuite, la formation professionnelle leur donne une issue, réduit la récidive. Enfin, la loi chinoise stipule « éducation prioritaire, punition accessoire ».)
\item \textbf{Relances types} :
\begin{itemize}
\item \zh{我不同意，因为惩罚不能解决根本原因，只有教育才能改变人。}
\item \zh{您能举个例子吗？比如杜拉的“微生物”进了职业学校后有多少回头？}
\item \zh{换个角度看，如果只靠坐牢，出来后更难找工作，反而更危险。}
\item \zh{也就是说，您的意思是建立更多专门学校而不是扩建监狱？}
\end{itemize}
\end{enumerate}
\end{corrige}

\begin{methode}[Check-list avant de passer à l'oral (examen ou débat)]
\begin{enumerate}
\item \square Mes 3 arguments sont notés avec connecteurs (\zh{首先/然后/最后}).
\item \square J'ai vérifié les tons des mots-clés (dictionnaire / prof / appli).
\item \square J'utilise \zh{您} pour m'adresser au jury / à l'enseignant / à l'adulte du jeu de rôle.
\item \square Je connais mes phrases d'ouverture (\zh{大家好，今天我要谈谈…}) et de conclusion (\zh{总之，我认为…谢谢大家。}).
\item \square J'ai un exemple concret (Cameroun ou Chine) pour appuyer mon point.
\item \square Je regarde mes interlocuteurs, je ne lis pas mes notes.
\end{enumerate}
\end{methode}

\newpage

\coursec{Activité d'intégration}

\coursec{Leçon 7 — Expression orale : la délinquance juvénile}

\courssub{Objectifs de l'activité}
Cette activité de synthèse vise à mobiliser l'ensemble des ressources linguistiques et communicationnelles travaillées dans la leçon 7 (vocabulaire thématique, formules fonctionnelles d'expression de l'opinion, de la concession et de la relance, structures grammaticales de l'argumentation, prononciation des tons) dans une situation de communication authentique, ancrée dans le contexte socioculturel camerounais. Les élèves devront être capables de :
\begin{itemize}
    \item structurer un exposé oral court (1 minute) de manière cohérente en utilisant des connecteurs logiques (\zh{首先、然后、最后} / \zh{一方面、另一方面}) ;
    \item exprimer un accord, un désaccord partiel ou total avec les formules appropriées (\zh{我同意、你说得有道理，但是……、我不同意}) ;
    \item relancer l'échange par des questions ouvertes (\zh{你觉得呢？、为什么呢？、怎么看？}) ;
    \item maîtriser la prononciation des tons sur le vocabulaire clé (notamment les tons 1/1/4 pour \zh{微生物}, les sandhi 3+3, les tons neutres) ;
    \item adopter une posture citoyenne et argumentée face à un problème de société réel.
\end{itemize}

\courssub{Situation}
\begin{textebox}[Situation déclenchante : Plateau TV « Jeunesse, Parlons-en ! »]
\textbf{Contexte :} La classe simule le plateau de l'émission « Jeunesse, Parlons-en ! » diffusée sur une chaîne nationale type CRTV ou Canal 2 International. L'animateur(trice) reçoit deux groupes de jeunes pour débattre d'un phénomène de société majeur à Douala : les « \zh{微生物} » (\zh{wēishēngwù}, Microbes), groupes de jeunes délinquants armés de machettes, couteaux, opérant dans les quartiers (New Bell, Bepanda, Akwa, Makepe).

\textbf{Document support (Lancement de l'animateur) :}
\zh{各位观众晚上好，欢迎收看《青年，我们说》。我是主持人 [Nom de l'élève]. 最近，杜阿拉的“微生物”现象引起了全社会的关注。青少年持刀抢劫、伤人，治安恶化。面对这一现象，大家议论纷纷：有人主张严厉惩罚，有人呼吁教育挽救。今晚的辩题是：青少年犯罪：应该惩罚还是应该教育？}
\emph{Guòwèi guānzhòng wǎnshàng hǎo, huānyíng shōukàn « Qīngnián, Wǒmen Shuō ». Wǒ shì zhǔchírén [Míngzi]. Zuìjìn, Dù'ālā de "Wēishēngwù" xiànxiàng yǐnqǐ le quán shèhuì de guānzhù. Qīngshàonián chí dāo qiǎngjié, shāngrén, zhì'ān èhuà. Miànduì zhè yī xiànxiàng, dàjiā yìlùn fēnfēn: yǒurén zhǔzhāng yánlì chéngfá, yǒurén hūyù jiàoyù wǎnjiù. Jīnwǎn de biàntí shì: Qīngshàonián fànzuì, yīnggāi chéngfá háishì yīnggāi jiàoyù?}

\textbf{Traduction :} Bonsoir chers téléspectateurs, bienvenue dans « La Jeunesse, on en parle ». Je suis l'animateur [Nom]. Récemment, le phénomène des « Microbes » à Douala suscite l'inquiétude de toute la société. Des jeunes armés volent et blessent, l'insécurité s'aggrave. Face à cela, les avis divergent : certains réclament une répression sévère, d'autres appellent à l'éducation et au sauvetage. Le sujet de ce soir : Délinquance juvénile : faut-il punir ou éduquer ?
\end{textebox}

\courssub{Consignes}
\begin{enumerate}
    \item \textbf{Constitution des groupes (2 min) :} La classe est divisée en 4 groupes : \textbf{Groupe A} (Équipe « Punition / Répression »), \textbf{Groupe B} (Équipe « Éducation / Réinsertion »), \textbf{Groupe C} (Animateur·trice + 1 co-animateur·trice), \textbf{Groupe D} (Jury d'évaluation : 3 à 4 élèves).
    \item \textbf{Préparation des arguments (5 min) :}
        \begin{itemize}
            \item Groupes A \& B : Rédiger 3 arguments structurés sur une fiche (mots-clés en caractères + pinyin). Utiliser obligatoirement : \zh{首先、然后、最后} (Premièrement, deuxièmement, enfin) ET au moins une structure de concession \zh{虽然……，但是……} ou \zh{你说得有道理，但是……}. Vocabulaire imposé : \zh{惩罚、法律、保护、受害者、教育、心理、家庭、社会、机会、改过自新}.
            \item Groupe C : Préparer la phrase d'accroche (cf. Situation), 2 questions de relance types (\zh{请问正方/反方，为什么……？、如果……，你们怎么看？}) et la phrase de conclusion type \zh{总之，今天的讨论很有意义……}.
            \item Groupe D : S'approprier la grille d'observation (voir Ressources).
        \end{itemize}
    \item \textbf{Déroulement du débat (15 min) :}
        \begin{enumerate}
            \item L'animateur lance le sujet (1 min).
            \item Tour 1 : Orateur 1 Groupe A (1 min) $\rightarrow$ Orateur 1 Groupe B (1 min).
            \item Relance de l'animateur + questions croisées (2 min).
            \item Tour 2 : Orateur 2 Groupe A (1 min) $\rightarrow$ Orateur 2 Groupe B (1 min).
            \item Débat libre modéré par l'animateur : \zh{你觉得呢？/ 为什么呢？} (3 min).
            \item Conclusion de l'animateur : \zh{总之……} (1 min).
        \end{enumerate}
    \item \textbf{Évaluation \& Débriefing (8 min) :}
        \begin{itemize}
            \item Le Jury (Groupe D) donne son avis argumenté en chinois simple (tons, formules, arguments, fluidité) et désigne l'équipe la plus convaincante.
            \item L'enseignant diffuse l'enregistrement vidéo/audio pour une auto-évaluation individuelle (grille élève).
        \end{itemize}
\end{enumerate}

\courssub{Ressources à mobiliser}
\begin{center}
\begin{tabularx}{\linewidth}{|X|X|X|X|}
\hline
\rowcolor{popblueL} \textbf{Caractères} & \textbf{Pinyin} & \textbf{Nature} & \textbf{Traduction} \\
\hline
青少年犯罪 & qīngshàonián fànzuì & Nom & Délinquance juvénile \\
微生物 & wēishēngwù & Nom & Microbes (terme camerounais pour gangs de jeunes) \\
惩罚 & chéngfá & Verbe/ Nom & Punir / Punition \\
严厉 & yánlì & Adj & Sévère, strict \\
法律 & fǎlǜ & Nom & Loi \\
保护 & bǎohù & Verbe & Protéger \\
受害者 & shòuhàizhě & Nom & Victime \\
教育 & jiàoyù & Verbe/ Nom & Éduquer / Éducation \\
挽救 & wǎnjiù & Verbe & Sauver, rattraper \\
心理 & xīnlǐ & Nom & Psychologie, mental \\
改过自新 & gǎiguòzìxīn & Idiome & Se corriger et recommencer à neuf (se réinsérer) \\
重新融入社会 & chóngxīn róngrù shèhuì & Locution & Réinsertion sociale \\
首先 / 然后 / 最后 & shǒuxiān / ránhòu / zuìhòu & Adv & D'abord / Ensuite / Enfin \\
虽然……，但是…… & suīrán……, dànshì…… & Conj & Bien que……, cependant…… \\
你说得有道理，但是…… & nǐ shuō de yǒu dàolǐ, dànshì…… & Phrase & Tu as raison, mais…… \\
我不同意 & wǒ bù tóngyì & Phrase & Je ne suis pas d'accord \\
你觉得呢？ & nǐ juéde ne? & Phrase & Qu'en penses-tu ? \\
为什么呢？ & wèishénme ne? & Phrase & Pourquoi ? \\
总之 & zǒngzhī & Adv & En bref, en conclusion \\
\hline
\end{tabularx}
\end{center}

\begin{propriete}[Structures argumentatives clés (Niveau HSK 2-3)]
\textbf{1. Énumérer des arguments :}
\zh{结构：首先，……；然后，……；最后，……}
\zh{例：首先，法律必须严厉惩罚犯罪；然后，保护受害者；最后，维护社会治安。}
\emph{Shǒuxiān, fǎlǜ bìxū yánlì chéngfá fànzuì; ránhòu, bǎohù shòuhàizhě; zuìhòu, wéihù shèhuì zhì'ān.}

\textbf{2. Exprimer une concession (argument adverse + son contre-argument) :}
\zh{结构：虽然 [Argument adverse] ，但是 [Mon argument principal] 。}
\zh{例：虽然他们还年轻，但是犯罪是不对的，必须负责任。}
\emph{Suīrán tāmen hái niánqīng, dànshì fànzuì shì bùduì de, bìxū fù zérèn.}

\textbf{3. Désaccord poli + contre-argument (Formule fonctionnelle) :}
\zh{结构：你说得有道理，但是 [Mon point de vue] 。}
\zh{例：你说得有道理，教育很重要，但是没有法律惩罚，坏人不怕。}
\emph{Nǐ shuō de yǒu dàolǐ, jiàoyù hěn zhòngyào, dànshì méiyǒu fǎlǜ chéngfá, huàirén bù pà.}

\textbf{4. Relancer / Demander l'avis :}
\zh{你觉得呢？} \emph{Nǐ juéde ne?} (Tu en penses quoi ?) \quad
\zh{为什么呢？} \emph{Wèishénme ne?} (Pourquoi ?) \quad
\zh{怎么看？} \emph{Zěnme kàn?} (Comment tu vois ça ?)
\end{propriete}

\begin{attention}[Pièges phonétiques et grammaticaux fréquents]
\begin{itemize}
    \item \textbf{Tons de \zh{微生物} :} \zh{微} (wēi, 1er ton), \zh{生} (shēng, 1er ton), \zh{物} (wù, 4e ton) $\rightarrow$ \textbf{wēi shēng wù (1-1-4)}. Attention à ne pas prononcer \zh{微} au 2e ton (wéi).
    \item \textbf{Sandhi 3e ton + 3e ton :} \zh{想改} (xiǎng gǎi) $\rightarrow$ \textbf{xiáng gǎi} (2-3). \zh{改过自新} : \zh{改} (gǎi, 3) + \zh{过} (guò, 4) $\rightarrow$ pas de sandhi (3-4).
    \item \textbf{Neutre \zh{了} (le) vs \zh{过} (guò) :} \zh{他犯了错} (Tā fàn le cuò - Il a fait une erreur, action accomplie) vs \zh{他犯过错} (Tā fàn guò cuò - Il a déjà fait une erreur, expérience passée). Pour le débat : \zh{他们犯过法} (Tāmen fàn guò fǎ - Ils ont déjà commis des délits).
    \item \textbf{Ordre des mots (Sujet + Aux + Verbe + Objet) :} Ne pas dire \zh{惩罚他们应该}. Dire : \zh{应该惩罚他们} (Yīnggāi chéngfá tāmen).
    \item \textbf{Nuances lexicales :} \zh{教育} ne signifie pas seulement « scolarité » mais « éducation morale, correction, rééducation ». \zh{惩罚} implique une sanction légale/judiciaire, pas seulement une punition parentale (\zh{惩罚} vs \zh{罚}).
\end{itemize}
\end{attention}

\courssub{Proposition de production et corrigé commenté}
\begin{exoresolu}[Modèle de débat simulé (Extrait)]
\textbf{Énoncé :} Produire un extrait du débat (Animateur, Équipe Punition, Équipe Éducation, Conclusion) en respectant les contraintes de temps, de vocabulaire et de structures.

\tcblower

\textbf{Solution détaillée (Script annoté) :}

\begin{textebox}[Animateur (主持人 Zhǔchírén)]
\zh{各位观众晚上好！欢迎收看《青年，我们说》。今天我们讨论一个沉重的话题：杜阿拉的“微生物”现象。青少年持刀抢劫，市民很害怕。辩题：青少年犯罪，应该惩罚还是应该教育？请正方先发言。}
\emph{Guòwèi guānzhòng wǎnshàng hǎo! Huānyíng shōukàn « Qīngnián, Wǒmen Shuō ». Jīntiān wǒmen tǎolùn yīgè chénzhòng de huàtí: Dù'ālā de "Wēishēngwù" xiànxiàng. Qīngshàonián chí dāo qiǎngjié, shìmín hěn hàipà. Biàntí: Qīngshàonián fànzuì, yīnggāi chéngfá háishì yīnggāi jiàoyù? Qǐng zhèngfāng xiān fāyán.}
« Bonsoir chers téléspectateurs ! Bienvenue dans « La Jeunesse, on en parle ». Aujourd'hui nous discutons d'un sujet grave : le phénomène "Microbes" à Douala. Des jeunes armés volent, les citoyens ont peur. Sujet : Délinquance juvénile, punir ou éduquer ? Je donne la parole à l'équipe "Pour la punition". »
\end{textebox}

\begin{textebox}[Équipe Punition - Orateur 1 (正方一辩 Zhèngfāng Yībiàn)]
\zh{我认为，必须严厉惩罚。首先，法律面前人人平等，犯罪就要坐牢，这是保护受害者的权利。然后，如果不惩罚，他们会变本加厉，社会治安会更乱。最后，严厉的惩罚能起到警示作用，让其他青少年不敢犯罪。}
\emph{Wǒ rènwéi, bìxū yánlì chéngfá. Shǒuxiān, fǎlǜ miànqián rénrén píngděng, fànzuì jiù yào zuò láo, zhè shì bǎohù shòuhàizhě de quánlì. Ránhòu, rúguǒ bù chéngfá, tāmen huì biàn běn jiā lì, shèhuì zhì'ān huì gèng luàn. Zuìhòu, yánlì de chéngfá néng qǐ dào jǐngshì zuòyòng, ràng qítā qīngshàonián bù gǎn fànzuì.}
« Je pense qu'il faut punir sévèrement. Primo, devant la loi tous sont égaux, crime = prison, ça protège les droits des victimes. Secundo, si on ne punit pas, ils iront crescendo, l'insécurité empire. Tertio, une punition sévère a un effet dissuasif, les autres jeunes n'oseront pas. »
\end{textebox}

\begin{textebox}[Équipe Éducation - Orateur 1 (反方一辩 Fǎngfāng Yībiàn)]
\zh{你说得有道理，惩罚是必要的。但是，我认为单纯惩罚不能解决根本问题。首先，这些“微生物”大多来自贫困家庭，缺乏父母管教，心理扭曲。然后，监狱只会让他们学到更坏的手段。最后，我们应该给他们机会，通过职业教育、心理辅导，帮他们改过自新，重新融入社会。}
\emph{Nǐ shuō de yǒu dàolǐ, chéngfá shì bìyào de. Dànshì, wǒ rènwéi dānchún chéngfá bù néng jiějué gēn běn wèn tí. Shǒuxiān, zhèxiē "Wēishēngwù" dàduō lái zì pín kùn jiātíng, quēfá fùmǔ guǎnjiào, xīnlǐ niǔqǔ. Ránhòu, jiānyù zhǐ huì ràng tāmen xué dào gèng huài de shǒu duàn. Zuìhòu, wǒmen yīnggāi gěi tāmen jīhuì, tōngguò zhí yè jiàoyù, xīnlǐ fǔdǎo, bāng tāmen gǎi guò zì xīn, chóng xīn róng rù shè huì.}
« Tu as raison, la punition est nécessaire. Mais je pense que la seule punition ne résout pas le fond. Primo, ces "Microbes" viennentmostly de familles pauvres, manquent d'encadrement parental, sont perturbés psychologiquement. Secundo, la prison ne fait qu'apprendre de pires méthodes. Tertio, il faut leur donner une chance : formation pro, soutien psy, les aider à se réinsérer, réintégrer la société. »
\end{textebox}

\begin{textebox}[Relance Animateur \& Débat libre]
\zh{主持人：正方，反方说监狱让他们变坏，你怎么看？}
\zh{正方二辩：虽然监狱环境不好，但是没有法律制裁，受害者谁来保护？你说得有道理，教育要跟上，但惩罚不能少。}
\zh{反方二辩：为什么非要二选一？我认为“惩罚为了更好地教育”。判刑的同时，必须强制职业培训。}
\emph{Zhǔchírén: Zhèngfāng, fǎngfāng shuō jiān yù ràng tāmen biàn huài, nǐ zěnme kàn? / Zhèngfāng Èrbiàn: Suīrán jiān yù huán jìng bù hǎo, dànshì méiyǒu fǎ lǜ zhì cái, shòu hài zhě shéi lái bǎo hù? Nǐ shuō de yǒu dàolǐ, jiào yù yào gēn shàng, dàn chéng fá bù néng shǎo. / Fǎngfāng Èrbiàn: Wèi shén me fēi yào èr xuǎn yī? Wǒ rèn wéi "chéng fá wèi le gèng hǎo de jiào yù". Pàn xíng de tóng shí, bì xū qiáng zhì zhí yè péi xùn.}
\end{textebox}

\begin{textebox}[Conclusion Animateur (主持人 总结)]
\zh{谢谢双方精彩的辩论。总之，面对“微生物”现象，单纯靠惩罚治标不治本，单纯靠教育又来不及保护市民。法律要严厉，教育要跟上，家庭、学校、社会要共同努力，给迷途少年一条回头路。谢谢观众，下期再见！}
\emph{Xièxie shuāngfāng jīngcǎi de biànlùn. Zǒngzhī, miànduì "Wēishēngwù" xiànxiàng, dānchún kào chéngfá zhì biāo bù zhì běn, dānchún kào jiàoyù yòu lái bù jí bǎohù shìmín. Fǎlǜ yào yánlì, jiàoyù yào gēn shàng, jiātíng, xuéxiào, shèhuì yào gòngtóng nǔlì, gěi mítú shàonián yī tiáo huítóu lù. Xièxie guānzhòng, xiàqī zàijiàn!}
« Merci aux deux camps pour ce débat brillant. En bref, face aux "Microbes", la seule punition traite le symptôme pas la cause, la seule éducation est trop lente pour protéger les citoyens. La loi doit être sévère, l'éducation doit suivre, famille, école, société doivent unir leurs efforts pour offrir un chemin de retour aux jeunes égarés. Merci téléspectateurs, à la prochaine ! »
\end{textebox}

\textbf{Commentaire des choix de langue (Pour l'enseignant) :}
\begin{itemize}
    \item \textbf{Registre :} Langage soutenu mais accessible (HSK 3), adapté à un plateau TV grand public camerounais. Utilisation de \zh{沉重的话题} (sujet grave), \zh{治标不治本} (traiter le symptôme pas la cause - chéngyǔ niveau 4 mais très courant à l'oral).
    \item \textbf{Connecteurs :} Respect strict de \zh{首先/然后/最后} pour la clarté expositive. Utilisation de \zh{虽然……但是……} (Orateur 2 réplique) et \zh{你说得有道理，但是……} (Orateur 1 réplique) pour modéliser la concession argumentative.
    \item \textbf{Vocabulaire ciblé :} \zh{微生物} (terme local intégré), \zh{变本加厉} (aller crescendo/empirer - idiome utile), \zh{警示作用} (effet dissuasif), \zh{心理辅导} (soutien psy), \zh{重新融入社会} (réinsertion).
    \item \textbf{Phonétique :} Attention aux tons de \zh{微生物} (1-1-4), \zh{变本加厉} (4-3-1-4), \zh{改过自新} (3-4-4-1). L'animateur doit modéliser l'intonation montante de \zh{你怎么看？} et la chute de \zh{总之……}.
    \item \textbf{Synthèse :} La conclusion de l'animateur synthétise les deux thèses (thèse/antithèse/synthèse) en réutilisant \zh{单纯……单纯……} (structure parallélisme) et \zh{要……要……要……} (impératif moral partagé).
\end{itemize}
\end{exoresolu}

\courssub{Bilan de l'activité}
Cette activité d'intégration a permis de vérifier l'atteinte des compétences orales du Module 1 dans une mise en situation complexe et motivante.
\begin{itemize}
    \item \textbf{Compétence lexicale :} Les élèves ont réinvesti le vocabulaire spécifique (délinquance, justice, éducation, psychologie, société) et les connecteurs logiques dans une production spontanée, sortant de l'exercice de répétition.
    \item \textbf{Compétence grammaticale :} L'usage fonctionnel des structures \zh{虽然……但是……}, \zh{你觉得呢？}, \zh{应该……} a été évalué en interaction réelle. La correction de l'ordre des mots (Suj + Adv + V + Obj) et de l'aspect (\zh{过} pour l'expérience) a pu être observée.
    \item \textbf{Compétence phonologique :} L'enregistrement vidéo/audio offre un support concret pour l'auto-correction des tons (notamment le sandhi du 3e ton, les tons neutres sur \zh{了、呢、的}) et de l'intonation interrogative/assertive.
    \item \textbf{Compétence socioculturelle :} L'ancrage dans le phénomène des « Microbes » à Douala a permis de lier apprentissage linguistique et éducation à la citoyenneté (réflexion sur causes sociales : pauvreté, échec scolaire, absence parentale ; solutions : justice, éducation, solidarité). La comparaison implicite avec la gestion chinoise de la délinquance juvénile (centres de rééducation \zh{少管所}, travail éducatif) peut être prolongée en devoir de recherche.
    \item \textbf{Pédagogie différentielle :} Les rôles (orateur, animateur, juré) permettent à chaque profil d'élève (à l'aise à l'oral, bon analyste, bon organisateur) de valoriser une compétence. La grille d'évaluation par les pairs (Groupe D) développe l'esprit critique et la méta-cognition linguistique.
\end{itemize}
\textbf{Prolongement possible :} Rédaction d'un article de journal scolaire (Expression écrite) ou d'une lettre ouverte au Ministre de la Jeunesse (Traduction/Redaction) reprenant les arguments du débat.

\newpage

\coursec{Méthodes et exercices résolus}

\coursec{Leçon 7 — Expression orale : Délinquance juvénile}

\courssub{Vocabulaire thématique : La délinquance juvénile}
\begin{definition}[Lexique HSK 3-4 — Thème : Société \& Droit]
\textbf{Consigne :} Mémorisez les caractères, le pinyin (tons sandhi appliqués à l'oral), la nature et la traduction. Travaillez les paires homophones / paronymes en colonne 4.
\begin{center}
\begin{tabularx}{\linewidth}{|X|X|X|X|}
\hline
\rowcolor{popblueL} \textbf{Caractères} & \textbf{Pinyin} & \textbf{Nature} & \textbf{Traduction / Note} \\
\hline
\zh{青少年} & qīngshàonián (qíngshàonián) & n. & adolescents, jeunes ; \textit{sandhi 3+3 $\rightarrow$ 2+3} \\
\zh{犯罪} & fànzuì (fánzuì) & v./n. & commettre un crime / criminalité ; \textit{sandhi 3+3 $\rightarrow$ 2+4} \\
\zh{青少年犯罪} & qīngshàonián fànzuì & n. & délinquance juvénile \\
\zh{管教} & guǎnjiào (guánjiào) & v. & discipliner, éduquer (souvent strict) ; \textit{sandhi 3+4 $\rightarrow$ 2+4} \\
\zh{预防} & yùfáng & v. & prévenir \\
\zh{矫正} & jiǎozhèng (jiáozhèng) & v./n. & corriger (comportement) / correction ; \textit{sandhi 3+4 $\rightarrow$ 2+4} \\
\zh{再犯} & zàifàn & v. & récidiver \\
\zh{震慑} & zhèndé & v. & dissuader, intimider \\
\zh{法治} & fǎzhì & n. & État de droit, gouvernance par la loi \\
\zh{法制} & fǎzhì & n. & système juridique, ensemble des lois \\
\zh{不良诱惑} & bùliáng yòuhuò & n. & mauvaises tentations, influences néfastes \\
\zh{心理扭曲} & xīnlǐ niǔqū & n. & déviation / distorsion psychologique \\
\zh{陪伴} & péibàn & v. & accompagner, passer du temps avec \\
\zh{联动} & liándòng & v. & agir en synergie, coordination multi-acteurs \\
\zh{社区} & shèqū & n. & communauté, quartier (unité administrative de base) \\
\zh{判处} & pànchǔ & v. & condamner à (une peine) — registre juridique \\
\zh{改过自新} & gǎiguòzìxīn & id. & corriger ses fautes et recommencer à neuf (chengyu) \\
\zh{一起} & yī qǐ & m. & classif. pour affaires, cas (juridiques) \\
\zh{宗} & zōng & m. & classif. pour dossiers, crimes graves \\
\hline
\end{tabularx}
\end{center}
\textbf{Paires à distinguer (tons + sens) :}
\begin{itemize}
    \item \zh{管教} (guǎnjiào, discipliner) vs \zh{官教} (guānjiào, enseignement officiel — rare).
    \item \zh{一律} (yīlǜ $\rightarrow$ yìlǜ, uniformément) vs \zh{依律} (yīlǜ, selon la loi).
    \item \zh{放任} (fàngrèn, laisser faire / négliger) vs \zh{放人} (fàngrén, libérer qqn).
    \item \zh{重罚} (zhòngfá, peine lourde) vs \zh{种法} (zhǒngfǎ, méthode de culture).
    \item \zh{法治} (fǎzhì, rule of law) vs \zh{法制} (fǎzhì, legal system).
\end{itemize}
\end{definition}

\courssub{Dialogue modèle 1 : Parler du problème (Registre courant)}
\begin{textebox}[Dialogue 1 — Deux lycéens discutent de l'actualité]
\zh{A : 最近新闻里青少年犯罪案件真不少。} (Zuìjìn xīnwén lǐ qīngshàonián fànzuì ànjiàn zhēn bùshǎo.) \\
\zh{B : 是啊，听说很多是因为家庭管教缺失。} (Shì a, tīngshuō hěnduō shì yīnwèi jiātíng guǎnjiào quēshī.) \\
\zh{A : 学校也得加强法治教育，不能光靠家庭。} (Xuéxiào yě děi jiāqiáng fǎzhì jiàoyù, bùnéng guāng kào jiātíng.) \\
\zh{B : 对，还得社区组织活动，比如踢足球，让孩子少接触不良诱惑。} (Duì, hái děi shèqū zǔzhī huódòng, bǐrú tī zúqiú, ràng háizi shǎo jiēchù bùliáng yòuhuò.) \\
\zh{A : 总之，预防青少年犯罪，全社会得联动。} (Zǒngzhī, yùfáng qīngshàonián fànzuì, quán shèhuì děi liándòng.)
\end{textebox}
\begin{aretenir}[Structures clés du dialogue 1]
\begin{itemize}
    \item \zh{真不少} (zhēn bùshǎo) : « vraiment pas peu » = « vraiment beaucoup » (litote courante).
    \item \zh{缺失} (quēshī) : « manquer, faire défaut » (registre plus formel que \zh{没有}).
    \item \zh{得} (děi) : « devoir, falloir » (oral, fort) — interchangeable avec \zh{必须} (bìxū) ou \zh{应该} (yīnggāi).
    \item \zh{光靠} (guāng kào) : « seulement compter sur / se fier uniquement à ».
    \item \zh{少接触} (shǎo jiēchù) : « réduire le contact avec » — structure V + O.
\end{itemize}
\end{aretenir}

\courssub{Formules fonctionnelles : Présenter, Opiner, Relancer}
\begin{propriete}[Boîte à outils pour l'oral interactif (Niveau Bac)]
\begin{center}
\begin{tabularx}{\linewidth}{|X|X|X|}
\hline
\rowcolor{popblueL} \textbf{Fonction} & \textbf{Formules chinoises (Caractères + Pinyin)} & \textbf{Équivalent français} \\
\hline
\textbf{Accroche / Lancer le sujet} & \zh{今天我们讨论……} (Jīntiān wǒmen tǎolùn…) & Aujourd'hui nous discutons de… \\
& \zh{关于……这个问题，我想说……} (Guānyú… zhège wèntí, wǒ xiǎng shuō…) & Concernant ce problème…, je veux dire… \\
\hline
\textbf{Donner son avis} & \zh{我认为 / 我看 / 据我看……} (Wǒ rènwéi / Wǒ kàn / Jù wǒ kàn…) & Je pense / À mon avis… \\
& \zh{在我看来……} (Zài wǒ kànlái…) & À mes yeux… \\
\hline
\textbf{Exprimer l'accord} & \zh{我赞同你的看法。} (Wǒ zàntóng nǐ de kànfǎ.) & J'approuve ton avis. \\
& \zh{你说得很对 / 完全正确。} (Nǐ shuō de hěn duì / Wánquán zhèngquè.) & Tu as raison / Tout à fait exact. \\
\hline
\textbf{Exprimer le désaccord (poli)} & \zh{我不太同意……} (Wǒ bù tài tóngyì…) & Je ne suis pas tout à fait d'accord… \\
& \zh{恐怕……} (Kǒngpà…) & Je crains que… / Il me semble que… (adoucit) \\
& \zh{这倒不一定。} (Zhè dào bù yídìng.) & Ce n'est pas forcément vrai. \\
\hline
\textbf{Nuancer / Concéder} & \zh{虽然……，但是……} (Suīrán…, dànshì…) & Bien que…, cependant… \\
& \zh{一方面……，另一方面……} (Yīfāngmiàn…, lìngyī fāngmiàn…) & D'un côté…, de l'autre… \\
\hline
\textbf{Relancer / Demander précision} & \zh{你为什么这么认为？} (Nǐ wèishénme zhème rènwéi?) & Pourquoi penses-tu cela ? \\
& \zh{有什么证据吗？} (Yǒu shénme zhèngjù ma?) & Y a-t-il des preuves ? \\
& \zh{能举个例子吗？} (Néng jǔ gè lìzi ma?) & Peux-tu donner un exemple ? \\
\hline
\textbf{Gagner du temps} & \zh{让我想想……} (Ràng wǒ xiǎngxiǎng…) & Laissez-moi réfléchir… \\
& \zh{这个问题挺复杂的。} (Zhège wèntí tǐng fùzá de.) & Cette question est assez complexe. \\
\hline
\textbf{Conclure} & \zh{总之 / 综上所述……} (Zǒngzhī / Zōngshàng shùshù…) & En somme / En conclusion… \\
& \zh{希望……} (Xīwàng…) & J'espère que… \\
\hline
\end{tabularx}
\end{center}
\end{propriete}

\courssub{Dialogue modèle 2 : Le débat — Punir ou éduquer ? (Registre formel / soutenu)}
\begin{textebox}[Dialogue 2 — Simulation procès / débat TV]
\zh{主持人 : 今天我们讨论一起青少年盗窃案。检察官请发言。} (Zhǔchírén: Jīntiān wǒmen tǎolùn yī qǐ qīngshàonián dàoqiè àn. Jiǎncháguān qǐng fāyán.) \\
\zh{检察官 : 我认为必须严厉惩罚，这样才能震慑别人，维护法律尊严。} (Jiǎncháguān: Wǒ rènwéi bìxū yánlì chéngfá, zhèyàng cáinéng zhèndé biérén, wéihù fǎlǜ zūnyán.) \\
\zh{辩护律师 : 我不太同意。虽然震慑重要，但重罚不等于预防再犯。} (Biànhù lǜshī: Wǒ bù tài tóngyì. Suīrán zhèndé zhòngyào, dàn zhòngfá bù děngyú yùfáng zàifàn.) \\
\zh{辩护律师 : 据统计，接受社区矫正的青少年再犯率远低于坐牢者。} (Biànhù lǜshī: Jù tǒngjì, jiēshòu shèqū jiǎozhèng de qīngshàonián zàifàn lǜ yuǎn dīyú zuòláo zhě.) \\
\zh{心理专家 : 一方面要惩罚，另一方面要疏导心理，找到犯罪根源。} (Xīnlǐ zhuānjiā: Yīfāngmiàn yào chéngfá, lìngyī fāngmiàn yào shūdǎo xīnlǐ, zhǎodào fànzuì gēnyuán.) \\
\zh{主持人 : 综合双方意见，建议法庭判处社区矫正，配合心理辅导。} (Zhǔchírén: Zōnghé shuāngfāng yìjiàn, jiànyì fǎtíng pànchǔ shèqū jiǎozhèng, pèihé xīnlǐ fǔdǎo.)
\end{textebox}
\begin{exemplebox}[Analyse linguistique du dialogue 2]
\begin{itemize}
    \item \zh{一起} (yī qǐ) : Classificateur correct pour « une affaire / un cas ».
    \item \zh{维护法律尊严} (wéihù fǎlǜ zūnyán) : Collocation fixe « défendre la dignité de la loi ».
    \item \zh{不等于} (bù děngyú) : « Ne pas équivaloir à » — structure logique précise.
    \item \zh{远低于} (yuǎn dī yú) : Comparatif « bien inférieur à » (HSK 4).
    \item \zh{疏导} (shūdǎo) : « guider, canaliser (émotions, foule) » — terme psy/éduca spécifique.
    \item \zh{配合} (pèihé) : « en combinaison avec, assorti de » — structure V + O comme complément de manière.
\end{itemize}
\end{exemplebox}

\courssub{Phonétique : Les tons à travailler (Focus HSK 3-4)}
\begin{methode}[Maîtriser les tons du vocabulaire thématique — Méthode « Karaoké »]
\textbf{Objectif :} Automatiser les sandhi (changements de tons) et distinguer les paires minimales pour un exposé fluide et intelligible.
\begin{enumerate}
    \item \textbf{Règle des 3e tons consécutifs (Sandhi obligatoire) :} 3 + 3 $\rightarrow$ 2 + 3.
    \begin{itemize}
        \item \zh{青少年} : qīng(1) shǎo(3) nián(2) $\rightarrow$ \textbf{qíng(2) shǎo(3) nián(2)}.
        \item \zh{管教} : guǎn(3) jiào(4) $\rightarrow$ \textbf{guán(2) jiào(4)}.
        \item \zh{犯罪} : fàn(4) zuì(4) — \textit{pas de sandhi 3+3 ici, mais dictionnaire souvent note 3e ton pour 犯 seul, 4e en composé. À l'oral : fán(2) zuì(4) par assimilation anticipée fréquente, mais standard : fàn(4) zuì(4).} \textbf{Standard enseigné : fàn zuì (4-4).}
        \item \zh{矫正} : jiǎo(3) zhèng(4) $\rightarrow$ \textbf{jiáo(2) zhèng(4)}.
    \end{itemize}
    \item \textbf{\cle{不} (bù) :} 4e ton $\rightarrow$ 2e ton (\textbf{bú}) devant un 4e ton.
    \begin{itemize}
        \item \zh{不稳定} : bù(4) wěn(3) dìng(4) $\rightarrow$ \textbf{bú(2) wěn(3) dìng(4)}.
        \item \zh{不严重} : \textbf{bú(2) yán(2) zhòng(4)}.
        \item \zh{不教育} : \textbf{bú(2) jiào(4) yù(4)}.
    \end{itemize}
    \item \textbf{\cle{一} (yī) :} 1er ton $\rightarrow$ 4e ton (\textbf{yì}) devant 1, 2, 3 ; $\rightarrow$ 2e ton (\textbf{yí}) devant 4.
    \begin{itemize}
        \item \zh{一律} : yī(1) lǜ(4) $\rightarrow$ \textbf{yì(4) lǜ(4)}.
        \item \zh{一个} : \textbf{yí(2) gè(4)}.
    \end{itemize}
    \item \textbf{Tons neutres (轻声) — Ne pas « chanter » :} \zh{问题} (wènti), \zh{孩子} (háizi), \zh{事情} (shìqing), \zh{办法} (bànfa), \zh{案件} (ànjiàn $\rightarrow$ ànjiàn souvent neutre 2e syllabe).
    \item \textbf{Entraînement quotidien (5 min) :}
    \begin{enumerate}
        \item Lecture exagérée « karaoké » (gestes main pour courbes de tons).
        \item Lecture vitesse normale, enregistrement.
        \item Comparaison avec modèle natif (application OnBuch+ / Pleco).
    \end{enumerate}
\end{enumerate}
\end{methode}

\courssub{Méthodes et exercices résolus}

\begin{methode}[Structurer un exposé oral simple sur la délinquance juvénile]
\textbf{Objectif :} Présenter un point de vue clair, organisé et prononcé correctement en 1 à 2 minutes (format Bac oral).
\begin{enumerate}
    \item \textbf{Accroche (开场白) :} Annoncer le thème avec un mot-clé fort.
    \textit{Exemple :} \zh{青少年犯罪问题日益严重。} (Qīngshàonián fànzuì wèntí rìyì yánzhòng.) « Le problème de la délinquance juvénile devient de plus en plus grave. »
    \item \textbf{Analyse des causes (分析原因) :} Utiliser la structure \cle{因为……所以……} (yīnwèi… suǒyǐ…) ou \cle{由于……因此……} (yóuyú… yīncǐ…). Citer 2 causes : famille, école, société, internet.
    \item \textbf{Conséquences (后果) :} Utiliser \cle{导致} (dǎozhì « entraîner ») ou \cle{影响} (yǐngxiǎng « affecter »).
    \textit{Exemple :} \zh{这导致社会不稳定。} (Zhè dǎozhì shèhuì bù wěndìng.)
    \item \textbf{Propositions / Conclusion (建议与总结) :} Introduire par \cle{建议} (jiànyì), \cle{应该} (yīnggāi), \cle{必须} (bìxū). Terminer par \cle{总之} (zǒngzhī) ou \cle{希望} (xīwàng).
    \item \textbf{Entraînement phonétique :} Répéter l'exposé 3 fois : 1. lent (tons isolés), 2. fluide (liaisons, tons sandhi), 3. vitesse naturelle avec regard vers le public.
\end{enumerate}
\end{methode}

\begin{exoresolu}[Réorganiser un exposé oral]
\textbf{Énoncé :} Les phrases ci-dessous forment un exposé mais sont dans le désordre. Remettez-les dans l'ordre logique (Accroche $\rightarrow$ Causes $\rightarrow$ Conséquences $\rightarrow$ Solutions $\rightarrow$ Conclusion). Écrivez les caractères, le pinyin et la traduction pour chaque étape.
\begin{enumerate}
    \item \zh{因此，我们必须加强家庭教育。}
    \item \zh{青少年犯罪已经成为一个热门话题。}
    \item \zh{因为很多父母忙于工作，缺乏对孩子的管教。}
    \item \zh{这导致很多年轻人走上犯罪道路。}
    \item \zh{总之，全社会都要关注青少年成长。}
\end{enumerate}

\tcblower
\textbf{Solution détaillée :}
\textbf{Ordre logique :} 2 $\rightarrow$ 3 $\rightarrow$ 4 $\rightarrow$ 1 $\rightarrow$ 5.

\begin{center}
\begin{tabularx}{\linewidth}{|X|X|X|X|}
\hline
\rowcolor{popblueL} \textbf{Ordre} & \textbf{Caractères} & \textbf{Pinyin} & \textbf{Traduction} \\
\hline
1 (Accroche) & \zh{青少年犯罪已经成为一个热门话题。} & Qīngshàonián fànzuì yǐjīng chéngwéi le yīgè rèmén huàtí. & La délinquance juvénile est devenue un sujet brûlant. \\
\hline
2 (Cause) & \zh{因为很多父母忙于工作，缺乏对孩子的管教。} & Yīnwèi hěnduō fùmǔ mángyú gōngzuò, quēfá duì háizi de guǎnjiào. & Parce que beaucoup de parents sont occupés par le travail, ils manquent de supervision sur leurs enfants. \\
\hline
3 (Conséquence) & \zh{这导致很多年轻人走上犯罪道路。} & Zhè dǎozhì hěnduō niánqīngrén zǒushàng fànzuì dàolù. & Cela entraîne beaucoup de jeunes sur la voie de la criminalité. \\
\hline
4 (Solution) & \zh{因此，我们必须加强家庭教育。} & Yīncǐ, wǒmen bìxū jiāqiáng jiātíng jiàoyù. & Par conséquent, nous devons renforcer l'éducation familiale. \\
\hline
5 (Conclusion) & \zh{总之，全社会都要关注青少年成长。} & Zǒngzhī, quán shèhuì dōu yào guānzhù qīngshàonián chéngzhǎng. & En somme, toute la société doit veiller à la croissance des adolescents. \\
\hline
\end{tabularx}
\end{center}

\textbf{Justifications grammaticales :}
\begin{itemize}
    \item Phrase 2 : Structure \cle{因为……，……} (conséquence implicite ou phrase suivante). \cle{缺乏} (quēfá) verbe transitif « manquer de ».
    \item Phrase 3 : \cle{导致} (dǎozhì) « entraîner » [Sujet = Cause, Objet = Conséquence]. \cle{走上} (zǒushàng) « s'engager sur (une voie) ».
    \item Phrase 4 : \cle{因此} (yīncǐ) « par conséquent » (adverbe de liaison en tête de proposition). \cle{加强} (jiāqiáng) « renforcer ».
    \item Phrase 5 : \cle{总之} (zǒngzhī) « en somme » pour conclure. \cle{都要} (dōu yào) « tous doivent ».
\end{itemize}
\textbf{Conseil Bac :} Apprenez ces 5 phrases comme un « squelette » universel, puis remplacez le vocabulaire spécifique (causes, solutions) selon le sujet tiré.
\end{exoresolu}

\begin{methode}[Participer à un débat : Exprimer accord, désaccord et nuance]
\textbf{Objectif :} Réagir spontanément aux arguments d'un camarade en utilisant les formules fonctionnelles officielles.
\begin{enumerate}
    \item \textbf{Exprimer l'accord (赞同) :} \cle{我赞同你的看法。} (Wǒ zàntóng nǐ de kànfǎ.) / \cle{你说得很对。} (Nǐ shuō de hěn duì.) / \cle{完全正确。} (Wánquán zhèngquè.)
    \item \textbf{Exprimer le désaccord poli (不同意/委婉) :} \cle{我不同意。} (Wǒ bù tóngyì — trop direct.) $\rightarrow$ Préférer : \cle{我不太同意……} (Wǒ bù tài tóngyì…) / \cle{恐怕……} (Kǒngpà… « Je crains que… ») / \cle{这倒不一定。} (Zhè dào bù yídìng. « Ce n'est pas forcément vrai. »)
    \item \textbf{Nuancer / Apporter un contre-argument (转折) :} \cle{一方面……，另一方面……} (Yīfāngmiàn… lìngyī fāngmiàn… « D'un côté… de l'autre… ») / \cle{虽然……，但是……} (Suīrán…, dànshì… « Bien que… cependant… »).
    \item \textbf{Relancer / Demander précision (追问) :} \cle{你为什么这么认为？} (Nǐ wèishénme zhème rènwéi?) / \cle{有什么证据吗？} (Yǒu shénme zhèngjù ma?) / \cle{能举个例子吗？} (Néng jǔ gè lìzi ma?)
    \item \textbf{Stratégie de survie :} Si vous séchez, utilisez \cle{让我想想……} (Ràng wǒ xiǎngxiǎng… « Laissez-moi réfléchir… ») pour gagner du temps sans perdre la parole.
\end{enumerate}
\end{methode}

\begin{exoresolu}[Compléter un débat : Punir ou éduquer ?]
\textbf{Énoncé :} Complétez le dialogue suivant en choisissant la formule fonctionnelle la plus adaptée parmi la boîte à outils. Écrivez la phrase complète en caractères, pinyin et traduction.
\textbf{Boîte à outils :} \zh{我不太同意} | \zh{另一方面} | \zh{能举个例子吗} | \zh{你说得很对} | \zh{恐怕}

\begin{textebox}[Dialogue à trous]
A : \zh{我认为应该严厉惩罚青少年罪犯，这样才能震慑别人。} (Wǒ rènwéi yīnggāi yánlì chéngfá qīngshàonián zuìfàn, zhèyàng cáinéng zhèndé biéren.)
B : \underline{\hspace{5cm}} \zh{，教育比惩罚更重要。} (…, jiàoyù bǐ chéngfá gèng zhòngyào.)
A : \underline{\hspace{5cm}} \zh{？} (…?)
B : \zh{比如挪威的监狱模式，重教化不重刑罚，再犯率很低。} (Bǐrú Nuówēi de jiānyù móshì, zhòng jiàohuà bù zhòng xíngfá, zàifàn lǜ hěn dī.)
A : \underline{\hspace{5cm}} \zh{。} (….) \zh{但是中国国情不同。} (Dànshì Zhōngguó guóqíng bùtóng.)
B : \underline{\hspace{5cm}} \zh{，我们可以借鉴理念。} (…, wǒmen kěyǐ jièjiàn lǐniàn.)
A : \underline{\hspace{5cm}} \zh{。} (….) \zh{这是一个很好的切入点。} (Zhè shì yīgè hěn hǎo de qiērùdiǎn.)
\end{textebox}

\tcblower
\textbf{Solution détaillée :}

\begin{center}
\begin{tabularx}{\linewidth}{|X|X|X|X|}
\hline
\rowcolor{popblueL} \textbf{Tour} & \textbf{Formule choisie} & \textbf{Phrase complète (Caractères + Pinyin)} & \textbf{Traduction} \\
\hline
B (1) & \zh{我不太同意} & \zh{我不太同意，教育比惩罚更重要。} (Wǒ bù tài tóngyì, jiàoyù bǐ chéngfá gèng zhòngyào.) & Je ne suis pas tout à fait d'accord, l'éducation est plus importante que la punition. \\
\hline
A (2) & \zh{能举个例子吗} & \zh{能举个例子吗？} (Néng jǔ gè lìzi ma?) & Pouvez-vous donner un exemple ? \\
\hline
B (3) & \zh{你说得很对} & \zh{你说得很对。} (Nǐ shuō de hěn duì.) & Tu as tout à fait raison. \\
\hline
A (4) & \zh{另一方面} & \zh{另一方面，我们可以借鉴理念。} (Lìngyī fāngmiàn, wǒmen kěyǐ jièjiàn lǐniàn.) & D'un autre côté, nous pouvons nous inspirer des concepts. \\
\hline
B (5) & \zh{恐怕} & \zh{恐怕这是一个很好的切入点。} (Kǒngpà zhè shì yīgè hěn hǎo de qiērùdiǎn.) & Je reconnais / Il se trouve que c'est un excellent point d'entrée. (\zh{恐怕} adoucit l'aveu). \\
\hline
\end{tabularx}
\end{center}

\textbf{Analyse des choix :}
\begin{itemize}
    \item \textbf{B(1) :} \zh{我不太同意} + virgule + argument contraire. Structure standard du désaccord constructif.
    \item \textbf{A(2) :} \zh{能举个例子吗？} Formule de relance pour demander une preuve (compétence « relancer »).
    \item \textbf{A(4) :} \zh{你说得很对} (accord) + \zh{但是} (transition) + \zh{另一方面} (introduction du contre-point). Maîtrise de la concession.
    \item \textbf{B(5) :} \zh{恐怕} (kǒngpà) ici ne signifie pas « j'ai peur » mais « il se peut que / je reconnais que », marqueur d'accord réticent ou de prise de conscience.
\end{itemize}
\textbf{Point de prononciation :} \zh{不太} (bù tài) $\rightarrow$ \textbf{bú tài} (règle du 3e ton + 4e ton : \cle{不} devient 2e ton devant 4e ton).
\end{exoresolu}

\begin{methode}[Correction de tons : Dictée tonale et discrimination]
\textbf{Objectif :} Valider l'application des sandhi et la distinction des quasi-homophones.
\begin{enumerate}
    \item Écrivez le pinyin \textbf{avec les tons corrects (après sandhi)} pour : \zh{青少年}, \zh{犯罪}, \zh{管教}, \zh{不稳定}, \zh{一律}, \zh{矫正}.
    \item Choisissez le bon caractère pour compléter la phrase (attention aux tons/sens) :
    \zh{我们要\_\_\_\_\_\_(管教/官教)青少年，不能\_\_\_\_\_\_(放任/放人)。}
    \zh{法律面前人人平等，\_\_\_\_\_\_(一律/依律)平等。}
\end{enumerate}
\end{methode}

\begin{exoresolu}[Correction de tons : Dictée tonale et discrimination]
\tcblower
\textbf{Solution détaillée :}

\textbf{1. Pinyin corrigé (tons réels prononcés) :}
\begin{center}
\begin{tabularx}{\linewidth}{|X|X|X|}
\hline
\rowcolor{popblueL} \textbf{Mot (Caractères)} & \textbf{Pinyin dictionnaire} & \textbf{Pinyin réel (Sandhi appliqué)} \\
\hline
\zh{青少年} & qīng shǎo nián (1-3-2) & \textbf{qíng shǎo nián} (2-3-2) \\
\zh{犯罪} & fàn zuì (4-4) & \textbf{fàn zuì} (4-4) \textit{(pas de sandhi 3+3 standard)} \\
\zh{管教} & guǎn jiào (3-4) & \textbf{guán jiào} (2-4) \\
\zh{不稳定} & bù wěn dìng (4-3-4) & \textbf{bú wěn dìng} (2-3-4) \\
\zh{一律} & yī lǜ (1-4) & \textbf{yì lǜ} (4-4) \\
\zh{矫正} & jiǎo zhèng (3-4) & \textbf{jiáo zhèng} (2-4) \\
\hline
\end{tabularx}
\end{center}
\textit{Rappel règle :} 3e ton + 3e ton $\rightarrow$ 2e ton + 3e ton. \cle{不} (4e) $\rightarrow$ 2e ton devant 4e ton. \cle{一} (1e) $\rightarrow$ 4e ton devant 1e, 2e, 3e ton ; 2e ton devant 4e ton.

\textbf{2. Choix des caractères (Homophones / Proches) :}
\begin{itemize}
    \item \zh{我们要\textbf{管教}(guǎnjiào)青少年，不能\textbf{放任}(fàngrèn)。}
    \textit{Traduction :} Nous devons éduquer/discipliner les jeunes, ne pas les laisser faire / négliger.
    \textit{Piège :} \zh{官教} (guānjiào) n'existe pas comme verbe courant. \zh{放人} (fàngrén) = « libérer qqn (de prison) », hors contexte.
    \item \zh{法律面前人人平等，\textbf{一律}(yīlǜ)平等。}
    \textit{Traduction :} Devant la loi, tous sont égaux, uniformément égaux.
    \textit{Piège :} \zh{依律} (yīlǜ) « selon la loi » (verbe + objet), ne fonctionne pas comme adverbe « uniformément » ici. \zh{一律} (yīlǜ) = « sans exception ».
\end{itemize}
\textbf{Conclusion :} La confusion \zh{管教}/\zh{官教} ou \zh{一律}/\zh{依律} coûte cher en compréhension orale et écrite. Travaillez les \textbf{paires minimales} par séries de 5 répétitions espacées.
\end{exoresolu}

\begin{methode}[Rédiger un paragraphe argumenté (Support écrit de l'oral / Épreuve écrite Bac)]
\textbf{Objectif :} Produire un texte court (80-100 caractères) cohérent, servant de notes pour l'exposé ou expression écrite.
\begin{enumerate}
    \item \textbf{Thèse claire (观点) :} Une phrase avec \cle{我认为} (wǒ rènwéi) ou \cle{在我看来} (zài wǒ kànlái). \textit{Ex :} \zh{我认为预防青少年犯罪关键在家庭。}
    \item \textbf{Argument 1 : Cause racine (原因) :} \cle{首先} (shǒuxiān) + \cle{因为……所以……}. \textit{Ex :} \zh{首先，缺乏父母陪伴导致心理扭曲。}
    \item \textbf{Argument 2 : Solution concrète (对策) :} \cle{其次} (qícì) + \cle{应该/建议} + Verbe + Objet. \textit{Ex :} \zh{其次，学校应开设法治教育课。}
    \item \textbf{Comparaison / Ouverture (对比/借鉴) :} \cle{比如} (bǐrú) + Exemple (Cameroun/Chine). \textit{Ex :} \zh{比如喀麦隆的社区足球项目就有效减少了街头青年。}
    \item \textbf{Conclusion percutante (总结) :} \cle{只有……才能……} (zhǐyǒu… cáinéng… « Ce n'est que si… qu'on peut… »). \textit{Ex :} \zh{只有家校社联动，才能守护未来。}
    \item \textbf{Vérification caractères :} Contrôlez les \cle{的/得/地}, les classificateurs (\cle{个/起/宗} cas), l'ordre [Sujet - Temps - Lieu - Manière - Verbe].
\end{enumerate}
\end{methode}

\begin{exoresolu}[Production écrite : « Prévenir plutôt que guérir »]
\textbf{Énoncé :} Rédigez un paragraphe d'environ 90 caractères en chinois simplifié sur le thème : \zh{如何预防青少年犯罪？} (Comment prévenir la délinquance juvénile ?). Utilisez au moins 5 mots du vocabulaire de la leçon et 3 connecteurs logiques. Présentez : Caractères | Pinyin | Traduction.

\tcblower
\textbf{Solution détaillée :}

\textbf{Proposition de rédaction (Modèle « Bac ») :}

\begin{textebox}[Rédaction modèle]
\zh{我认为预防青少年犯罪，关键在于家庭、学校与社会联动。} \\
\zh{首先，父母应加强管教，多陪伴孩子，预防心理扭曲。} \\
\zh{其次，学校必须开设法治教育，增强青少年法制观念。} \\
\zh{再者，社区可组织文体活动，如足球赛，减少不良诱惑。} \\
\zh{只有全社会共同努力，才能守护祖国的花朵健康成长。}
\end{textebox}

\begin{center}
\begin{tabularx}{\linewidth}{|X|X|X|}
\hline
\rowcolor{popblueL} \textbf{Caractères} & \textbf{Pinyin} & \textbf{Traduction} \\
\hline
\zh{我认为预防青少年犯罪，关键在于家庭、学校与社会联动。} & Wǒ rènwéi yùfáng qīngshàonián fànzuì, guānjiàn zài yú jiātíng, xuéxiào yǔ shèhuì liándòng. & Je pense que pour prévenir la délinquance juvénile, la clé réside dans la synergie famille-école-société. \\
\hline
\zh{首先，父母应加强管教，多陪伴孩子，预防心理扭曲。} & Shǒuxiān, fùmǔ yīng jiāqiáng guǎnjiào, duō péibàn háizi, yùfáng xīnlǐ niǔqū. & Premièrement, les parents doivent renforcer la discipline, accompagner davantage les enfants, prévenir la déviation psychologique. \\
\hline
\zh{其次，学校必须开设法治教育，增强青少年法制观念。} & Qícì, xuéxiào bìxū kāishè fǎzhì jiàoyù, zēngqiáng qīngshàonián fǎzhì guānniàn. & Deuxièmement, les écoles doivent ouvrir des cours d'éducation légale, renforcer la conscience juridique des jeunes. \\
\hline
\zh{再者，社区可组织文体活动，如足球赛，减少不良诱惑。} & Zài'ěr, shèqū kě zǔzhī wéntǐ huódòng, rú zúqiúsài, jiǎnshǎo bùliáng yòuhuò. & En outre, la communauté peut organiser des activités culturelles/sportives, ex: match de foot, réduire les mauvaises tentations. \\
\hline
\zh{只有全社会共同努力，才能守护祖国的花朵健康成长。} & Zhǐyǒu quán shèhuì gòngtóng nǔlì, cáinéng shǒuhù zǔguó de huāduǒ jiànkāng chéngzhǎng. & Ce n'est que si toute la société s'unit qu'on peut protéger les fleurs de la patrie pour qu'elles grandissent sainement. \\
\hline
\end{tabularx}
\end{center}

\textbf{Analyse structurelle (Grille d'évaluation) :}
\begin{itemize}
    \item \textbf{Vocabulaire cible (6) :} \zh{预防} (yùfáng), \zh{青少年犯罪} (qīngshàonián fànzuì), \zh{管教} (guǎnjiào), \zh{法治教育} (fǎzhì jiàoyù), \zh{不良诱惑} (bùliáng yòuhuò), \zh{联动} (liándòng).
    \item \textbf{Connecteurs (4) :} \zh{首先} (shǒuxiān), \zh{其次} (qícì), \zh{再者} (zài'ěr), \zh{只有……才能……} (zhǐyǒu… cáinéng…).
    \item \textbf{Grammaire :} \zh{应/必须/可} (modaux), \zh{加强/开设/组织/增强/减少/守护} (verbes d'action précis), \zh{在于} (zài yú « résider dans »).
    \item \textbf{Socioculturel :} Référence au football (\zh{足球赛}) et à l'image des « fleurs de la patrie » (\zh{祖国的花朵}) très usitée en Chine pour les jeunes.
\end{itemize}
\end{exoresolu}

\courssub{Jeux de rôle et débat guidé : Simulation « Le Tribunal des Jeunes »}
\begin{methode}[Animer un jeu de rôle : « Le Tribunal des Jeunes » (Simulation interactive)]
\textbf{Objectif :} Mettre en interaction 3-4 élèves dans une situation communicationnelle authentique (jugement, médiation, débat TV) pour évaluer l'interaction orale (compétence « S'exprimer oralement en interaction »).
\begin{enumerate}
    \item \textbf{Distribution des rôles (角色分配) :}
    \begin{itemize}
        \item Rôle A : \zh{法官/主持人} (Fǎguān / Zhǔchírén) — Structure le débat, pose les questions, synthétise, rend le verdict.
        \item Rôle B : \zh{检察官/家长代表} (Jiǎncháguān / Jiāzhǎng dàibiǎo) — Accusation / Exige sévérité / Protection victimes.
        \item Rôle C : \zh{辩护律师/心理专家} (Biànhù lǜshī / Xīnlǐ zhuānjiā) — Défense / Plaide l'éducation, le contexte, la réinsertion.
        \item Rôle D (optionnel) : \zh{青少年/当事人} (Qīngshàonián / Dāngshìrén) — S'excuse, explique son acte, promet de changer.
    \end{itemize}
    \item \textbf{Préparation silencieuse (5 min) :} Chaque élève note 3 arguments clés et 2 formules de relance sur une fiche (caractères + pinyin). \textbf{Interdit d'écrire le texte complet.}
    \item \textbf{Déroulement (8-10 min) :}
    \begin{enumerate}
        \item \textbf{Ouverture (Juge) :} \zh{今天我们讨论一起青少年盗窃案。检察官请发言。}
        \item \textbf{Confrontation (Procureur vs Avocat) :} Échanges vifs. Le Juge coupe si hors-sujet : \zh{请围绕“惩罚还是教育”发言。} (Qǐng wéirao “chéngfá háishì jiàoyù” fāyán.)
        \item \textbf{Témoignage (Jeune - si présent) :} \zh{我知错了，我以后一定好好学习。} (Wǒ zhīcuò le, wǒ yǐhòu yídìng hǎohǎo xuéxí.)
        \item \textbf{Verdict / Synthèse (Juge) :} \zh{综合双方意见，建议……} (Zōnghé shuāngfāng yìjiàn, jiànyì…)
    \end{enumerate}
    \item \textbf{Évaluation par les pairs (Grille simplifiée) :}
    \begin{center}
    \begin{tabularx}{\linewidth}{|X|X|}
    \hline
    \rowcolor{popblueL} \textbf{Critère} & \textbf{Indicateurs (Oui / En partie / Non)} \\
    \hline
    Prononciation (Tons) & Tons de base corrects, sandhi 3+3, \zh{不}/\zh{一} appliqués. \\
    Fluidité & Débit naturel, peu de pauses longues, connecteurs utilisés. \\
    Vocabulaire précis & Emploi de \zh{管教, 矫正, 震慑, 再犯, 法治, 联动} etc. \\
    Interaction (Relance) & Utilise \zh{能举个例子吗？}, \zh{你为什么这么认为？}, \zh{恐怕……} \\
    Respect du tour de parole & Écoute active, ne coupe pas, réagit à l'argument précédent. \\
    \hline
    \end{tabularx}
    \end{center}
\end{enumerate}
\end{methode}

\begin{exoresolu}[Simulation : Plaidoirie « Éduquer ou Punir »]
\textbf{Énoncé :} Vous êtes l'Avocat de la défense (Rôle C). Le Procureur vient de dire : \zh{必须重判，否则无法震慑！} (Il faut une peine lourde, sinon pas de dissuasion !). Répondez en 3 phrases en utilisant : 1. Une concession (\zh{虽然……但是……}), 2. Un argument factuel (\zh{据统计……}), 3. Une proposition alternative (\zh{建议……}). Écrivez votre intervention complète.

\tcblower
\textbf{Solution détaillée :}

\textbf{Intervention de l'Avocat (Rôle C) :}

\begin{textebox}[Plaidoirie]
\zh{虽然震慑很重要，但是重罚不等于预防再犯。} \\
\zh{据统计，接受矫正教育的青少年再犯率远低于单纯坐牢者。} \\
\zh{建议法庭判处社区矫正，配合心理疏导，给孩子改过自新的机会。}
\end{textebox}

\begin{center}
\begin{tabularx}{\linewidth}{|X|X|X|}
\hline
\rowcolor{popblueL} \textbf{Caractères} & \textbf{Pinyin} & \textbf{Traduction} \\
\hline
\zh{虽然震慑很重要，但是重罚不等于预防再犯。} & Suīrán zhèndé hěn zhòngyào, dànshì zhòngfá bù děngyú yùfáng zàifàn. & Bien que la dissuasion soit importante, une peine lourde n'équivaut pas à prévenir la récidive. \\
\hline
\zh{据统计，接受矫正教育的青少年再犯率远低于单纯坐牢者。} & Jù tǒngjì, jiēshòu jiǎozhèng jiàoyù de qīngshàonián zàifàn lǜ yuǎn dīyú dānchún zuòláo zhě. & Selon les statistiques, le taux de récidive des jeunes ayant reçu une éducation correctionnelle est bien inférieur à celui de ceux qui font juste de la prison. \\
\hline
\zh{建议法庭判处社区矫正，配合心理疏导，给孩子改过自新的机会。} & Jiànyì fǎtíng pànchǔ shèqū jiǎozhèng, pèihé xīnlǐ shūdǎo, gěi háizi gǎiguòzìxīn de jīhuì. & Je suggère au tribunal de prononcer une correctionnelle communautaire, avec un soutien psychologique, pour donner à l'enfant une chance de se corriger. \\
\hline
\end{tabularx}
\end{center}

\textbf{Points de grammaire validés (Niveau HSK 4 / Bac) :}
\begin{itemize}
    \item \zh{虽然……但是……} (Structure de concession incontournable).
    \item \zh{据统计} (Jù tǒngjì) « Selon les stats » — marqueur d'argument d'autorité (oral formel).
    \item \zh{远低于} (yuǎn dī yú) « bien inférieur à » — comparatif supérieur.
    \item \zh{判处} (pànchǔ) « condamner à (peine) » — verbe juridique précis.
    \item \zh{配合} (pèihé) « en combinaison avec / assorti de » — structure V + O comme complément de manière.
    \item \zh{改过自新} (gǎiguòzìxīn) — Chengyu (成语) « corriger ses fautes et recommencer à neuf », très valorisé à l'oral.
\end{itemize}
\textbf{Conseil prononciation :} \zh{重罚} (zhòngfá) vs \zh{种法} (zhǒngfǎ). \zh{再犯} (zàifàn 4-4 « récidiver ») vs \zh{在饭} (zài fàn « pendant le repas »). Articulez les finales \zh{-an} (\zh{统计} tǒngjì) et \zh{-ang} (\zh{矫正} jiǎozhèng) distinctement.
\end{exoresolu}

\courssub{Élément socioculturel : Cameroun \& Chine}
\begin{savaistu}[Justice des mineurs : Approches comparées]
\textbf{Cameroun (Contexte local) :}
\begin{itemize}
    \item \textbf{Cadre légal :} Loi n°2016/007 du 12 juillet 2016 relative à la protection de l'enfant. Âge de la responsabilité pénale : 10 ans (non imputabilité absolue < 10 ans, atténuée 10-18 ans).
    \item \textbf{Structures :} Centres de rééducation (ex: Centre de Rééducation de Betamba), Tribunaux pour enfants. Mesures : réprimande, remise à parents, liberté surveillée, placement en centre.
    \item \textbf{Réalité terrain :} Surpopulation carcérale, manque d'éducateurs spécialisés, rôle crucial des ONG et chefs traditionnels (médiation coutumière). Le football (\zh{足球}) est un levier majeur d'insertion (programmes FECAFOOT / MINJEC).
\end{itemize}
\textbf{Chine (Référence HSK) :}
\begin{itemize}
    \item \textbf{Cadre légal :} Loi sur la protection des mineurs (révisée 2021), Code pénal (Amendement XI, 2020 : responsabilité pénale abaissée à 12 ans pour crimes graves \zh{故意杀人, 故意伤害致死} sur décision du Parquet suprême).
    \item \textbf{Philosophie :} \zh{教育、感化、挽救} (Éduquer, transformer par l'influence morale, sauver). Priorité à la \zh{社区矫正} (correction communautaire) vs incarcération.
    \item \textbf{École :} Cours obligatoires de \zh{法治教育} (éducation à l'État de droit) dès le primaire. \zh{法治副校长} (Vice-principal chargé du droit) : policiers/procureurs nommés dans les écoles.
    \item \textbf{Vocabulaire spécifique :} \zh{不良行为} (comportements déviants, < crime), \zh{严重不良行为} (comportements gravement déviants), \zh{专门学校} (écoles spécialisées de correction).
\end{itemize}
\textbf{Point de convergence :} Les deux pays évoluent vers une \textbf{justice restaurative} (réparation du lien social) plutôt que purement punitive, avec l'implication de la \zh{社区} (communauté/quartier).
\end{savaistu}

\courssub{Exercices d'entraînement (Travail personnel / Devoirs)}
\begin{exercice}[Exercice 1 : Vocabulaire \& Tons — Compléter le tableau]
Complétez le tableau ci-dessous pour les mots du thème. Attention aux tons sandhi à l'oral (colonne 3).
\begin{center}
\begin{tabularx}{\linewidth}{|X|X|X|X|}
\hline
\rowcolor{popblueL} \textbf{Caractères} & \textbf{Pinyin dictionnaire} & \textbf{Pinyin oral (Sandhi)} & \textbf{Traduction} \\
\hline
\zh{预防} & & & \\
\zh{心理扭曲} & & & \\
\zh{陪伴} & & & \\
\zh{判处} & & & \\
\zh{改过自新} & & & \\
\hline
\end{tabularx}
\end{center}
\end{exercice}

\begin{exercice}[Exercice 2 : Grammaire — Choisir le bon mot-outil / classificateur]
Remplissez les vides avec \zh{的 / 得 / 地} ou le bon classificateur (\zh{起 / 宗 / 个 / 位}).
\begin{enumerate}
    \item 这是一\_\_\_\_\_\_严重的青少年犯罪案件。
    \item 法官判处\_\_\_\_\_\_三年有期徒刑。
    \item 我们要严厉\_\_\_\_\_\_惩罚犯罪分子。
    \item 这位\_\_\_\_\_\_律师辩护得很有力。
    \item 孩子们快乐\_\_\_\_\_\_地踢足球。
\end{enumerate}
\end{exercice}

\begin{exercice}[Exercice 3 : Traduction — Thème (Français $\rightarrow$ Chinois)]
Traduisez en chinois (caractères + pinyin). Utilisez le vocabulaire de la leçon.
\begin{enumerate}
    \item Prévenir la délinquance juvénile nécessite la coopération de toute la société.
    \item Bien que la peine soit lourde, elle ne peut pas résoudre le problème fondamental.
    \item Selon les statistiques, le taux de récidive des jeunes ayant reçu une éducation correctionnelle est plus bas.
    \item Les parents devraient passer plus de temps avec leurs enfants pour prévenir les déviances psychologiques.
\end{enumerate}
\end{exercice}

\begin{exercice}[Exercice 4 : Expression orale — Préparation débat]
\textbf{Sujet :} \zh{“降低刑事责任年龄能否有效遏制青少年犯罪？”} (Abaisser l'âge de responsabilité pénale peut-il endiguer efficacement la délinquance juvénile ?).
\textbf{Consigne :} Préparez une fiche d'arguments (face A : POUR / face B : CONTRE) avec 3 arguments par face. Chaque argument = 1 phrase chinoise (caractères + pinyin) + 1 connecteur logique. Utilisez au moins 3 mots du vocabulaire par face. Durée de l'exposé visé : 1 min 30.
\end{exercice}

\courssub{Corrigés des exercices}
\begin{corrige}[Exercice 1]
\begin{center}
\begin{tabularx}{\linewidth}{|X|X|X|X|}
\hline
\rowcolor{popblueL} \textbf{Caractères} & \textbf{Pinyin dictionnaire} & \textbf{Pinyin oral (Sandhi)} & \textbf{Traduction} \\
\hline
\zh{预防} & yù fáng (4-2) & yù fáng (4-2) & prévenir \\
\zh{心理扭曲} & xīn lǐ niǔ qū (1-3-3-1) & xīn lǐ niú qū (1-3-2-1) & déviation psychologique \\
\zh{陪伴} & péi bàn (2-4) & péi bàn (2-4) & accompagner \\
\zh{判处} & pàn chǔ (4-3) & pàn chǔ (4-3) & condamner à \\
\zh{改过自新} & gǎi guò zì xīn (3-4-4-1) & gǎi guò zì xīn (3-4-4-1) & se corriger / tourner la page \\
\hline
\end{tabularx}
\end{center}
\textit{Note :} \zh{扭曲} niǔqū (3-1) $\rightarrow$ niúqū (2-1) par sandhi 3+3 sur \zh{扭} si isolé, mais ici \zh{心理} (2-3) précède, donc \zh{扭} reste 3e ton standard ou devient 2e selon débit. Standard enseigné : \textbf{niǔqū} (3-1) ou \textbf{niúqū} (2-1). Les deux acceptés à l'oral.
\end{corrige}

\begin{corrige}[Exercice 2]
\begin{enumerate}
    \item \zh{起} (classif. pour affaire/cas juridique). \textit{Phrase :} 这是一\textbf{起}严重的青少年犯罪案件。
    \item \zh{的} (complément de nombre/quantité après verbe \zh{判处} + durée). \textit{Structure :} V + \textbf{的} + Durée + Objet. \textit{Phrase :} 法官判处\textbf{的}三年有期徒刑。 (Ou \zh{判处三年有期徒刑} sans \zh{的} si pas d'attribut sur la durée. Ici \zh{三年} attribut de \zh{徒刑}, pas de \zh{的} nécessaire. \textbf{Piège !} Correction : \textbf{Ø} (zéro). \zh{法官判处三年有期徒刑。} Si phrase : \zh{判处的这三年} alors \zh{的}. \textbf{Réponse attendue :} \textbf{Ø} (rien) ou \zh{了} (aspect accompli). \textit{Reformulation exercice :} "Le juge A CONDAMNÉ..." $\rightarrow$ \zh{法官判处\textbf{了}三年...}).
    \item \zh{地} (adverbe \zh{严厉地} modifiant \zh{惩罚}). \textit{Phrase :} 我们要严厉\textbf{地}惩罚犯罪分子。
    \item \zh{位} (classif. honorifique pour personne / professionnel). \textit{Phrase :} 这\textbf{位}律师辩护得很有力。
    \item \zh{地} (adverbe \zh{快乐地} modifiant \zh{踢}). \textit{Phrase :} 孩子们快乐\textbf{地}踢足球。
\end{enumerate}
\end{corrige}

\begin{corrige}[Exercice 3]
\begin{enumerate}
    \item \zh{预防青少年犯罪需要全社会联动。} (Yùfáng qīngshàonián fànzuì xūyào quán shèhuì liándòng.)
    \item \zh{虽然刑罚很重，但是不能解决根本问题。} (Suīrán xíngfá hěn zhòng, dànshì bùnéng jiějué gēnběn wèntí.)
    \item \zh{据统计，接受矫正教育的青少年再犯率比较低。} (Jù tǒngjì, jiēshòu jiǎozhèng jiàoyù de qīngshàonián zàifàn lǜ bǐjiào dī.)
    \item \zh{父母应该多陪伴孩子，预防心理扭曲。} (Fùmǔ yīnggāi duō péibàn háizi, yùfáng xīnlǐ niǔqū.)
\end{enumerate}
\end{corrige}

\begin{corrige}[Exercice 4 — Pistes de correction (Non exhaustives)]
\textbf{Face A : POUR (降低年龄有效)}
\begin{itemize}
    \item \zh{首先，降低责任年龄能震慑恶性犯罪青少年。} (Shǒuxiān, jiàngdī zérèn niánlíng néng zhèndé èxìng fànzuì qīngshàonián.)
    \item \zh{其次，严重暴力案件必须由刑法制裁，体现公平正义。} (Qícì, yánzhòng bàolì ànjiàn bìxū yóu xíngfǎ zhìcái, tǐxiàn gōngpíng zhèngyì.)
    \item \zh{再者，这倒逼家庭加强管教，落实监护责任。} (Zài'ěr, zhè dǎobī jiātíng jiāqiáng guǎnjiào, luòshí jiānhù zérèn.)
\end{itemize}
\textbf{Face B : CONTRE (降低年龄无效/有害)}
\begin{itemize}
    \item \zh{首先，刑罚不等于矫正，监狱反而可能成为犯罪学校。} (Shǒuxiān, xíngfá bù děngyú jiǎozhèng, jiānyù fǎn'ér kěnéng chéngwéi fànzuì xuéxiào.)
    \item \zh{其次，青少年心智未成熟，应以教育挽救为主，而非报复。} (Qícì, qīngshàonián xīnzhì wèi chéngshú, yǐ yǐ jiàoyù wǎnjiù wéi zhǔ, érfēi bàofù.)
    \item \zh{再者，国际趋势是完善社区矫正和心理疏导体系。} (Zài'ěr, guójì qūshì shì wánshàn shèqū jiǎozhèng hé xīnlǐ shūdǎo tǐxì.)
\end{itemize}
\end{corrige}

\begin{attention}[Les 5 erreurs qui coûtent des points au Bac (Oral \& Écrit)]
\begin{enumerate}
    \item \textbf{Tons : Confusion \zh{青} (qīng 1) / \zh{轻} (qīng 1) / \zh{情} (qíng 2) / \zh{请} (qǐng 3).}
    \textit{Piège :} \zh{青少年} (qīngshàonián) vs \zh{轻少年} (inexistant) vs \zh{情少年} (inexistant). \textbf{Règle :} \zh{青} = 1er ton (haut plat). \zh{请} = 3e ton (bas descendant-remontant). En sandhi \zh{青少年} $\rightarrow$ \textbf{qíngshàonián}. Ne dites pas \zh{qīngshàonián} (plat-plat) ni \zh{qǐngshàonián}.
    
    \item \textbf{Ordre des mots : Complément de lieu/temps APRÈS le verbe (calque français).}
    \textit{Faux :} \zh{我在家庭建议加强教育。} (Je *à la famille* suggère...).
    \textit{Vrai :} \zh{我建议在家庭加强教育。} (Sujet + Verbe + \textbf{在/从/到} Lieu + Verbe/Objet). \textbf{Règle :} En chinois, le circonstanciel de lieu (où se passe l'action) précède le verbe (ou suit \zh{在}), jamais en début de phrase comme en français sauf topicalisation marquée.
    
    \item \textbf{Mots-outils : \zh{的} (de) vs \zh{得} (de) vs \zh{地} (de).}
    \textit{Contexte leçon :} \zh{严厉\textbf{地}惩罚} (sévèrement punir - adverbe \zh{地} + V), \zh{严厉\textbf{的}惩罚} (la punition sévère - attribut \zh{的} + N), \zh{惩罚\textbf{得}严厉} (punir sévèrement - complément de degré V + \zh{得} + Adj).
    \textit{Erreur type :} \zh{我们应该严厉的惩罚他们。} $\rightarrow$ Corrigez : \zh{严厉\textbf{地}惩罚} ou \zh{惩罚\textbf{得}严厉}.
    
    \item \textbf{Classificateurs : \zh{个} (gè) abusif pour « cas / affaire / incident ».}
    \textit{Correct :} \zh{一\textbf{起}青少年犯罪案件} (yī \textbf{qǐ}... — \zh{起} pour affaires/cas juridiques), \zh{一\textbf{宗}罪案} (yī \textbf{zōng}... — \zh{宗} pour dossiers/crimes graves). \zh{一个犯罪} est incorrect (sonne comme « un crime » objet concret).
    
    \item \textbf{Voix passive / Disposition : \zh{被} (bèi) vs \zh{把} (bǎ) vs \zh{让} (ràng).}
    \textit{Contexte :} « Le jeune a été éduqué par ses parents. »
    \textit{Faux :} \zh{青少年被父母教育了。} (Passif \zh{被} marquant souvent la souffrance/subie — nuance négative).
    \textit{Mieux (neutre/factuel, action aboutie) :} \zh{父母\textbf{把}青少年\textbf{教育}好了。} (Disposition \zh{把} : action aboutie sur l'objet + complément de résultat \zh{好}).
    \textit{Oral courant (causatif) :} \zh{父母\textbf{让}青少年\textbf{受教育}了。} (\zh{让} causatif).
    \textbf{Attention :} Au Bac, \zh{被} implique souvent une nuance négative (volé, blessé, puni). Pour « éduquer / corriger / aider », privilégiez \zh{把} + V + 好/完 ou structure active.
\end{enumerate}
\end{attention}

\newpage

\coursec{Exercices}

\courssub{Je vérifie mes connaissances}

\begin{exercice}[QCM : vocabulaire et structures]
Chaque item ne comporte qu'une seule réponse correcte. Choisis la bonne proposition.
\begin{enumerate}
    \item Comment dit-on « délinquance juvénile » en chinois ?
    \begin{enumerate}
        \item 青少年犯罪
        \item 青少年教育
        \item 青少年问题
        \item 青少年心理
    \end{enumerate}
    \item Quelle structure permet d'exprimer « Je pense que… / À mon avis… » de façon formelle ?
    \begin{enumerate}
        \item 我觉得
        \item 我认为
        \item 我猜
        \item 我希望
    \end{enumerate}
    \item Pour relancer le débat et demander l'avis de l'autre, on utilise :
    \begin{enumerate}
        \item 你好吗？
        \item 你觉得呢？
        \item 你是谁？
        \item 你做什么？
    \end{enumerate}
    \item La phrase « 小偷被警察抓了 » (Le voleur a été attrapé par la police) illustre la voix passive avec :
    \begin{enumerate}
        \item 把
        \item 被
        \item 让
        \item 给
    \end{enumerate}
    \item Quel connecteur logique signifie « enfin / en dernier lieu » dans un exposé structuré ?
    \begin{enumerate}
        \item 首先
        \item 然后
        \item 最后
        \item 另外
    \end{enumerate}
\end{enumerate}
\end{exercice}

\begin{exercice}[Vrai ou Faux justifié]
Lis le court texte ci-dessous, puis détermine si les affirmations sont vraies (V) ou fausses (F). Justifie ta réponse en citant un extrait du texte (caractères + pinyin).
\begin{textebox}[Mini-texte : Le rôle de la famille]
预防青少年犯罪，家庭教育是关键。父母应该多陪伴孩子，少打骂。如果家庭和睦，孩子就不容易走上歪路。学校和社会也要配合，共同解决问题。
(Yùfáng qīngshàonián fànzuì, jiātíng jiàoyù shì guānjiàn. Fùmǔ yīnggāi duō péibàn háizi, shǎo dǎmà. Rúguǒ jiātíng hémù, háizi jiù bù róngyì zǒushàng wāilù. Xuéxiào hé shèhuì yě yào pèihé, gòngtóng jiějué wèntí.)
\par\emph{Prévenir la délinquance juvénile, l'éducation familiale est la clé. Les parents devraient plus accompagner les enfants, moins les battre et gronder. Si la famille est harmonieuse, les enfants n'empruntent pas facilement la mauvaise voie. L'école et la société doivent aussi coopérer pour résoudre le problème ensemble.}
\end{textebox}
\begin{enumerate}
    \item L'éducation familiale est le facteur principal pour prévenir la délinquance.
    \item Le texte encourage les punitions corporelles (打骂) pour corriger l'enfant.
    \item L'école et la société n'ont pas de rôle à jouer selon le texte.
    \item Un climat familial harmonieux (家庭和睦) protège l'enfant.
\end{enumerate}
\end{exercice}

\begin{exercice}[Vocabulaire : associer et compléter]
Complète le tableau ci-dessous. Écris les caractères chinois, le pinyin avec tons ou la traduction française manquante.
\begin{center}
\begin{tabularx}{\linewidth}{|X|X|X|X|}
\hline
\rowcolor{popblueL} \textbf{Caractères} & \textbf{Pinyin} & \textbf{Nature} & \textbf{Traduction} \\
\hline
青少年犯罪 &  & n. &  \\
\hline
 & fàn zuì & v. & commettre un crime / délit \\
\hline
原因 &  & n. &  \\
\hline
 & jiě jué & v. &  \\
\hline
预防 &  & v. &  \\
\hline
 &  & n./adj. & éducation / éducatif \\
\hline
惩罚 & chéng fá &  &  \\
\hline
 &  & conj. & au contraire / inversement \\
\hline
\end{tabularx}
\end{center}
\end{exercice}

\begin{exercice}[Pinyin et tons : dictée silencieuse]
Écris le pinyin complet avec les marques de tons (ex: \emph{qīngshàonián}) pour les 10 mots suivants, présentés ici en caractères seulement.
\begin{enumerate}
    \item 法律
    \item 犯罪
    \item 原因
    \item 解决
    \item 预防
    \item 教育
    \item 惩罚
    \item 家庭
    \item 社会
    \item 和睦
\end{enumerate}
\end{exercice}

\courssub{Je m'entraîne}

\begin{exercice}[Traduction : français $\rightarrow$ chinois]
Traduis les phrases suivantes en chinois (caractères + pinyin). Utilise le vocabulaire et les structures de la leçon (\emph{比}, \emph{被}, \emph{首先/然后/最后}, \emph{我认为}).
\begin{enumerate}
    \item La prévention est plus importante que la punition.
    \item Le voleur a été attrapé par la police hier.
    \item À mon avis, la cause principale est le manque d'éducation familiale.
    \item D'abord, nous devons analyser les causes ; ensuite, nous proposons des solutions ; enfin, nous agissons ensemble.
    \item Ce jeune a commis un crime parce qu'il n'avait pas de travail.
\end{enumerate}
\end{exercice}

\begin{exercice}[Transformation : voix passive avec \textbf{被}]
Transforme les phrases actives suivantes en phrases passives avec \emph{被}. Fais attention à l'ordre des mots : Sujet + 被 + Agent + Verbe (+ Complément de résultat/degré).
\begin{enumerate}
    \item 父母教育孩子。
    \item 学校解决了问题。
    \item 社会帮助青少年。
    \item 法律惩罚罪犯。
    \item 大家预防犯罪。
\end{enumerate}
\end{exercice}

\begin{exercice}[Texte à trous : formules fonctionnelles du débat]
Complète le dialogue avec les expressions données dans l'encadré. Chaque expression s'utilise une seule fois. Écris les caractères chinois dans les blancs.
\begin{center}
\begin{tabular}{|c|}
\hline
\textbf{我认为 \quad 你觉得呢 \quad 总之 \quad 一方面 \quad 另一方面 \quad 我不同意} \\
\hline
\end{tabular}
\end{center}
\begin{textebox}[Débat : Punir ou éduquer ?]
A: 关于青少年犯罪，\underline{\hspace{3cm}}，预防比惩罚重要。
B: \underline{\hspace{3cm}}。\underline{\hspace{3cm}}，惩罚也是必要的；\underline{\hspace{3cm}}，教育更根本。
A: \underline{\hspace{3cm}}。没有法律约束，教育很难生效。
B: \underline{\hspace{3cm}}，我们要法律和教育结合。
\par\emph{A : Concernant la délinquance juvénile, \dots, la prévention est plus importante que la punition.}
\par\emph{B : \dots . \dots, la punition est aussi nécessaire ; \dots, l'éducation est plus fondamentale.}
\par\emph{A : \dots . Sans contrainte légale, l'éducation peine à faire effet.}
\par\emph{B : \dots, nous devons combiner la loi et l'éducation.}
\end{textebox}
\end{exercice}

\begin{exercice}[Remise en ordre \& Appariements]
\begin{enumerate}
    \item Remets les mots dans l'ordre pour former des phrases correctes (écris caractères + pinyin).
    \begin{enumerate}
        \item 犯罪 / 青少年 / 预防 / 关键 / 是 / 教育
        \item 家庭 / 和睦 / 如果 / 孩子 / 不 / 犯罪 / 容易
        \item 被 / 小偷 / 警察 / 抓住 / 了
    \end{enumerate}
    \item Apparie le début de phrase (1-4) à sa fin logique (a-d).
    \begin{tabular}{ll}
    1. 我认为预防为主， & a. 才能根本解决问题。 \\
    2. 只有家校社配合， & b. 惩罚为辅。 \\
    3. 青少年犯罪的原因很复杂， & c. 我们要综合分析。 \\
    4. 面对困难， & d. 我们不能放弃。 \\
    \end{tabular}
\end{enumerate}
\end{exercice}

\courssub{Je me prépare au Bac}

\begin{exercice}[Compréhension de l'écrit et langue — Type Bac]
Lis le texte original ci-dessous, puis réponds aux questions en français. Les réponses courtes sont acceptées pour la compréhension ; la langue demande des citations précises.
\begin{textebox}[Texte original : La délinquance à Yaoundé et à Beijing — regards croisés]
近年来，喀麦隆首都雅温得和中国首都北京都面临青少年犯罪上升的挑战。在雅温得，许多专家认为，贫困、辍学和家庭破裂是主要原因。一些年轻人因为找不到工作，加入了街头帮派，抢劫、吸毒现象时有发生。社区领袖呼吁加强职业培训，给年轻人希望。
在北京，情况略有不同。虽然经济发达，但学业压力大、网络成瘾、独生子女娇惯成为新诱因。专家建议：\textbf{首先}，减轻学生课业负担；\textbf{然后}，加强心理健康教育；\textbf{最后}，完善未成年人保护法。两国经验表明：\textbf{单靠严厉惩罚不能解决根本问题}，预防和教育才是长久之计。
(Jìnnián, Kāmàilóng shǒudū Yǎwēndé hé Zhōngguó shǒudū Běijīng dōu miànlín qīngshàonián fànzuì shàngshēng de tiǎozhàn. Zài Yǎwēndé, xǔduō zhuānjiā rènwéi, pínkùn, xuèxué hé jiātíng pòliè shì zhǔyào yuányīn. Yīxiē niánqīngrén yīnwèi zhǎobudào gōngzuò, jiārùle jiētou bāngpài, qiǎngjié, xīdú xiànxiàng shíyǒu fāshēng. Shèqū lǐngxiù hūyù jiāqiáng zhíyè péixùn, gěi niánqīngrén xīwàng. Zài Běijīng, qíngkuàng lüeyǒu bùtóng. Suīrán jīngjì fādá, dàn xuéyè yālì dà, wǎngluò chéngyǐn, dúshēngzǐnǚ jiāohuàn chéngwéi xīn yòuyīn. Zhuānjiā jiànyì: Shǒuxiǎn, jiǎnqīng xuéshēng kèyè fùdān; Ránhòu, jiāqiáng xīnlǐ jiànkāng jiàoyù; Zuìhòu, wánshàn wèichéngniánrén bǎohù fǎ. Liǎngguó jīngyàn biǎomíng: Dān kào yánlì chéngfá bùnéng jiějué gēnběn wèntí, yùfáng hé jiàoyù cái shì chángjiǔ zhījì.)
\par\emph{Depuis quelques années, Yaoundé (capitale du Cameroun) et Beijing (capitale de la Chine) font face au défi de la hausse de la délinquance juvénile. À Yaoundé, beaucoup d'experts estiment que la pauvreté, le décrochage scolaire et la rupture familiale sont les causes principales. Certains jeunes, ne trouvant pas de travail, rejoignent des gangs de rue ; vols et toxicomanie surviennent parfois. Les leaders communautaires appellent au renforcement de la formation professionnelle pour donner de l'espoir aux jeunes. À Beijing, la situation est légèrement différente. Bien que l'économie soit développée, la pression scolaire, la cyberdépendance et le gâtisme des enfants uniques deviennent de nouveaux facteurs. Experts suggèrent : \textbf{Premièrement}, alléger la charge scolaire ; \textbf{Deuxièmement}, renforcer l'éducation à la santé mentale ; \textbf{Troisièmement}, perfectionner la loi sur la protection des mineurs. L'expérience des deux pays montre : \textbf{Compter seulement sur une punition sévère ne peut résoudre le problème de fond}, la prévention et l'éducation sont la solution durable.}
\end{textebox}
\textbf{Compréhension (6 points)}
\begin{enumerate}
    \item Cite deux causes de la délinquance spécifiques à Yaoundé mentionnées dans le texte.
    \item Quelle solution proposent les leaders communautaires à Yaoundé ?
    \item Cite trois causes spécifiques à Beijing.
    \item Quelles sont les trois mesures préconisées par les experts chinois (utilise les mots \emph{首先}, \emph{然后}, \emph{最后}) ?
    \item Quelle conclusion commune tire l'auteur des deux expériences ?
\end{enumerate}
\textbf{Langue (4 points)}
\begin{enumerate}
    \setcounter{enumi}{5}
    \item Repère dans le texte la phrase qui exprime la voix passive « ne peut être résolu par... » (structure \emph{被} ou \emph{靠} passif) et recopie-la en caractères.
    \item Trouve dans le texte le mot qui signifie « décrochage scolaire » et écris son pinyin.
    \item Donne l'antonyme (contraire) de \emph{严厉} (yánlì) présent dans le texte ou de ton savoir.
    \item Transforme la phrase « 专家建议减轻学生课业负担 » en utilisant la structure \emph{建议 + 某人 + 做某事} (suggérer à qqn de faire qqch).
\end{enumerate}
\end{exercice}

\begin{exercice}[Thème et Version — Type Bac]
\textbf{A. Thème (Français $\rightarrow$ Chinois) — 5 points}
Traduis le paragraphe suivant en chinois (caractères + pinyin). Soigne les structures (\emph{因为...所以...}, \emph{不仅...而且...}, \emph{被}) et le vocabulaire de la leçon.
\emph{« Au Cameroun, la délinquance juvénile inquiète la société. Parce que beaucoup de jeunes quittent l'école tôt et ne trouvent pas de travail, ils commettent des vols. La police les arrête souvent, mais la prison ne règle pas le problème. Je pense que l'éducation et la formation professionnelle sont plus utiles que la punition. Seule la coopération entre la famille, l'école et la société peut protéger nos jeunes. »}

\textbf{B. Version (Chinois $\rightarrow$ Français) — 5 points}
Traduis le texte suivant en français courant et correct.
\begin{textebox}
预防青少年犯罪，不能只靠警察抓人，也不能只靠法院判刑。关键在于“预防”二字。家庭要给孩子爱和规矩，学校要教知识也教做人，社会要提供机会不歧视。大家一起努力，青少年才有光明的未来。
(Yùfáng qīngshàonián fànzuì, bùnéng zhǐ kào jǐngchá zhuārén, yě bùnéng zhǐ kào fǎyuàn pànxíng. Guānjiàn zài yú “yùfáng” èr zì. Jiātíng yào gěi háizi ài hé guīju, xuéxiào yào jiāo zhīshì yě jiāo zuòrén, shèhuì yào tígōng jīhuì bù qíshì. Dàjiā yīqǐ nǔlì, qīngshàonián cái yǒu guāngmíng de wèilái.)
\end{textebox}
\end{exercice}

\begin{exercice}[Expression orale / Écrite : Exposé et Débat — Type Bac Oral]
\textbf{Sujet :} \emph{« Les causes de la délinquance juvénile au Cameroun et pistes de solutions »}.
\textbf{Consignes :}
\begin{enumerate}
    \item Prépare un exposé oral de \textbf{2 minutes maximum} (environ \textbf{150--200 caractères chinois}).
    \item Ton exposé doit être structuré avec les connecteurs : \textbf{首先} (Premièrement), \textbf{然后} (Deuxièmement / Ensuite), \textbf{最后} (Enfin).
    \item Tu dois utiliser \textbf{au moins 5 mots} du glossaire de la leçon (ex: \emph{青少年犯罪, 原因, 贫困, 辍学, 预防, 教育, 职业培训, 家庭, 社会, 解决}).
    \item Utilise au moins une phrase d'opinion (\emph{我认为 / 我觉得 / 建议}) et une phrase de relance pour engager le débat (\emph{你觉得呢？ / 大家怎么看？}).
    \item Écris ton texte en caractères chinois et en pinyin sur ta copie. Tu seras évalué sur la prononciation (tons), la fluidité, la correction grammaticale et la richesse lexicale.
\end{enumerate}
\end{exercice}

\begin{experience}[Simulation d'examen oral : Jeu de rôle noté]
\textbf{Binôme :} Élève A = Journaliste de CRTV / CCTV ; Élève B = Chef de quartier / Expert en prévention.
\textbf{Scénario :} Un reportage sur la hausse des vols à l'arraché près des lycées.
\textbf{Contraintes :} Minimum \textbf{6 échanges} (3 par personne). Obligation d'inclure : \textbf{2 relances} (\emph{你觉得呢？}, \emph{还有什么建议？}), \textbf{1 désaccord poli} (\emph{我不太同意……因为……}, \emph{这个观点有道理，但是……}). Utiliser le vocabulaire : \emph{抢劫, 治安, 巡逻, 监控, 预防}.
\textbf{Déroulement :} 5 min préparation $\rightarrow$ 3 min passage oral $\rightarrow$ 2 min feedback prof/pairs (grille : tons, vocab, structures, interaction).
\end{experience}


\newpage

\coursec{Corrigés des exercices}

\coursec{Corrigés détaillés — Leçon 7 : Expression orale – Délinquance juvénile}

\begin{corrige}[Exercice 1 — QCM : vocabulaire et structures]
\begin{enumerate}
    \item \textbf{Réponse a :} \zh{青少年犯罪} (\emph{qīngshàonián fànzuì}) = « délinquance juvénile ». Les autres options signifient : b) \zh{青少年教育} (\emph{qīngshàonián jiàoyù}) éducation des jeunes, c) \zh{青少年问题} (\emph{qīngshàonián wèntí}) problèmes des jeunes, d) \zh{青少年心理} (\emph{qīngshàonián xīnlǐ}) psychologie des jeunes.
    \item \textbf{Réponse b :} \zh{我认为} (\emph{wǒ rènwéi}) est la formule formelle d'opinion (« je considère que »). \zh{我觉得} (\emph{wǒ juéde}) est correct mais plus oral ; \zh{我猜} (\emph{wǒ cāi}) = « je devine » ; \zh{我希望} (\emph{wǒ xīwàng}) = « j'espère ».
    \item \textbf{Réponse b :} \zh{你觉得呢？} (\emph{Nǐ juéde ne?}) = « Et toi, qu'en penses-tu ? » — la relance typique du débat.
    \item \textbf{Réponse b :} \zh{被} (\emph{bèi}) marque la voix passive : Sujet (patient) + \zh{被} + Agent + Verbe. \zh{把} (\emph{bǎ}) est la construction de disposition (active), \zh{让} (\emph{ràng})/\zh{给} (\emph{gěi}) peuvent avoir un emploi passif familier mais la marque standard est \zh{被}.
    \item \textbf{Réponse c :} \zh{最后} (\emph{zuìhòu}) = « enfin ». \zh{首先} (\emph{shǒuxiān}) = d'abord, \zh{然后} (\emph{ránhòu}) = ensuite, \zh{另外} (\emph{lìngwài}) = de plus.
\end{enumerate}
\end{corrige}

\begin{corrige}[Exercice 2 — Vrai ou Faux justifié]
\begin{enumerate}
    \item \textbf{VRAI.} Justification : \zh{家庭教育是关键} (\emph{jiātíng jiàoyù shì guānjiàn}) — « l'éducation familiale est la clé ».
    \item \textbf{FAUX.} Le texte dit le contraire : \zh{父母应该多陪伴孩子，少打骂} (\emph{fùmǔ yīnggāi duō péibàn háizi, shǎo dǎmà}) — « les parents doivent accompagner davantage les enfants et moins les frapper/gronder ».
    \item \textbf{FAUX.} Justification : \zh{学校和社会也要配合，共同解决问题} (\emph{xuéxiào hé shèhuì yě yào pèihé, gòngtóng jiějué wèntí}) — l'école et la société doivent coopérer.
    \item \textbf{VRAI.} Justification : \zh{如果家庭和睦，孩子就不容易走上歪路} (\emph{rúguǒ jiātíng hémù, háizi jiù bù róngyì zǒushàng wāilù}) — une famille harmonieuse éloigne l'enfant de la mauvaise voie.
\end{enumerate}
\end{corrige}

\begin{corrige}[Exercice 3 — Vocabulaire : associer et compléter]
\begin{center}
\begin{tabularx}{\linewidth}{|X|X|X|X|}
\hline
\rowcolor{popblueL} \textbf{Caractères} & \textbf{Pinyin} & \textbf{Nature} & \textbf{Traduction} \\
\hline
青少年犯罪 & qīngshàonián fànzuì & n. & délinquance juvénile \\
\hline
犯罪 & fànzuì & v. & commettre un crime / délit \\
\hline
原因 & yuányīn & n. & cause, raison \\
\hline
解决 & jiějué & v. & résoudre \\
\hline
预防 & yùfáng & v. & prévenir \\
\hline
教育 & jiàoyù & n./adj. & éducation / éducatif \\
\hline
惩罚 & chéngfá & v./n. & punir / punition \\
\hline
反而 & fǎnér & conj. & au contraire / inversement \\
\hline
\end{tabularx}
\end{center}
\textbf{Points de vigilance :} \zh{犯罪} : attention au 4\textsuperscript{e} ton sur \emph{fàn} et au 4\textsuperscript{e} ton sur \emph{zuì} ; \zh{原因} (\emph{yuányīn}) ne se confond pas avec \zh{因为} (\emph{yīnwèi}, « parce que ») ; \zh{反而} introduit un résultat contraire à l'attente, ne pas l'employer comme simple « mais » (\zh{但是}).
\end{corrige}

\begin{corrige}[Exercice 4 — Pinyin et tons : dictée silencieuse]
\begin{enumerate}
    \item 法律 : \emph{fǎlǜ} (3\textsuperscript{e} + 4\textsuperscript{e} ton ; attention au \emph{ü} de \emph{lǜ})
    \item 犯罪 : \emph{fànzuì} (4 + 4)
    \item 原因 : \emph{yuányīn} (2 + 1)
    \item 解决 : \emph{jiějué} (3 + 2)
    \item 预防 : \emph{yùfáng} (4 + 2)
    \item 教育 : \emph{jiàoyù} (4 + 4)
    \item 惩罚 : \emph{chéngfá} (2 + 2) — piège fréquent : ne pas lire \emph{chéngfà} (4\textsuperscript{e} ton)
    \item 家庭 : \emph{jiātíng} (1 + 2)
    \item 社会 : \emph{shèhuì} (4 + 4)
    \item 和睦 : \emph{hémù} (2 + 4)
\end{enumerate}
\textbf{Rappel (sablière) :} devant un autre 3\textsuperscript{e} ton, le premier 3\textsuperscript{e} ton se prononce comme un 2\textsuperscript{e} ton (ex. : \zh{你好} \emph{nǐhǎo} $\rightarrow$ \emph{níhǎo}), mais l'écriture pinyin conserve le ton d'origine. \zh{法律} (\emph{fǎlǜ}) n'est \emph{pas} concerné car le second ton est un 4\textsuperscript{e} ton.
\end{corrige}

\begin{corrige}[Exercice 5 — Traduction : français $\rightarrow$ chinois]
\begin{enumerate}
    \item \zh{预防比惩罚更重要。} \emph{Yùfáng bǐ chéngfá gèng zhòngyào.} — Structure comparative : A + \zh{比} + B + (更) + adjectif.
    \item \zh{小偷昨天被警察抓住了。} \emph{Xiǎotōu zuótiān bèi jǐngchá zhuāzhù le.} — Passif : patient + \zh{被} + agent + verbe + complément de résultat \zh{住} + \zh{了}.
    \item \zh{我认为，主要原因是缺少家庭教育。} \emph{Wǒ rènwéi, zhǔyào yuányīn shì quēshǎo jiātíng jiàoyù.}
    \item \zh{首先，我们要分析原因；然后，我们提出解决办法；最后，我们一起行动。} \emph{Shǒuxiān, wǒmen yào fēnxī yuányīn; ránhòu, wǒmen tíchū jiějué bànfǎ; zuìhòu, wǒmen yīqǐ xíngdòng.}
    \item \zh{这个年轻人因为没有工作，所以犯了罪。} \emph{Zhège niánqīngrén yīnwèi méiyǒu gōngzuò, suǒyǐ fàn le zuì.} — Corrélation \zh{因为}…\zh{所以}… ; en chinois, on garde souvent les deux (contrairement au français où « parce que » exclut « donc »).
\end{enumerate}
\textbf{Points de vigilance :} le complément de temps \zh{昨天} se place avant \zh{被} (et avant le verbe) ; ne jamais traduire « que » du comparatif par un mot : \zh{比} suffit.
\end{corrige}

\begin{corrige}[Exercice 6 — Transformation : voix passive avec 被]
Règle : Patient + \zh{被} + Agent + Verbe (+ complément obligatoire : \zh{了}, résultat, degré…).
\begin{enumerate}
    \item \zh{孩子被父母教育了。} \emph{Háizi bèi fùmǔ jiàoyù le.}
    \item \zh{问题被学校解决了。} \emph{Wèntí bèi xuéxiào jiějué le.}
    \item \zh{青少年被社会帮助了。} \emph{Qīngshàonián bèi shèhuì bāngzhù le.} (plus naturel : \zh{青少年受到社会的帮助。} \emph{Qīngshàonián shòudào shèhuì de bāngzhù.})
    \item \zh{罪犯被法律惩罚了。} \emph{Zuìfàn bèi fǎlǜ chéngfá le.}
    \item \zh{犯罪被大家预防了。} \emph{Fànzuì bèi dàjiā yùfáng le.} (phrase grammaticale mais peu idiomatique ; préférer \zh{犯罪应该被大家共同预防。} \emph{Fànzuì yīnggāi bèi dàjiā gòngtóng yùfáng.})
\end{enumerate}
\textbf{Points de vigilance :} le verbe au passif avec \zh{被} ne peut presque jamais rester « nu » : il faut un complément (\zh{了}, résultat, degré) ; l'agent suit immédiatement \zh{被} ; on ne dit jamais \zh{被父母教育孩子} (ordre français).
\end{corrige}

\begin{corrige}[Exercice 7 — Texte à trous : formules fonctionnelles du dialogue de débat]
\textbf{Expressions à placer (chacune une fois) :} \zh{我认为}、\zh{我不同意}、\zh{一方面}、\zh{另一方面}、\zh{你觉得呢}、\zh{总之}。

\textbf{Dialogue complet corrigé :}
\begin{textebox}[Dialogue complété]
A : 关于青少年犯罪，\textbf{我认为}，预防比惩罚重要。\\
B : \textbf{我不同意}。\textbf{一方面}，惩罚也是必要的；\textbf{另一方面}，教育更根本。\\
A : \textbf{你觉得呢}。没有法律约束，教育很难生效。\\
B : \textbf{总之}，我们要法律和教育结合。
\end{textebox}

\textbf{Traduction :}
A : « En ce qui concerne la délinquance juvénile, je pense que la prévention est plus importante que la punition. »
B : « Je ne suis pas d'accord. D'un côté, la punition est aussi nécessaire ; de l'autre, l'éducation est plus fondamentale. »
A : « Qu'en penses-tu ? Sans contrainte légale, l'éducation peine à faire effet. »
B : « En somme, nous devons combiner loi et éducation. »

\textbf{Analyse :} \zh{我认为} ouvre l'exposé ; \zh{我不同意} marque le désaccord ; \zh{一方面}…\zh{另一方面} structurent la nuance ; \zh{你觉得呢} sert ici de relance rhétorique avant l'argument d'A ; \zh{总之} conclut la synthèse de B.
\end{corrige}

\begin{corrige}[Exercice 8 — Remise en ordre \& Appariements]
\textbf{1. Remise en ordre :}
\begin{enumerate}
    \item \zh{预防青少年犯罪，教育是关键。} \emph{Yùfáng qīngshàonián fànzuì, jiàoyù shì guānjiàn.} « Pour prévenir la délinquance juvénile, l'éducation est la clé. »
    \item \zh{如果家庭和睦，孩子不容易犯罪。} \emph{Rúguǒ jiātíng hémù, háizi bù róngyì fànzuì.} « Si la famille est harmonieuse, l'enfant ne bascule pas facilement dans la délinquance. »
    \item \zh{小偷被警察抓住了。} \emph{Xiǎotōu bèi jǐngchá zhuāzhù le.} « Le voleur a été attrapé par la police. »
\end{enumerate}

\textbf{2. Appariements :}
\begin{itemize}
    \item 1 \rightarrow b : \zh{我认为预防为主，惩罚为辅。} (\emph{Wǒ rènwéi yùfáng wéi zhǔ, chéngfá wéi fǔ.}) « La prévention en priorité, la punition en appoint. »
    \item 2 \rightarrow a : \zh{只有家校社配合，才能根本解决问题。} (\emph{Zhǐyǒu jiā xiào shè pèihé, cáinéng gēnběn jiějué wèntí.}) Corrélation \zh{只有}…\zh{才}…
    \item 3 \rightarrow c : Causes complexes $\rightarrow$ analyse globale (\zh{原因复杂，需全面分析}).
    \item 4 \rightarrow d : Face aux difficultés, ne pas abandonner (\zh{遇到困难，不放弃}).
\end{itemize}
\end{corrige}

\begin{corrige}[Exercice 9 — Compréhension de l'écrit et langue — Type Bac]
\textbf{Barème indicatif :} Compréhension 6 pts (1 + 1 + 1,5 + 1,5 + 1) ; Langue 4 pts (1 pt par question).

\textbf{Compréhension :}
\begin{enumerate}
    \item Deux causes propres à Yaoundé : la pauvreté (\zh{贫困} \emph{pínkùn}) et le décrochage scolaire (\zh{辍学} \emph{chuòxué}) ; on accepte aussi la rupture familiale (\zh{家庭破裂} \emph{jiātíng pòliè}) et le chômage (\zh{失业} \emph{shīyè}).
    \item Les leaders communautaires proposent de renforcer la formation professionnelle (\zh{加强职业培训} \emph{jiāqiáng zhíyè péixùn}) pour redonner espoir aux jeunes.
    \item À Beijing : la pression scolaire (\zh{学业压力大} \emph{xuéyè yālì dà}), la cyberdépendance (\zh{网络成瘾} \emph{wǎngluò chéngyǐn}), le gâtisme des enfants uniques (\zh{独生子女娇惯} \emph{dúshēngzǐnǚ jiāoguàn}).
    \item \zh{首先} (\emph{shǒuxiān}) : alléger la charge scolaire (\zh{减轻学生课业负担} \emph{jiǎnqīng xuéshēng kèyè fùdān}) ; \zh{然后} (\emph{ránhòu}) : renforcer l'éducation à la santé mentale (\zh{加强心理健康教育} \emph{jiāqiáng xīnlǐ jiànkāng jiàoyù}) ; \zh{最后} (\emph{zuìhòu}) : perfectionner la loi de protection des mineurs (\zh{完善未成年人保护法} \emph{wánshàn wèichéngniánrén bǎohù fǎ}).
    \item Conclusion commune : la punition sévère seule ne résout pas le problème de fond ; seules la prévention et l'éducation constituent une solution durable (\zh{预防和教育才是长久之计} \emph{yùfáng hé jiàoyù cái shì chángjiǔ zhī jì}).
\end{enumerate}

\textbf{Langue :}
\begin{enumerate}
    \setcounter{enumi}{5}
    \item \zh{单靠严厉惩罚不能解决根本问题} (\emph{Dān kào yánlì chéngfá bùnéng jiějué gēnběn wèntí}) — structure active avec \zh{靠} (\emph{kào}, « compter sur ») : « Compter uniquement sur la punition sévère ne peut résoudre… » (pas de passif \zh{被} ici).
    \item « Décrochage scolaire » : \zh{辍学} (\emph{chuòxué}).
    \item Antonyme de \zh{严厉} (\emph{yánlì}, sévère) : \zh{宽容} (\emph{kuānróng}, tolérant) ou \zh{温和} (\emph{wēnhé}, doux).
    \item \zh{专家建议学生减轻课业负担。} \emph{Zhuānjiā jiànyì xuéshēng jiǎnqīng kèyè fùdān.} — Structure : \zh{建议} + personne + action (construction en éventail).
\end{enumerate}
\textbf{Points de vigilance :} en question 6, ne pas inventer une phrase avec \zh{被} absente du texte : il fallait repérer la tournure avec \zh{靠} ; pour 9, l'objet de \zh{建议} devient le sujet logique du second verbe.
\end{corrige}

\begin{corrige}[Exercice 10 — Thème et Version — Type Bac]
\textbf{Barème indicatif :} 5 pts par partie (1 pt par phrase/tronçon, tolérance 0,5 pour une faute de ton mineure).

\textbf{A. Thème — proposition de traduction :}
\begin{textebox}[Version chinoise]
在喀麦隆，青少年犯罪让社会担忧。因为很多年轻人早早辍学，找不到工作，所以他们偷窃。警察常常抓住他们，但是监狱不能解决问题。我认为，教育和职业培训比惩罚更有用。只有家庭、学校和社会合作，才能保护我们的年轻人。
\end{textebox}
\emph{Zài Kāmàilóng, qīngshàonián fànzuì ràng shèhuì dānyōu. Yīnwèi hěn duō niánqīngrén zǎozǎo chuòxué, zhǎobudào gōngzuò, suǒyǐ tāmen tōuqiè. Jǐngchá chángcháng zhuāzhù tāmen, dànshì jiānyù bùnéng jiějué wèntí. Wǒ rènwéi, jiàoyù hé zhíyè péixùn bǐ chéngfá gèng yǒuyòng. Zhǐyǒu jiātíng, xuéxiào hé shèhuì hézuò, cái néng bǎohù wǒmen de niánqīngrén.}

\textbf{Points de vigilance :} « inquiète la société » rendu par \zh{让社会担忧} (\emph{ràng shèhuì dānyōu}) ; « la police les arrête » : \zh{抓住} (\emph{zhuāzhù}) plus naturel que \zh{逮捕} (\emph{dǎibǔ}) au niveau HSK 3 ; « seule… peut » exige la corrélation \zh{只有}…\zh{才}… (\emph{zhǐyǒu}…\emph{cái}…).

\textbf{B. Version — proposition de traduction :}
« Pour prévenir la délinquance juvénile, on ne peut pas compter uniquement sur la police pour arrêter les gens, ni uniquement sur les tribunaux pour prononcer des peines. La clé réside dans le mot "prévention". La famille doit donner à l'enfant de l'amour et des règles ; l'école doit enseigner les savoirs mais aussi le savoir-être ; la société doit offrir des opportunités sans discriminer. Si tout le monde fait des efforts ensemble, les jeunes auront un avenir prometteur. »

\textbf{Points de vigilance :} \zh{教做人} (\emph{jiào zuòrén}) = « apprendre à se comporter / à être un homme (un adulte) », ne pas traduire littéralement « enseigner à faire l'homme » ; \zh{二字} (\emph{èrzì}) = « ces deux caractères » (pré-vention) ; \zh{光明的未来} (\emph{guāngmíng de wèilái}) = « avenir lumineux/prometteur ».
\end{corrige}

\begin{corrige}[Exercice 11 — Expression orale / Écrite : Exposé et Débat — Type Bac Oral]
\textbf{Barème indicatif (10 pts) :} prononciation et tons 3 pts ; grammaire et structures 3 pts ; richesse du vocabulaire (5 mots du glossaire minimum) 2 pts ; structure (\zh{首先}/\zh{然后}/\zh{最后}) et interaction (opinion + relance) 2 pts.

\textbf{Proposition de rédaction modèle (environ 180 caractères) :}
\begin{textebox}[Exposé modèle]
大家好！今天我想谈谈喀麦隆的青少年犯罪问题。首先，我认为主要原因有三个：贫困、辍学和缺少家庭教育。很多年轻人找不到工作，就容易走上歪路。然后，单靠惩罚不能解决根本问题，预防才是关键。最后，我建议：家庭多陪伴孩子，学校加强教育，社会提供职业培训。只有大家一起努力，青少年才有光明的未来。你觉得呢？谢谢大家！
\end{textebox}
\emph{Dàjiā hǎo! Jīntiān wǒ xiǎng tántan Kāmàilóng de qīngshàonián fànzuì wèntí. Shǒuxiān, wǒ rènwéi zhǔyào yuányīn yǒu sān ge: pínkùn, chuòxué hé quēshǎo jiātíng jiàoyù. Hěn duō niánqīngrén zhǎobudào gōngzuò, jiù róngyì zǒushàng wāilù. Ránhòu, dān kào chéngfá bùnéng jiějué gēnběn wèntí, yùfáng cái shì guānjiàn. Zuìhòu, wǒ jiànyì: jiātíng duō péibàn háizi, xuéxiào jiāqiáng jiàoyù, shèhuì tígōng zhíyè péixùn. Zhǐyǒu dàjiā yīqǐ nǔlì, qīngshàonián cái yǒu guāngmíng de wèilái. Nǐ juéde ne? Xièxie dàjiā!}

\textbf{Traduction :} « Bonjour à tous ! Aujourd'hui, je voudrais parler du problème de la délinquance juvénile au Cameroun. Premièrement, je pense qu'il y a trois causes principales : la pauvreté, le décrochage scolaire et le manque d'éducation familiale. Beaucoup de jeunes ne trouvent pas de travail et basculent facilement dans la mauvaise voie. Deuxièmement, compter sur la seule punition ne résout pas le problème de fond : la prévention est la clé. Enfin, je suggère : que les familles accompagnent davantage leurs enfants, que l'école renforce l'éducation, que la société offre une formation professionnelle. Ce n'est qu'en travaillant tous ensemble que les jeunes auront un avenir prometteur. Qu'en pensez-vous ? Merci à tous ! »

\textbf{Points de vigilance pour l'oral :} soigner les tons de \zh{犯罪} (\emph{fànzuì}, 4+4) et \zh{预防} (\emph{yùfáng}, 4+2) ; marquer une courte pause après chaque connecteur (\zh{首先}、\zh{然后}、\zh{最后}) ; la relance finale \zh{你觉得呢？} doit être montante (intonation interrogative naturelle).
\end{corrige}

\begin{corrige}[Exercice 12 — Simulation d'examen oral : Jeu de rôle noté]
\textbf{Grille d'évaluation indicative (10 pts) :} tons et prononciation 3 pts ; vocabulaire imposé employé correctement (\zh{抢劫}、\zh{治安}、\zh{巡逻}、\zh{监控}、\zh{预防}) 2 pts ; structures (relances, désaccord poli) 3 pts ; fluidité et interaction 2 pts.

\textbf{Dialogue modèle (6 échanges) :}
\begin{textebox}[Jeu de rôle : Journaliste / Chef de quartier]
A (journaliste) : \zh{您好！最近学校附近抢劫案件增多，您认为主要原因是什么？} \\
\emph{Nín hǎo! Zuìjìn xuéxiào fùjìn qiǎngjié ànjiàn zēngduō, nín rènwéi zhǔyào yuányīn shì shénme?} \\
B (chef de quartier) : \zh{我认为主要是一些青少年辍学后没有工作，治安压力就大了。} \\
\emph{Wǒ rènwéi zhǔyào shì yīxiē qīngshàonián chuòxué hòu méiyǒu gōngzuò, zhì'ān yālì jiù dà le.} \\
A : \zh{你觉得呢，只靠警察巡逻够吗？} \\
\emph{Nǐ juéde ne, zhǐ kào jǐngchá xúnluó gòu ma?} \\
B : \zh{我不太同意只靠巡逻，因为预防更重要。我们还要安装监控，教育孩子。} \\
\emph{Wǒ bù tài tóngyì zhǐ kào xúnluó, yīnwèi yùfáng gèng zhòngyào. Wǒmen hái yào ānzhuāng jiànkòng, jiàoyù háizi.} \\
A : \zh{这个观点有道理，但是装监控需要很多钱，还有什么建议？} \\
\emph{Zhège guāndiǎn yǒu dàolǐ, dànshì zhuāng jiànkòng xūyào hěn duō qián, hái yǒu shénme jiànyì?} \\
B : \zh{总之，警察巡逻、社区预防和家庭学校教育要结合，治安才能真正改善。} \\
\emph{Zǒngzhī, jǐngchá xúnluó, shèqū yùfáng hé jiātíng xuéxiào jiàoyù yào jiéhé, zhì'ān cái néng zhēnzhèng gǎishàn.}
\end{textebox}

\textbf{Points de vigilance :} le désaccord poli exige \zh{我不太同意…因为…} ou \zh{…有道理，但是…} ; chaque relance (\zh{你觉得呢？}, \zh{还有什么建议？}) doit être clairement prononcée ; vérifier que les 5 mots imposés apparaissent bien au moins une fois (\zh{抢劫、治安、巡逻、监控、预防}).
\end{corrige}

\newpage

\coursec{Fiche bilan}

\begin{popfigure}
\begin{tikzpicture}[mindmap, grow cyclic, every node/.style={concept, align=center, font=\small},
root concept/.append style={concept color=popblue, font=\bfseries},
level 1/.append style={level distance=3.1cm, sibling angle=60},
level 1 concept/.append style={font=\footnotesize}]
\node[root concept, align=center]{Délinquance juvénile\\ \zh{青少年犯罪}\\ \emph{qīngshàonián fànzuì}}
child[concept color=poporange]{node[align=center]{Lexique\\ \zh{犯罪 偷 抢}\\ \zh{法律 警察}}}
child[concept color=popgreen]{node[align=center]{Dialogue 1\\ Parler du\\ problème\\ \zh{越来越严重}}}
child[concept color=poppurple]{node[align=center]{Formules\\ Présenter, opiner\\ \zh{我认为…}\\ \zh{我同意}}}
child[concept color=poppink]{node[align=center]{Dialogue 2\\ Punir ou\\ éduquer ?\\ \zh{应该…还是应该…}}}
child[concept color=popgold]{node[align=center]{Phonétique\\ Les 4 tons\\ + ton neutre}}
child[concept color=popblue!60]{node[align=center]{Débat guidé\\ Jeux de rôle\\ \zh{你的看法呢？}}};
\end{tikzpicture}
\legende{Carte mentale de la leçon : l'expression orale sur la délinquance juvénile.}
\end{popfigure}

\begin{bilan}
\textbf{1. Lexique clé à connaître par cœur}

\begin{center}
\begin{tabularx}{\linewidth}{|X|X|X|X|}
\hline
\rowcolor{popblueL} Caractères & Pinyin & Nature & Traduction \\
\hline
青少年 & qīngshàonián & nom & adolescent, jeune \\
\hline
犯罪 & fànzuì & verbe/nom & commettre un crime ; crime \\
\hline
偷 & tōu & verbe & voler (dérober) \\
\hline
抢 & qiǎng & verbe & voler (avec force), arracher \\
\hline
法律 & fǎlǜ & nom & loi, droit \\
\hline
警察 & jǐngchá & nom & police, policier \\
\hline
教育 & jiàoyù & nom/verbe & éducation ; éduquer \\
\hline
惩罚 & chéngfá & nom/verbe & punition ; punir \\
\hline
原因 & yuányīn & nom & cause, raison \\
\hline
解决 & jiějué & verbe & résoudre \\
\hline
严重 & yánzhòng & adjectif & grave, sérieux \\
\hline
社会 & shèhuì & nom & société \\
\hline
\end{tabularx}
\end{center}

\textbf{2. Structures et formules fonctionnelles}

\begin{center}
\begin{tabularx}{\linewidth}{|X|X|X|}
\hline
\rowcolor{popblueL} Fonction & Formule (caractères + pinyin) & Traduction \\
\hline
Présenter un exposé & \zh{今天我要谈的是…} Jīntiān wǒ yào tán de shì… & Aujourd'hui, je vais parler de… \\
\hline
Donner son avis & \zh{我认为 / 我觉得…} Wǒ rènwéi / wǒ juéde… & Je pense que… / Je trouve que… \\
\hline
Être d'accord & \zh{我同意你的看法。} Wǒ tóngyì nǐ de kànfǎ. & Je suis d'accord avec ton avis. \\
\hline
Ne pas être d'accord & \zh{我不太同意，因为…} Wǒ bù tài tóngyì, yīnwèi… & Je ne suis pas tout à fait d'accord, parce que… \\
\hline
Relancer le débat & \zh{你的看法呢？} Nǐ de kànfǎ ne ? & Et toi, quel est ton avis ? \\
\hline
Opposer deux solutions & \zh{应该A还是应该B？} Yīnggāi A háishi yīnggāi B ? & Faut-il A ou faut-il B ? \\
\hline
Décrire une aggravation & \zh{越来越严重} yuè lái yuè yánzhòng & de plus en plus grave \\
\hline
\end{tabularx}
\end{center}

\textbf{3. Points de grammaire à maîtriser}
\begin{itemize}
\item \cle{越来越 + adjectif} : exprime une évolution progressive. Ex. \zh{这个问题越来越严重。} \emph{Zhège wèntí yuè lái yuè yánzhòng.} « Ce problème devient de plus en plus grave. »
\item \cle{应该} + verbe : exprime le devoir, le conseil. Ex. \zh{我们应该教育他们。} \emph{Wǒmen yīnggāi jiàoyù tāmen.} « Nous devrions les éduquer. »
\item \cle{因为…所以…} : exprime la cause et la conséquence. Ex. \zh{因为他们没有钱，所以去偷东西。} \emph{Yīnwèi tāmen méiyǒu qián, suǒyǐ qù tōu dōngxi.} « Parce qu'ils n'ont pas d'argent, ils volent des choses. »
\item \cle{A还是B？} : question alternative (choix), sans \zh{吗}. Ex. \zh{应该惩罚还是应该教育？} « Faut-il punir ou éduquer ? »
\item \cle{我觉得 / 我认为} + phrase complète : introduit une opinion personnelle.
\end{itemize}

\textbf{4. Phonétique : les tons à travailler}
\begin{itemize}
\item 1\textsuperscript{er} ton haut et plat : \zh{偷} tōu, \zh{青} qīng.
\item 2\textsuperscript{e} ton montant : \zh{年} nián, \zh{来} lái.
\item 3\textsuperscript{e} ton descendant-montant : \zh{警} jǐng, \zh{法} fǎ.
\item 4\textsuperscript{e} ton descendant sec : \zh{罪} zuì, \zh{教} jiào.
\item Attention aux suites : \zh{犯罪} fànzuì (4\textsuperscript{e} + 4\textsuperscript{e}), \zh{严重} yánzhòng (2\textsuperscript{e} + 4\textsuperscript{e}).
\end{itemize}

\textbf{5. Méthode express : réussir un exposé oral}
\begin{enumerate}
\item Annoncer le sujet : \zh{今天我要谈的是青少年犯罪。}
\item Décrire le problème avec \zh{越来越} et \zh{因为…所以…}.
\item Donner ton avis : \zh{我认为…}.
\item Proposer une solution avec \zh{应该}.
\item Conclure et relancer : \zh{你的看法呢？}
\end{enumerate}
\end{bilan}

\begin{aretenir}[Checklist avant le Bac]
\begin{itemize}
\item[$\square$] Je connais par cœur le lexique du thème : \zh{青少年、犯罪、偷、抢、法律、警察、教育、惩罚、原因、解决、严重、社会}.
\item[$\square$] Je sais présenter un exposé avec \zh{今天我要谈的是…}.
\item[$\square$] Je sais donner mon avis avec \zh{我认为} et \zh{我觉得}.
\item[$\square$] Je sais exprimer l'accord (\zh{我同意}) et le désaccord poli (\zh{我不太同意，因为…}).
\item[$\square$] Je sais relancer un camarade avec \zh{你的看法呢？}.
\item[$\square$] Je maîtrise la structure \zh{越来越 + adjectif}.
\item[$\square$] Je maîtrise la question alternative \zh{A还是B？} (sans \zh{吗}).
\item[$\square$] Je sais utiliser \zh{因为…所以…} pour expliquer une cause.
\item[$\square$] Je prononce correctement les 4 tons sur les mots clés (\zh{犯罪、严重、警察、法律}).
\item[$\square$] Je sais opposer deux solutions : \zh{应该惩罚还是应该教育？}.
\item[$\square$] Je peux tenir un exposé de 5 phrases sans lire mes notes.
\item[$\square$] Je sais écrire les caractères essentiels : \zh{罪、警、罚、教}.
\end{itemize}
\end{aretenir}

\findecours

\end{document}
